\documentclass[11pt,a4paper]{article}
\usepackage[a4paper,margin=20mm]{geometry}
\usepackage{fontspec}
\setmainfont{Noto Serif}
\setsansfont{Noto Sans}
\usepackage{amsmath,amssymb}
\usepackage{booktabs,tabularx,array}
\usepackage[table]{xcolor}
\usepackage{hyperref}
\definecolor{primary}{RGB}{27,79,114}
\definecolor{accent}{RGB}{40,116,166}
\definecolor{rowtint}{RGB}{235,245,251}
\hypersetup{colorlinks=true,linkcolor=primary,urlcolor=primary}
\setlength{\parskip}{5pt plus 2pt minus 1pt}
\setlength{\parindent}{0pt}
\begin{document}
\begin{center}
  {\LARGE\bfseries\color{primary} PHIẾU HỌC TẬP TRẮC NGHIỆM TOÁN 9}\\[4pt]
  {\normalsize\itshape\color{accent} Bộ sách: Kết nối tri thức với cuộc sống (Tập 1 \& Tập 2 – 10 Chương)}\\[4pt]
  {\small Cấu trúc mỗi chương: 20 câu Dễ (Câu 1--20) $\bullet$ 20 câu Trung bình (Câu 21--40) $\bullet$ 10 câu Khó (Câu 41--50) -- Tổng 50 câu/chương}
\end{center}
\vspace{6pt}\hrule\vspace{10pt}
\section*{\color{primary} CHƯƠNG I: PHƯƠNG TRÌNH VÀ HỆ HAI PHƯƠNG TRÌNH BẬC NHẤT HAI ẨN}
\textit{Trọng tâm kiến thức SGK: Khái niệm phương trình và hệ hai phương trình bậc nhất hai ẩn, giải hệ phương trình bằng phương pháp thế, cộng đại số và giải bài toán bằng cách lập hệ phương trình.}\par\vspace{4pt}
\subsection*{\color{accent} Phần I. Mức độ Dễ (Nhận biết -- Câu 1 đến Câu 20)}
\textbf{\color{primary} Câu 1.} Cặp số nào sau đây là nghiệm của phương trình bậc nhất hai ẩn 2x + y = 7?\par
\hspace*{1.5em}\textbf{A.} (0; 2)\quad\quad \textbf{B.} (3; 2)\quad\quad \textbf{C.} (1; 2)\quad\quad \textbf{D.} (2; 3)\par\vspace{2pt}
\textbf{\color{primary} Câu 2.} Thực hiện yêu cầu nhận biết sau: Nghiệm (x; y) của hệ phương trình: x + y = 9 và x - y = 3 là.\par
\hspace*{1.5em}\textbf{A.} (3; 6)\quad\quad \textbf{B.} (6; 3)\quad\quad \textbf{C.} (6; 3) (1)\quad\quad \textbf{D.} (5; 4)\par\vspace{2pt}
\textbf{\color{primary} Câu 3.} Trong bài học thuộc Chương I, xét nội dung: Phương trình bậc nhất hai ẩn ax + by = c (với a $\ne$ 0 hoặc b $\ne$ 0) luôn có bao nhiêu nghiệm. Đáp án đúng là:\par
\hspace*{1.5em}\textbf{A.} Vô số nghiệm\quad\quad \textbf{B.} Vô nghiệm\quad\quad \textbf{C.} Hai nghiệm phân biệt\quad\quad \textbf{D.} Một nghiệm duy nhất\par\vspace{2pt}
\textbf{\color{primary} Câu 4.} Dựa vào kiến thức chuẩn SGK của Chương I (Phương trình và hệ hai phương trình bậc nhất hai ẩn), cặp số nào sau đây là nghiệm của phương trình bậc nhất hai ẩn 2x + y = 13:\par
\hspace*{1.5em}\textbf{A.} (4; 2)\quad\quad \textbf{B.} (0; 5)\quad\quad \textbf{C.} (5; 3)\quad\quad \textbf{D.} (3; 5)\par\vspace{2pt}
\textbf{\color{primary} Câu 5.} Chọn kết quả đúng nhất khi xét: Nghiệm (x; y) của hệ phương trình: x + y = 15 và x - y = 3 là:\par
\hspace*{1.5em}\textbf{A.} (8; 7)\quad\quad \textbf{B.} (6; 9)\quad\quad \textbf{C.} (9; 6)\quad\quad \textbf{D.} (12; 3)\par\vspace{2pt}
\textbf{\color{primary} Câu 6.} Khi làm bài tập cơ bản về phương trình và hệ hai phương trình bậc nhất hai ẩn, một học sinh gặp câu hỏi: Phương trình bậc nhất hai ẩn ax + by = c (với a $\ne$ 0 hoặc b $\ne$ 0) luôn có bao nhiêu nghiệm. Kết quả nào dưới đây đúng?\par
\hspace*{1.5em}\textbf{A.} Vô nghiệm\quad\quad \textbf{B.} Hai nghiệm phân biệt\quad\quad \textbf{C.} Vô số nghiệm\quad\quad \textbf{D.} Một nghiệm duy nhất\par\vspace{2pt}
\textbf{\color{primary} Câu 7.} Xác định giá trị hoặc khẳng định đúng cho bài toán sau: Cặp số nào sau đây là nghiệm của phương trình bậc nhất hai ẩn 2x + y = 19:\par
\hspace*{1.5em}\textbf{A.} (0; 8)\quad\quad \textbf{B.} (8; 3)\quad\quad \textbf{C.} (3; 8)\quad\quad \textbf{D.} (7; 2)\par\vspace{2pt}
\textbf{\color{primary} Câu 8.} Áp dụng trực tiếp định nghĩa và quy tắc trong Chương I: Nghiệm (x; y) của hệ phương trình: x + y = 21 và x - y = 3 là:\par
\hspace*{1.5em}\textbf{A.} (18; 3)\quad\quad \textbf{B.} (12; 9)\quad\quad \textbf{C.} (11; 10)\quad\quad \textbf{D.} (9; 12)\par\vspace{2pt}
\textbf{\color{primary} Câu 9.} Kiểm tra nhanh kiến thức nền tảng: Phương trình bậc nhất hai ẩn ax + by = c (với a $\ne$ 0 hoặc b $\ne$ 0) luôn có bao nhiêu nghiệm?\par
\hspace*{1.5em}\textbf{A.} Vô số nghiệm\quad\quad \textbf{B.} Một nghiệm duy nhất\quad\quad \textbf{C.} Hai nghiệm phân biệt\quad\quad \textbf{D.} Vô nghiệm\par\vspace{2pt}
\textbf{\color{primary} Câu 10.} Cho dữ kiện nhận biết sau đây: Cặp số nào sau đây là nghiệm của phương trình bậc nhất hai ẩn 2x + y = 25. Phương án chính xác là:\par
\hspace*{1.5em}\textbf{A.} (0; 11)\quad\quad \textbf{B.} (10; 2)\quad\quad \textbf{C.} (11; 3)\quad\quad \textbf{D.} (3; 11)\par\vspace{2pt}
\textbf{\color{primary} Câu 11.} Hãy chỉ ra kết quả đúng trong bốn lựa chọn A, B, C, D đối với yêu cầu: Nghiệm (x; y) của hệ phương trình: x + y = 27 và x - y = 3 là:\par
\hspace*{1.5em}\textbf{A.} (12; 15)\quad\quad \textbf{B.} (15; 12)\quad\quad \textbf{C.} (24; 3)\quad\quad \textbf{D.} (14; 13)\par\vspace{2pt}
\textbf{\color{primary} Câu 12.} Câu hỏi ôn tập lý thuyết và kỹ năng cơ bản (Chương I): Phương trình bậc nhất hai ẩn ax + by = c (với a $\ne$ 0 hoặc b $\ne$ 0) luôn có bao nhiêu nghiệm:\par
\hspace*{1.5em}\textbf{A.} Một nghiệm duy nhất\quad\quad \textbf{B.} Vô số nghiệm\quad\quad \textbf{C.} Vô nghiệm\quad\quad \textbf{D.} Hai nghiệm phân biệt\par\vspace{2pt}
\textbf{\color{primary} Câu 13.} Theo nội dung sách giáo khoa Kết nối tri thức với cuộc sống (Chương I), cặp số nào sau đây là nghiệm của phương trình bậc nhất hai ẩn 2x + y = 31:\par
\hspace*{1.5em}\textbf{A.} (3; 14)\quad\quad \textbf{B.} (13; 2)\quad\quad \textbf{C.} (14; 3)\quad\quad \textbf{D.} (0; 14)\par\vspace{2pt}
\textbf{\color{primary} Câu 14.} Tìm đáp số đúng của bài toán cơ bản sau: Nghiệm (x; y) của hệ phương trình: x + y = 33 và x - y = 3 là:\par
\hspace*{1.5em}\textbf{A.} (17; 16)\quad\quad \textbf{B.} (15; 18)\quad\quad \textbf{C.} (18; 15)\quad\quad \textbf{D.} (30; 3)\par\vspace{2pt}
\textbf{\color{primary} Câu 15.} Xét đối tượng toán học trong Chương I: Phương trình bậc nhất hai ẩn ax + by = c (với a $\ne$ 0 hoặc b $\ne$ 0) luôn có bao nhiêu nghiệm. Giá trị hoặc tính chất tương ứng là:\par
\hspace*{1.5em}\textbf{A.} Hai nghiệm phân biệt\quad\quad \textbf{B.} Vô số nghiệm\quad\quad \textbf{C.} Một nghiệm duy nhất\quad\quad \textbf{D.} Vô nghiệm\par\vspace{2pt}
\textbf{\color{primary} Câu 16.} Hoàn thành phép xác định cơ bản sau đây: Cặp số nào sau đây là nghiệm của phương trình bậc nhất hai ẩn 2x + y = 37:\par
\hspace*{1.5em}\textbf{A.} (3; 17)\quad\quad \textbf{B.} (16; 2)\quad\quad \textbf{C.} (0; 17)\quad\quad \textbf{D.} (17; 3)\par\vspace{2pt}
\textbf{\color{primary} Câu 17.} Trong các khẳng định hoặc giá trị dưới đây, đâu là đáp án đúng cho câu hỏi: Nghiệm (x; y) của hệ phương trình: x + y = 39 và x - y = 3 là?\par
\hspace*{1.5em}\textbf{A.} (20; 19)\quad\quad \textbf{B.} (36; 3)\quad\quad \textbf{C.} (21; 18)\quad\quad \textbf{D.} (18; 21)\par\vspace{2pt}
\textbf{\color{primary} Câu 18.} Bài tập khởi động Chương I: Phương trình bậc nhất hai ẩn ax + by = c (với a $\ne$ 0 hoặc b $\ne$ 0) luôn có bao nhiêu nghiệm:\par
\hspace*{1.5em}\textbf{A.} Vô số nghiệm\quad\quad \textbf{B.} Hai nghiệm phân biệt\quad\quad \textbf{C.} Một nghiệm duy nhất\quad\quad \textbf{D.} Vô nghiệm\par\vspace{2pt}
\textbf{\color{primary} Câu 19.} Quan sát dữ kiện đề bài: Cặp số nào sau đây là nghiệm của phương trình bậc nhất hai ẩn 2x + y = 43. Hãy chọn phương án đúng:\par
\hspace*{1.5em}\textbf{A.} (3; 20)\quad\quad \textbf{B.} (19; 2)\quad\quad \textbf{C.} (0; 20)\quad\quad \textbf{D.} (20; 3)\par\vspace{2pt}
\textbf{\color{primary} Câu 20.} Sử dụng công thức hoặc tính chất đầu bài của Chương I để giải quyết yêu cầu: Nghiệm (x; y) của hệ phương trình: x + y = 45 và x - y = 3 là:\par
\hspace*{1.5em}\textbf{A.} (23; 22)\quad\quad \textbf{B.} (21; 24)\quad\quad \textbf{C.} (42; 3)\quad\quad \textbf{D.} (24; 21)\par\vspace{2pt}
\subsection*{\color{accent} Phần II. Mức độ Trung bình (Thông hiểu -- Câu 21 đến Câu 40)}
\textbf{\color{primary} Câu 21.} Biết hệ phương trình 2x + by = 8 và ax - y = -1 có nghiệm là (x; y) = (1; 3). Khi đó tổng a + b bằng:\par
\hspace*{1.5em}\textbf{A.} 4\quad\quad \textbf{B.} 3\quad\quad \textbf{C.} 5\quad\quad \textbf{D.} 2\par\vspace{2pt}
\textbf{\color{primary} Câu 22.} Vận dụng quy tắc trong Chương I (Phương trình và hệ hai phương trình bậc nhất hai ẩn) để giải bài toán: Giải hệ phương trình: 2x + 3y = 12 và x - y = 1. Giá trị của x + y bằng:\par
\hspace*{1.5em}\textbf{A.} 3\quad\quad \textbf{B.} 7\quad\quad \textbf{C.} 5\quad\quad \textbf{D.} 6\par\vspace{2pt}
\textbf{\color{primary} Câu 23.} Biến đổi và tính toán theo kiến thức Chương I: Với giá trị nào của m thì hệ phương trình mx + 2y = 5 và 3x + y = 1 vô nghiệm. Kết quả thu được là:\par
\hspace*{1.5em}\textbf{A.} m = 6\quad\quad \textbf{B.} m = 2\quad\quad \textbf{C.} m = -6\quad\quad \textbf{D.} m = 3\par\vspace{2pt}
\textbf{\color{primary} Câu 24.} Xét bài toán mức độ thông hiểu sau: Biết hệ phương trình 2x + by = 14 và ax - y = 5 có nghiệm là (x; y) = (4; 3). Khi đó tổng a + b bằng. Đáp án đúng là:\par
\hspace*{1.5em}\textbf{A.} 4\quad\quad \textbf{B.} 2\quad\quad \textbf{C.} 3\quad\quad \textbf{D.} 5\par\vspace{2pt}
\textbf{\color{primary} Câu 25.} Để hoàn thành phiếu học tập Chương I, hãy xác định: Giải hệ phương trình: 2x + 3y = 18 và x - y = 4. Giá trị của x + y bằng:\par
\hspace*{1.5em}\textbf{A.} 6\quad\quad \textbf{B.} 12\quad\quad \textbf{C.} 8\quad\quad \textbf{D.} 10\par\vspace{2pt}
\textbf{\color{primary} Câu 26.} Phân tích dữ kiện bài toán: Với giá trị nào của m thì hệ phương trình mx + 2y = 5 và 3x + y = 1 vô nghiệm. Giá trị cần tìm bằng:\par
\hspace*{1.5em}\textbf{A.} m = 6\quad\quad \textbf{B.} m = 3\quad\quad \textbf{C.} m = 2\quad\quad \textbf{D.} m = -6\par\vspace{2pt}
\textbf{\color{primary} Câu 27.} Áp dụng tính chất đã học ở Chương I vào bài tập sau: Biết hệ phương trình 2x + by = 20 và ax - y = 11 có nghiệm là (x; y) = (7; 3). Khi đó tổng a + b bằng:\par
\hspace*{1.5em}\textbf{A.} 3\quad\quad \textbf{B.} 2\quad\quad \textbf{C.} 4\quad\quad \textbf{D.} 5\par\vspace{2pt}
\textbf{\color{primary} Câu 28.} Thực hiện các bước suy luận hợp lý để giải quyết yêu cầu: Giải hệ phương trình: 2x + 3y = 24 và x - y = 7. Giá trị của x + y bằng:\par
\hspace*{1.5em}\textbf{A.} 11\quad\quad \textbf{B.} 9\quad\quad \textbf{C.} 13\quad\quad \textbf{D.} 18\par\vspace{2pt}
\textbf{\color{primary} Câu 29.} Cho bài toán tự luyện (Phương trình và hệ hai phương trình bậc nhất hai ẩn): Với giá trị nào của m thì hệ phương trình mx + 2y = 5 và 3x + y = 1 vô nghiệm. Phương án nào sau đây chính xác?\par
\hspace*{1.5em}\textbf{A.} m = 3\quad\quad \textbf{B.} m = -6\quad\quad \textbf{C.} m = 2\quad\quad \textbf{D.} m = 6\par\vspace{2pt}
\textbf{\color{primary} Câu 30.} Kết hợp các phép biến đổi trong Chương I, hãy tìm kết quả của bài toán: Biết hệ phương trình 2x + by = 26 và ax - y = 17 có nghiệm là (x; y) = (10; 3). Khi đó tổng a + b bằng:\par
\hspace*{1.5em}\textbf{A.} 2\quad\quad \textbf{B.} 4\quad\quad \textbf{C.} 3\quad\quad \textbf{D.} 5\par\vspace{2pt}
\textbf{\color{primary} Câu 31.} Một bài kiểm tra định kỳ nội dung Chương I nêu câu hỏi: Giải hệ phương trình: 2x + 3y = 30 và x - y = 10. Giá trị của x + y bằng. Hãy chọn đáp số đúng:\par
\hspace*{1.5em}\textbf{A.} 14\quad\quad \textbf{B.} 12\quad\quad \textbf{C.} 24\quad\quad \textbf{D.} 16\par\vspace{2pt}
\textbf{\color{primary} Câu 32.} Tìm lời giải đúng cho bài toán thông hiểu: Với giá trị nào của m thì hệ phương trình mx + 2y = 5 và 3x + y = 1 vô nghiệm:\par
\hspace*{1.5em}\textbf{A.} m = 3\quad\quad \textbf{B.} m = 2\quad\quad \textbf{C.} m = -6\quad\quad \textbf{D.} m = 6\par\vspace{2pt}
\textbf{\color{primary} Câu 33.} Sử dụng mối liên hệ giữa các đại lượng trong Chương I: Biết hệ phương trình 2x + by = 32 và ax - y = 23 có nghiệm là (x; y) = (13; 3). Khi đó tổng a + b bằng:\par
\hspace*{1.5em}\textbf{A.} 2\quad\quad \textbf{B.} 3\quad\quad \textbf{C.} 4\quad\quad \textbf{D.} 5\par\vspace{2pt}
\textbf{\color{primary} Câu 34.} Hãy tính toán và chọn phương án đúng đối với bài tập: Giải hệ phương trình: 2x + 3y = 36 và x - y = 13. Giá trị của x + y bằng:\par
\hspace*{1.5em}\textbf{A.} 15\quad\quad \textbf{B.} 30\quad\quad \textbf{C.} 17\quad\quad \textbf{D.} 19\par\vspace{2pt}
\textbf{\color{primary} Câu 35.} Dựa trên các tính chất trọng tâm của Chương I, xét bài toán: Với giá trị nào của m thì hệ phương trình mx + 2y = 5 và 3x + y = 1 vô nghiệm:\par
\hspace*{1.5em}\textbf{A.} m = -6\quad\quad \textbf{B.} m = 3\quad\quad \textbf{C.} m = 6\quad\quad \textbf{D.} m = 2\par\vspace{2pt}
\textbf{\color{primary} Câu 36.} Thực hành giải toán (Phương trình và hệ hai phương trình bậc nhất hai ẩn): Biết hệ phương trình 2x + by = 38 và ax - y = 29 có nghiệm là (x; y) = (16; 3). Khi đó tổng a + b bằng. Đáp án nào dưới đây thỏa mãn?\par
\hspace*{1.5em}\textbf{A.} 5\quad\quad \textbf{B.} 2\quad\quad \textbf{C.} 3\quad\quad \textbf{D.} 4\par\vspace{2pt}
\textbf{\color{primary} Câu 37.} Kiểm tra kỹ năng biến đổi và suy luận qua bài toán: Giải hệ phương trình: 2x + 3y = 42 và x - y = 16. Giá trị của x + y bằng:\par
\hspace*{1.5em}\textbf{A.} 22\quad\quad \textbf{B.} 20\quad\quad \textbf{C.} 36\quad\quad \textbf{D.} 18\par\vspace{2pt}
\textbf{\color{primary} Câu 38.} Cho giả thiết thuộc phạm vi Chương I: Với giá trị nào của m thì hệ phương trình mx + 2y = 5 và 3x + y = 1 vô nghiệm. Kết luận đúng là:\par
\hspace*{1.5em}\textbf{A.} m = 2\quad\quad \textbf{B.} m = 6\quad\quad \textbf{C.} m = -6\quad\quad \textbf{D.} m = 3\par\vspace{2pt}
\textbf{\color{primary} Câu 39.} Rút gọn hoặc tính giá trị theo yêu cầu sau của Chương I: Biết hệ phương trình 2x + by = 44 và ax - y = 35 có nghiệm là (x; y) = (19; 3). Khi đó tổng a + b bằng:\par
\hspace*{1.5em}\textbf{A.} 2\quad\quad \textbf{B.} 4\quad\quad \textbf{C.} 5\quad\quad \textbf{D.} 3\par\vspace{2pt}
\textbf{\color{primary} Câu 40.} Xác định phương án đúng sau khi thực hiện phép tính / suy luận: Giải hệ phương trình: 2x + 3y = 48 và x - y = 19. Giá trị của x + y bằng:\par
\hspace*{1.5em}\textbf{A.} 23\quad\quad \textbf{B.} 25\quad\quad \textbf{C.} 21\quad\quad \textbf{D.} 42\par\vspace{2pt}
\subsection*{\color{accent} Phần III. Mức độ Khó (Vận dụng \& Vận dụng cao -- Câu 41 đến Câu 50)}
\textbf{\color{primary} Câu 41.} Hai vòi nước cùng chảy vào một bể cạn thì sau 6 giờ đầy bể. Nếu vòi thứ nhất chảy trong 3 giờ và vòi thứ hai chảy trong 6 giờ thì được 3/4 bể (khi hai vòi cùng chảy 6 giờ đầy bể). Thời gian vòi thứ hai chảy một mình đầy bể là:\par
\hspace*{1.5em}\textbf{A.} 12 giờ\quad\quad \textbf{B.} 24 giờ\quad\quad \textbf{C.} 9 giờ\quad\quad \textbf{D.} 18 giờ\par\vspace{2pt}
\textbf{\color{primary} Câu 42.} Bài toán vận dụng (Chương I – Phương trình và hệ hai phương trình bậc nhất hai ẩn): Cho hệ phương trình x + y = m + 3 và x - y = m - 3. Tìm m để nghiệm (x; y) của hệ thỏa mãn x$\cdot$y = 6. Kết quả đúng là:\par
\hspace*{1.5em}\textbf{A.} m = 3\quad\quad \textbf{B.} m = -2\quad\quad \textbf{C.} m = 0\quad\quad \textbf{D.} m = 2\par\vspace{2pt}
\textbf{\color{primary} Câu 43.} Vận dụng kiến thức nâng cao trong phạm vi Chương I để giải bài toán: Hai vòi nước cùng chảy vào một bể cạn thì sau 6 giờ đầy bể. Nếu vòi thứ nhất chảy trong 3 giờ và vòi thứ hai chảy trong 6 giờ thì được 3/4 bể (khi hai vòi cùng chảy 6 giờ đầy bể). Thời gian vòi thứ hai chảy một mình đầy bể là:\par
\hspace*{1.5em}\textbf{A.} 24 giờ\quad\quad \textbf{B.} 12 giờ\quad\quad \textbf{C.} 18 giờ\quad\quad \textbf{D.} 9 giờ\par\vspace{2pt}
\textbf{\color{primary} Câu 44.} Phân tích và giải bài toán tổng hợp sau của Chương I: Cho hệ phương trình x + y = m + 5 và x - y = m - 5. Tìm m để nghiệm (x; y) của hệ thỏa mãn x$\cdot$y = 10:\par
\hspace*{1.5em}\textbf{A.} m = 0\quad\quad \textbf{B.} m = -2\quad\quad \textbf{C.} m = 5\quad\quad \textbf{D.} m = 2\par\vspace{2pt}
\textbf{\color{primary} Câu 45.} Thử thách tư duy trong phạm vi kiến thức Chương I: Hai vòi nước cùng chảy vào một bể cạn thì sau 6 giờ đầy bể. Nếu vòi thứ nhất chảy trong 3 giờ và vòi thứ hai chảy trong 6 giờ thì được 3/4 bể (khi hai vòi cùng chảy 6 giờ đầy bể). Thời gian vòi thứ hai chảy một mình đầy bể là. Đáp số chính xác là:\par
\hspace*{1.5em}\textbf{A.} 9 giờ\quad\quad \textbf{B.} 12 giờ\quad\quad \textbf{C.} 18 giờ\quad\quad \textbf{D.} 24 giờ\par\vspace{2pt}
\textbf{\color{primary} Câu 46.} Sử dụng phương pháp giải toán của Chương I để tìm đáp án cho bài toán: Cho hệ phương trình x + y = m + 7 và x - y = m - 7. Tìm m để nghiệm (x; y) của hệ thỏa mãn x$\cdot$y = 14:\par
\hspace*{1.5em}\textbf{A.} m = 7\quad\quad \textbf{B.} m = -2\quad\quad \textbf{C.} m = 0\quad\quad \textbf{D.} m = 2\par\vspace{2pt}
\textbf{\color{primary} Câu 47.} Xét bài toán vận dụng thực tiễn / nâng cao thuộc Chương I: Hai vòi nước cùng chảy vào một bể cạn thì sau 6 giờ đầy bể. Nếu vòi thứ nhất chảy trong 3 giờ và vòi thứ hai chảy trong 6 giờ thì được 3/4 bể (khi hai vòi cùng chảy 6 giờ đầy bể). Thời gian vòi thứ hai chảy một mình đầy bể là:\par
\hspace*{1.5em}\textbf{A.} 18 giờ\quad\quad \textbf{B.} 12 giờ\quad\quad \textbf{C.} 9 giờ\quad\quad \textbf{D.} 24 giờ\par\vspace{2pt}
\textbf{\color{primary} Câu 48.} Tìm kết quả cuối cùng của bài toán phân loại học sinh khá giỏi (Chương I): Cho hệ phương trình x + y = m + 9 và x - y = m - 9. Tìm m để nghiệm (x; y) của hệ thỏa mãn x$\cdot$y = 18:\par
\hspace*{1.5em}\textbf{A.} m = 9\quad\quad \textbf{B.} m = 2\quad\quad \textbf{C.} m = 0\quad\quad \textbf{D.} m = -2\par\vspace{2pt}
\textbf{\color{primary} Câu 49.} Kết hợp các tính chất đã học trong Chương I, hãy giải quyết bài toán: Hai vòi nước cùng chảy vào một bể cạn thì sau 6 giờ đầy bể. Nếu vòi thứ nhất chảy trong 3 giờ và vòi thứ hai chảy trong 6 giờ thì được 3/4 bể (khi hai vòi cùng chảy 6 giờ đầy bể). Thời gian vòi thứ hai chảy một mình đầy bể là:\par
\hspace*{1.5em}\textbf{A.} 12 giờ\quad\quad \textbf{B.} 9 giờ\quad\quad \textbf{C.} 18 giờ\quad\quad \textbf{D.} 24 giờ\par\vspace{2pt}
\textbf{\color{primary} Câu 50.} Câu hỏi vận dụng cao cuối Chương I: Cho hệ phương trình x + y = m + 11 và x - y = m - 11. Tìm m để nghiệm (x; y) của hệ thỏa mãn x$\cdot$y = 22. Phương án đúng là:\par
\hspace*{1.5em}\textbf{A.} m = 2\quad\quad \textbf{B.} m = -2\quad\quad \textbf{C.} m = 0\quad\quad \textbf{D.} m = 11\par\vspace{2pt}
\subsection*{\color{primary} Bảng đáp án Chương I: Phương trình và hệ hai phương trình bậc nhất hai ẩn}
\begin{center}
\renewcommand{\arraystretch}{1.25}
\begin{tabularx}{\textwidth}{*{10}{>{\centering\arraybackslash}X}}
\toprule
\textbf{\color{primary} Cột 1} & \textbf{\color{primary} Cột 2} & \textbf{\color{primary} Cột 3} & \textbf{\color{primary} Cột 4} & \textbf{\color{primary} Cột 5} & \textbf{\color{primary} Cột 6} & \textbf{\color{primary} Cột 7} & \textbf{\color{primary} Cột 8} & \textbf{\color{primary} Cột 9} & \textbf{\color{primary} Cột 10} \\
\midrule
\textbf{1.D} & \textbf{2.B} & \textbf{3.A} & \textbf{4.C} & \textbf{5.C} & \textbf{6.C} & \textbf{7.B} & \textbf{8.B} & \textbf{9.A} & \textbf{10.C} \\
\rowcolor{rowtint} \textbf{11.B} & \textbf{12.B} & \textbf{13.C} & \textbf{14.C} & \textbf{15.B} & \textbf{16.D} & \textbf{17.C} & \textbf{18.A} & \textbf{19.D} & \textbf{20.D} \\
\textbf{21.A} & \textbf{22.C} & \textbf{23.A} & \textbf{24.A} & \textbf{25.C} & \textbf{26.A} & \textbf{27.C} & \textbf{28.A} & \textbf{29.D} & \textbf{30.B} \\
\rowcolor{rowtint} \textbf{31.A} & \textbf{32.D} & \textbf{33.C} & \textbf{34.C} & \textbf{35.C} & \textbf{36.D} & \textbf{37.B} & \textbf{38.B} & \textbf{39.B} & \textbf{40.A} \\
\textbf{41.A} & \textbf{42.D} & \textbf{43.B} & \textbf{44.D} & \textbf{45.B} & \textbf{46.D} & \textbf{47.B} & \textbf{48.B} & \textbf{49.A} & \textbf{50.A} \\
\bottomrule
\end{tabularx}
\end{center}\vspace{10pt}
\section*{\color{primary} CHƯƠNG II: PHƯƠNG TRÌNH VÀ BẤT PHƯƠNG TRÌNH BẬC NHẤT MỘT ẨN}
\textit{Trọng tâm kiến thức SGK: Phương trình quy về phương trình bậc nhất một ẩn (phương trình tích, phương trình chứa ẩn ở mẫu), bất đẳng thức và tính chất, bất phương trình bậc nhất một ẩn.}\par\vspace{4pt}
\subsection*{\color{accent} Phần I. Mức độ Dễ (Nhận biết -- Câu 1 đến Câu 20)}
\textbf{\color{primary} Câu 1.} Tập nghiệm của phương trình tích (x - 2)(x + 5) = 0 là:\par
\hspace*{1.5em}\textbf{A.} S = \{2; -5\}\quad\quad \textbf{B.} S = \{2\}\quad\quad \textbf{C.} S = \{-2; 5\}\quad\quad \textbf{D.} S = \{2; 5\}\par\vspace{2pt}
\textbf{\color{primary} Câu 2.} Thực hiện yêu cầu nhận biết sau: Điều kiện xác định của phương trình 3/(x - 2) + 2/(x + 4) = 1 là.\par
\hspace*{1.5em}\textbf{A.} x $\ne$ -2 và x $\ne$ 4\quad\quad \textbf{B.} x $\ne$ 2 và x $\ne$ -4\quad\quad \textbf{C.} x > 2\quad\quad \textbf{D.} x $\ne$ 2\par\vspace{2pt}
\textbf{\color{primary} Câu 3.} Trong bài học thuộc Chương II, xét nội dung: Nghiệm của bất phương trình 2x - 10 > 0 là. Đáp án đúng là:\par
\hspace*{1.5em}\textbf{A.} x $\ge$ 5\quad\quad \textbf{B.} x > -5\quad\quad \textbf{C.} x > 5\quad\quad \textbf{D.} x < 5\par\vspace{2pt}
\textbf{\color{primary} Câu 4.} Dựa vào kiến thức chuẩn SGK của Chương II (Phương trình và bất phương trình bậc nhất một ẩn), tập nghiệm của phương trình tích (x - 5)(x + 8) = 0 là:\par
\hspace*{1.5em}\textbf{A.} S = \{5; 8\}\quad\quad \textbf{B.} S = \{-5; 8\}\quad\quad \textbf{C.} S = \{5; -8\}\quad\quad \textbf{D.} S = \{5\}\par\vspace{2pt}
\textbf{\color{primary} Câu 5.} Chọn kết quả đúng nhất khi xét: Điều kiện xác định của phương trình 3/(x - 5) + 2/(x + 4) = 1 là:\par
\hspace*{1.5em}\textbf{A.} x $\ne$ 5\quad\quad \textbf{B.} x $\ne$ 5 và x $\ne$ -4\quad\quad \textbf{C.} x $\ne$ -5 và x $\ne$ 4\quad\quad \textbf{D.} x > 5\par\vspace{2pt}
\textbf{\color{primary} Câu 6.} Khi làm bài tập cơ bản về phương trình và bất phương trình bậc nhất một ẩn, một học sinh gặp câu hỏi: Nghiệm của bất phương trình 2x - 16 > 0 là. Kết quả nào dưới đây đúng?\par
\hspace*{1.5em}\textbf{A.} x $\ge$ 8\quad\quad \textbf{B.} x < 8\quad\quad \textbf{C.} x > 8\quad\quad \textbf{D.} x > -8\par\vspace{2pt}
\textbf{\color{primary} Câu 7.} Xác định giá trị hoặc khẳng định đúng cho bài toán sau: Tập nghiệm của phương trình tích (x - 8)(x + 11) = 0 là:\par
\hspace*{1.5em}\textbf{A.} S = \{8; 11\}\quad\quad \textbf{B.} S = \{8; -11\}\quad\quad \textbf{C.} S = \{8\}\quad\quad \textbf{D.} S = \{-8; 11\}\par\vspace{2pt}
\textbf{\color{primary} Câu 8.} Áp dụng trực tiếp định nghĩa và quy tắc trong Chương II: Điều kiện xác định của phương trình 3/(x - 8) + 2/(x + 4) = 1 là:\par
\hspace*{1.5em}\textbf{A.} x $\ne$ 8\quad\quad \textbf{B.} x $\ne$ -8 và x $\ne$ 4\quad\quad \textbf{C.} x > 8\quad\quad \textbf{D.} x $\ne$ 8 và x $\ne$ -4\par\vspace{2pt}
\textbf{\color{primary} Câu 9.} Kiểm tra nhanh kiến thức nền tảng: Nghiệm của bất phương trình 2x - 22 > 0 là:\par
\hspace*{1.5em}\textbf{A.} x < 11\quad\quad \textbf{B.} x > -11\quad\quad \textbf{C.} x $\ge$ 11\quad\quad \textbf{D.} x > 11\par\vspace{2pt}
\textbf{\color{primary} Câu 10.} Cho dữ kiện nhận biết sau đây: Tập nghiệm của phương trình tích (x - 11)(x + 14) = 0 là. Phương án chính xác là:\par
\hspace*{1.5em}\textbf{A.} S = \{-11; 14\}\quad\quad \textbf{B.} S = \{11\}\quad\quad \textbf{C.} S = \{11; 14\}\quad\quad \textbf{D.} S = \{11; -14\}\par\vspace{2pt}
\textbf{\color{primary} Câu 11.} Hãy chỉ ra kết quả đúng trong bốn lựa chọn A, B, C, D đối với yêu cầu: Điều kiện xác định của phương trình 3/(x - 11) + 2/(x + 4) = 1 là:\par
\hspace*{1.5em}\textbf{A.} x $\ne$ 11\quad\quad \textbf{B.} x $\ne$ 11 và x $\ne$ -4\quad\quad \textbf{C.} x > 11\quad\quad \textbf{D.} x $\ne$ -11 và x $\ne$ 4\par\vspace{2pt}
\textbf{\color{primary} Câu 12.} Câu hỏi ôn tập lý thuyết và kỹ năng cơ bản (Chương II): Nghiệm của bất phương trình 2x - 28 > 0 là:\par
\hspace*{1.5em}\textbf{A.} x < 14\quad\quad \textbf{B.} x > -14\quad\quad \textbf{C.} x $\ge$ 14\quad\quad \textbf{D.} x > 14\par\vspace{2pt}
\textbf{\color{primary} Câu 13.} Theo nội dung sách giáo khoa Kết nối tri thức với cuộc sống (Chương II), tập nghiệm của phương trình tích (x - 14)(x + 17) = 0 là:\par
\hspace*{1.5em}\textbf{A.} S = \{-14; 17\}\quad\quad \textbf{B.} S = \{14; 17\}\quad\quad \textbf{C.} S = \{14; -17\}\quad\quad \textbf{D.} S = \{14\}\par\vspace{2pt}
\textbf{\color{primary} Câu 14.} Tìm đáp số đúng của bài toán cơ bản sau: Điều kiện xác định của phương trình 3/(x - 14) + 2/(x + 4) = 1 là:\par
\hspace*{1.5em}\textbf{A.} x $\ne$ -14 và x $\ne$ 4\quad\quad \textbf{B.} x $\ne$ 14 và x $\ne$ -4\quad\quad \textbf{C.} x $\ne$ 14\quad\quad \textbf{D.} x > 14\par\vspace{2pt}
\textbf{\color{primary} Câu 15.} Xét đối tượng toán học trong Chương II: Nghiệm của bất phương trình 2x - 34 > 0 là. Giá trị hoặc tính chất tương ứng là:\par
\hspace*{1.5em}\textbf{A.} x $\ge$ 17\quad\quad \textbf{B.} x > 17\quad\quad \textbf{C.} x > -17\quad\quad \textbf{D.} x < 17\par\vspace{2pt}
\textbf{\color{primary} Câu 16.} Hoàn thành phép xác định cơ bản sau đây: Tập nghiệm của phương trình tích (x - 17)(x + 20) = 0 là:\par
\hspace*{1.5em}\textbf{A.} S = \{17; -20\}\quad\quad \textbf{B.} S = \{17; 20\}\quad\quad \textbf{C.} S = \{-17; 20\}\quad\quad \textbf{D.} S = \{17\}\par\vspace{2pt}
\textbf{\color{primary} Câu 17.} Trong các khẳng định hoặc giá trị dưới đây, đâu là đáp án đúng cho câu hỏi: Điều kiện xác định của phương trình 3/(x - 17) + 2/(x + 4) = 1 là?\par
\hspace*{1.5em}\textbf{A.} x $\ne$ 17 và x $\ne$ -4\quad\quad \textbf{B.} x $\ne$ 17\quad\quad \textbf{C.} x $\ne$ -17 và x $\ne$ 4\quad\quad \textbf{D.} x > 17\par\vspace{2pt}
\textbf{\color{primary} Câu 18.} Bài tập khởi động Chương II: Nghiệm của bất phương trình 2x - 40 > 0 là:\par
\hspace*{1.5em}\textbf{A.} x > -20\quad\quad \textbf{B.} x < 20\quad\quad \textbf{C.} x $\ge$ 20\quad\quad \textbf{D.} x > 20\par\vspace{2pt}
\textbf{\color{primary} Câu 19.} Quan sát dữ kiện đề bài: Tập nghiệm của phương trình tích (x - 20)(x + 23) = 0 là. Hãy chọn phương án đúng:\par
\hspace*{1.5em}\textbf{A.} S = \{20\}\quad\quad \textbf{B.} S = \{20; -23\}\quad\quad \textbf{C.} S = \{-20; 23\}\quad\quad \textbf{D.} S = \{20; 23\}\par\vspace{2pt}
\textbf{\color{primary} Câu 20.} Sử dụng công thức hoặc tính chất đầu bài của Chương II để giải quyết yêu cầu: Điều kiện xác định của phương trình 3/(x - 20) + 2/(x + 4) = 1 là:\par
\hspace*{1.5em}\textbf{A.} x $\ne$ 20 và x $\ne$ -4\quad\quad \textbf{B.} x $\ne$ -20 và x $\ne$ 4\quad\quad \textbf{C.} x $\ne$ 20\quad\quad \textbf{D.} x > 20\par\vspace{2pt}
\subsection*{\color{accent} Phần II. Mức độ Trung bình (Thông hiểu -- Câu 21 đến Câu 40)}
\textbf{\color{primary} Câu 21.} Giải bất phương trình: -3x + 12 $\le$ 0.:\par
\hspace*{1.5em}\textbf{A.} x $\ge$ 4\quad\quad \textbf{B.} x $\le$ 4\quad\quad \textbf{C.} x $\le$ -4\quad\quad \textbf{D.} x $\ge$ -4\par\vspace{2pt}
\textbf{\color{primary} Câu 22.} Vận dụng quy tắc trong Chương II (Phương trình và bất phương trình bậc nhất một ẩn) để giải bài toán: Cho a < b. Khẳng định nào sau đây là ĐÚNG:\par
\hspace*{1.5em}\textbf{A.} -2a + 4 > -2b + 4\quad\quad \textbf{B.} a - 4 > b - 4\quad\quad \textbf{C.} -2a + 4 < -2b + 4\quad\quad \textbf{D.} 2a + 4 > 2b + 4\par\vspace{2pt}
\textbf{\color{primary} Câu 23.} Biến đổi và tính toán theo kiến thức Chương II: Nghiệm của phương trình chứa ẩn ở mẫu (2x + 5)/(x - 1) = 3 (với x $\ne$ 1) là. Kết quả thu được là:\par
\hspace*{1.5em}\textbf{A.} x = 1\quad\quad \textbf{B.} x = -8\quad\quad \textbf{C.} x = 8\quad\quad \textbf{D.} x = 2\par\vspace{2pt}
\textbf{\color{primary} Câu 24.} Xét bài toán mức độ thông hiểu sau: Giải bất phương trình: -3x + 21 $\le$ 0.. Đáp án đúng là:\par
\hspace*{1.5em}\textbf{A.} x $\ge$ -7\quad\quad \textbf{B.} x $\le$ 7\quad\quad \textbf{C.} x $\le$ -7\quad\quad \textbf{D.} x $\ge$ 7\par\vspace{2pt}
\textbf{\color{primary} Câu 25.} Để hoàn thành phiếu học tập Chương II, hãy xác định: Cho a < b. Khẳng định nào sau đây là ĐÚNG:\par
\hspace*{1.5em}\textbf{A.} -2a + 7 > -2b + 7\quad\quad \textbf{B.} -2a + 7 < -2b + 7\quad\quad \textbf{C.} 2a + 7 > 2b + 7\quad\quad \textbf{D.} a - 7 > b - 7\par\vspace{2pt}
\textbf{\color{primary} Câu 26.} Phân tích dữ kiện bài toán: Nghiệm của phương trình chứa ẩn ở mẫu (2x + 8)/(x - 1) = 3 (với x $\ne$ 1) là. Giá trị cần tìm bằng:\par
\hspace*{1.5em}\textbf{A.} x = 11\quad\quad \textbf{B.} x = 5\quad\quad \textbf{C.} x = 1\quad\quad \textbf{D.} x = -11\par\vspace{2pt}
\textbf{\color{primary} Câu 27.} Áp dụng tính chất đã học ở Chương II vào bài tập sau: Giải bất phương trình: -3x + 30 $\le$ 0.:\par
\hspace*{1.5em}\textbf{A.} x $\ge$ 10\quad\quad \textbf{B.} x $\le$ -10\quad\quad \textbf{C.} x $\le$ 10\quad\quad \textbf{D.} x $\ge$ -10\par\vspace{2pt}
\textbf{\color{primary} Câu 28.} Thực hiện các bước suy luận hợp lý để giải quyết yêu cầu: Cho a < b. Khẳng định nào sau đây là ĐÚNG:\par
\hspace*{1.5em}\textbf{A.} a - 10 > b - 10\quad\quad \textbf{B.} 2a + 10 > 2b + 10\quad\quad \textbf{C.} -2a + 10 < -2b + 10\quad\quad \textbf{D.} -2a + 10 > -2b + 10\par\vspace{2pt}
\textbf{\color{primary} Câu 29.} Cho bài toán tự luyện (Phương trình và bất phương trình bậc nhất một ẩn): Nghiệm của phương trình chứa ẩn ở mẫu (2x + 11)/(x - 1) = 3 (với x $\ne$ 1) là. Phương án nào sau đây chính xác?\par
\hspace*{1.5em}\textbf{A.} x = 14\quad\quad \textbf{B.} x = 8\quad\quad \textbf{C.} x = -14\quad\quad \textbf{D.} x = 1\par\vspace{2pt}
\textbf{\color{primary} Câu 30.} Kết hợp các phép biến đổi trong Chương II, hãy tìm kết quả của bài toán: Giải bất phương trình: -3x + 39 $\le$ 0.:\par
\hspace*{1.5em}\textbf{A.} x $\ge$ 13\quad\quad \textbf{B.} x $\le$ 13\quad\quad \textbf{C.} x $\le$ -13\quad\quad \textbf{D.} x $\ge$ -13\par\vspace{2pt}
\textbf{\color{primary} Câu 31.} Một bài kiểm tra định kỳ nội dung Chương II nêu câu hỏi: Cho a < b. Khẳng định nào sau đây là ĐÚNG. Hãy chọn đáp số đúng:\par
\hspace*{1.5em}\textbf{A.} 2a + 13 > 2b + 13\quad\quad \textbf{B.} a - 13 > b - 13\quad\quad \textbf{C.} -2a + 13 > -2b + 13\quad\quad \textbf{D.} -2a + 13 < -2b + 13\par\vspace{2pt}
\textbf{\color{primary} Câu 32.} Tìm lời giải đúng cho bài toán thông hiểu: Nghiệm của phương trình chứa ẩn ở mẫu (2x + 14)/(x - 1) = 3 (với x $\ne$ 1) là:\par
\hspace*{1.5em}\textbf{A.} x = 11\quad\quad \textbf{B.} x = 17\quad\quad \textbf{C.} x = -17\quad\quad \textbf{D.} x = 1\par\vspace{2pt}
\textbf{\color{primary} Câu 33.} Sử dụng mối liên hệ giữa các đại lượng trong Chương II: Giải bất phương trình: -3x + 48 $\le$ 0.:\par
\hspace*{1.5em}\textbf{A.} x $\ge$ -16\quad\quad \textbf{B.} x $\le$ 16\quad\quad \textbf{C.} x $\le$ -16\quad\quad \textbf{D.} x $\ge$ 16\par\vspace{2pt}
\textbf{\color{primary} Câu 34.} Hãy tính toán và chọn phương án đúng đối với bài tập: Cho a < b. Khẳng định nào sau đây là ĐÚNG:\par
\hspace*{1.5em}\textbf{A.} a - 16 > b - 16\quad\quad \textbf{B.} -2a + 16 > -2b + 16\quad\quad \textbf{C.} 2a + 16 > 2b + 16\quad\quad \textbf{D.} -2a + 16 < -2b + 16\par\vspace{2pt}
\textbf{\color{primary} Câu 35.} Dựa trên các tính chất trọng tâm của Chương II, xét bài toán: Nghiệm của phương trình chứa ẩn ở mẫu (2x + 17)/(x - 1) = 3 (với x $\ne$ 1) là:\par
\hspace*{1.5em}\textbf{A.} x = 14\quad\quad \textbf{B.} x = -20\quad\quad \textbf{C.} x = 1\quad\quad \textbf{D.} x = 20\par\vspace{2pt}
\textbf{\color{primary} Câu 36.} Thực hành giải toán (Phương trình và bất phương trình bậc nhất một ẩn): Giải bất phương trình: -3x + 57 $\le$ 0.. Đáp án nào dưới đây thỏa mãn?\par
\hspace*{1.5em}\textbf{A.} x $\ge$ -19\quad\quad \textbf{B.} x $\ge$ 19\quad\quad \textbf{C.} x $\le$ 19\quad\quad \textbf{D.} x $\le$ -19\par\vspace{2pt}
\textbf{\color{primary} Câu 37.} Kiểm tra kỹ năng biến đổi và suy luận qua bài toán: Cho a < b. Khẳng định nào sau đây là ĐÚNG:\par
\hspace*{1.5em}\textbf{A.} -2a + 19 > -2b + 19\quad\quad \textbf{B.} a - 19 > b - 19\quad\quad \textbf{C.} 2a + 19 > 2b + 19\quad\quad \textbf{D.} -2a + 19 < -2b + 19\par\vspace{2pt}
\textbf{\color{primary} Câu 38.} Cho giả thiết thuộc phạm vi Chương II: Nghiệm của phương trình chứa ẩn ở mẫu (2x + 20)/(x - 1) = 3 (với x $\ne$ 1) là. Kết luận đúng là:\par
\hspace*{1.5em}\textbf{A.} x = 17\quad\quad \textbf{B.} x = 23\quad\quad \textbf{C.} x = 1\quad\quad \textbf{D.} x = -23\par\vspace{2pt}
\textbf{\color{primary} Câu 39.} Rút gọn hoặc tính giá trị theo yêu cầu sau của Chương II: Giải bất phương trình: -3x + 66 $\le$ 0.:\par
\hspace*{1.5em}\textbf{A.} x $\ge$ 22\quad\quad \textbf{B.} x $\le$ 22\quad\quad \textbf{C.} x $\le$ -22\quad\quad \textbf{D.} x $\ge$ -22\par\vspace{2pt}
\textbf{\color{primary} Câu 40.} Xác định phương án đúng sau khi thực hiện phép tính / suy luận: Cho a < b. Khẳng định nào sau đây là ĐÚNG:\par
\hspace*{1.5em}\textbf{A.} -2a + 22 > -2b + 22\quad\quad \textbf{B.} a - 22 > b - 22\quad\quad \textbf{C.} -2a + 22 < -2b + 22\quad\quad \textbf{D.} 2a + 22 > 2b + 22\par\vspace{2pt}
\subsection*{\color{accent} Phần III. Mức độ Khó (Vận dụng \& Vận dụng cao -- Câu 41 đến Câu 50)}
\textbf{\color{primary} Câu 41.} Số nghiệm nguyên dương của bất phương trình 3x - 5 $\le$ x + 7 là:\par
\hspace*{1.5em}\textbf{A.} 6 nghiệm (1)\quad\quad \textbf{B.} 6 nghiệm\quad\quad \textbf{C.} 7 nghiệm\quad\quad \textbf{D.} 5 nghiệm\par\vspace{2pt}
\textbf{\color{primary} Câu 42.} Bài toán vận dụng (Chương II – Phương trình và bất phương trình bậc nhất một ẩn): Cho hai số thực dương x, y thỏa mãn x + y = 10. Giá trị lớn nhất của tích P = x$\cdot$y bằng. Kết quả đúng là:\par
\hspace*{1.5em}\textbf{A.} 50\quad\quad \textbf{B.} 26\quad\quad \textbf{C.} 25\quad\quad \textbf{D.} 10\par\vspace{2pt}
\textbf{\color{primary} Câu 43.} Vận dụng kiến thức nâng cao trong phạm vi Chương II để giải bài toán: Số nghiệm nguyên dương của bất phương trình 3x - 5 $\le$ x + 11 là:\par
\hspace*{1.5em}\textbf{A.} 9 nghiệm\quad\quad \textbf{B.} 7 nghiệm\quad\quad \textbf{C.} 8 nghiệm\quad\quad \textbf{D.} 10 nghiệm\par\vspace{2pt}
\textbf{\color{primary} Câu 44.} Phân tích và giải bài toán tổng hợp sau của Chương II: Cho hai số thực dương x, y thỏa mãn x + y = 14. Giá trị lớn nhất của tích P = x$\cdot$y bằng:\par
\hspace*{1.5em}\textbf{A.} 14\quad\quad \textbf{B.} 50\quad\quad \textbf{C.} 98\quad\quad \textbf{D.} 49\par\vspace{2pt}
\textbf{\color{primary} Câu 45.} Thử thách tư duy trong phạm vi kiến thức Chương II: Số nghiệm nguyên dương của bất phương trình 3x - 5 $\le$ x + 15 là. Đáp số chính xác là:\par
\hspace*{1.5em}\textbf{A.} 10 nghiệm\quad\quad \textbf{B.} 9 nghiệm\quad\quad \textbf{C.} 11 nghiệm\quad\quad \textbf{D.} 14 nghiệm\par\vspace{2pt}
\textbf{\color{primary} Câu 46.} Sử dụng phương pháp giải toán của Chương II để tìm đáp án cho bài toán: Cho hai số thực dương x, y thỏa mãn x + y = 18. Giá trị lớn nhất của tích P = x$\cdot$y bằng:\par
\hspace*{1.5em}\textbf{A.} 162\quad\quad \textbf{B.} 18\quad\quad \textbf{C.} 82\quad\quad \textbf{D.} 81\par\vspace{2pt}
\textbf{\color{primary} Câu 47.} Xét bài toán vận dụng thực tiễn / nâng cao thuộc Chương II: Số nghiệm nguyên dương của bất phương trình 3x - 5 $\le$ x + 19 là:\par
\hspace*{1.5em}\textbf{A.} 11 nghiệm\quad\quad \textbf{B.} 13 nghiệm\quad\quad \textbf{C.} 18 nghiệm\quad\quad \textbf{D.} 12 nghiệm\par\vspace{2pt}
\textbf{\color{primary} Câu 48.} Tìm kết quả cuối cùng của bài toán phân loại học sinh khá giỏi (Chương II): Cho hai số thực dương x, y thỏa mãn x + y = 22. Giá trị lớn nhất của tích P = x$\cdot$y bằng:\par
\hspace*{1.5em}\textbf{A.} 22\quad\quad \textbf{B.} 121\quad\quad \textbf{C.} 242\quad\quad \textbf{D.} 122\par\vspace{2pt}
\textbf{\color{primary} Câu 49.} Kết hợp các tính chất đã học trong Chương II, hãy giải quyết bài toán: Số nghiệm nguyên dương của bất phương trình 3x - 5 $\le$ x + 23 là:\par
\hspace*{1.5em}\textbf{A.} 22 nghiệm\quad\quad \textbf{B.} 15 nghiệm\quad\quad \textbf{C.} 13 nghiệm\quad\quad \textbf{D.} 14 nghiệm\par\vspace{2pt}
\textbf{\color{primary} Câu 50.} Câu hỏi vận dụng cao cuối Chương II: Cho hai số thực dương x, y thỏa mãn x + y = 26. Giá trị lớn nhất của tích P = x$\cdot$y bằng. Phương án đúng là:\par
\hspace*{1.5em}\textbf{A.} 26\quad\quad \textbf{B.} 169\quad\quad \textbf{C.} 170\quad\quad \textbf{D.} 338\par\vspace{2pt}
\subsection*{\color{primary} Bảng đáp án Chương II: Phương trình và bất phương trình bậc nhất một ẩn}
\begin{center}
\renewcommand{\arraystretch}{1.25}
\begin{tabularx}{\textwidth}{*{10}{>{\centering\arraybackslash}X}}
\toprule
\textbf{\color{primary} Cột 1} & \textbf{\color{primary} Cột 2} & \textbf{\color{primary} Cột 3} & \textbf{\color{primary} Cột 4} & \textbf{\color{primary} Cột 5} & \textbf{\color{primary} Cột 6} & \textbf{\color{primary} Cột 7} & \textbf{\color{primary} Cột 8} & \textbf{\color{primary} Cột 9} & \textbf{\color{primary} Cột 10} \\
\midrule
\textbf{1.A} & \textbf{2.B} & \textbf{3.C} & \textbf{4.C} & \textbf{5.B} & \textbf{6.C} & \textbf{7.B} & \textbf{8.D} & \textbf{9.D} & \textbf{10.D} \\
\rowcolor{rowtint} \textbf{11.B} & \textbf{12.D} & \textbf{13.C} & \textbf{14.B} & \textbf{15.B} & \textbf{16.A} & \textbf{17.A} & \textbf{18.D} & \textbf{19.B} & \textbf{20.A} \\
\textbf{21.A} & \textbf{22.A} & \textbf{23.C} & \textbf{24.D} & \textbf{25.A} & \textbf{26.A} & \textbf{27.A} & \textbf{28.D} & \textbf{29.A} & \textbf{30.A} \\
\rowcolor{rowtint} \textbf{31.C} & \textbf{32.B} & \textbf{33.D} & \textbf{34.B} & \textbf{35.D} & \textbf{36.B} & \textbf{37.A} & \textbf{38.B} & \textbf{39.A} & \textbf{40.A} \\
\textbf{41.B} & \textbf{42.C} & \textbf{43.C} & \textbf{44.D} & \textbf{45.A} & \textbf{46.D} & \textbf{47.D} & \textbf{48.B} & \textbf{49.D} & \textbf{50.B} \\
\bottomrule
\end{tabularx}
\end{center}\vspace{10pt}
\section*{\color{primary} CHƯƠNG III: CĂN BẬC HAI VÀ CĂN BẬC BA}
\textit{Trọng tâm kiến thức SGK: Căn bậc hai và căn thức bậc hai, điều kiện xác định, khai căn bậc hai với phép nhân và phép chia, biến đổi đơn giản và rút gọn biểu thức chứa căn thức bậc hai, căn bậc ba.}\par\vspace{4pt}
\subsection*{\color{accent} Phần I. Mức độ Dễ (Nhận biết -- Câu 1 đến Câu 20)}
\textbf{\color{primary} Câu 1.} Điều kiện xác định của căn thức $\sqrt{x - 3}$ là:\par
\hspace*{1.5em}\textbf{A.} x $\ne$ 3\quad\quad \textbf{B.} x $\ge$ 3\quad\quad \textbf{C.} x > 3\quad\quad \textbf{D.} x $\le$ 3\par\vspace{2pt}
\textbf{\color{primary} Câu 2.} Thực hiện yêu cầu nhận biết sau: Giá trị của căn bậc ba $\sqrt[3]{-27}$ bằng.\par
\hspace*{1.5em}\textbf{A.} Không tồn tại\quad\quad \textbf{B.} 3\quad\quad \textbf{C.} -3\quad\quad \textbf{D.} -9\par\vspace{2pt}
\textbf{\color{primary} Câu 3.} Trong bài học thuộc Chương III, xét nội dung: Đưa thừa số ra ngoài dấu căn của $\sqrt{16  \cdot  5}$ ta được. Đáp án đúng là:\par
\hspace*{1.5em}\textbf{A.} 8$\sqrt{5}$\quad\quad \textbf{B.} 4$\sqrt{5}$\quad\quad \textbf{C.} 16$\sqrt{5}$\quad\quad \textbf{D.} 5$\sqrt{4}$\par\vspace{2pt}
\textbf{\color{primary} Câu 4.} Dựa vào kiến thức chuẩn SGK của Chương III (Căn bậc hai và căn bậc ba), điều kiện xác định của căn thức $\sqrt{x - 6}$ là:\par
\hspace*{1.5em}\textbf{A.} x $\ne$ 6\quad\quad \textbf{B.} x $\ge$ 6\quad\quad \textbf{C.} x $\le$ 6\quad\quad \textbf{D.} x > 6\par\vspace{2pt}
\textbf{\color{primary} Câu 5.} Chọn kết quả đúng nhất khi xét: Giá trị của căn bậc ba $\sqrt[3]{-216}$ bằng:\par
\hspace*{1.5em}\textbf{A.} -6\quad\quad \textbf{B.} 6\quad\quad \textbf{C.} -36\quad\quad \textbf{D.} Không tồn tại\par\vspace{2pt}
\textbf{\color{primary} Câu 6.} Khi làm bài tập cơ bản về căn bậc hai và căn bậc ba, một học sinh gặp câu hỏi: Đưa thừa số ra ngoài dấu căn của $\sqrt{49  \cdot  5}$ ta được. Kết quả nào dưới đây đúng?\par
\hspace*{1.5em}\textbf{A.} 5$\sqrt{7}$\quad\quad \textbf{B.} 14$\sqrt{5}$\quad\quad \textbf{C.} 7$\sqrt{5}$\quad\quad \textbf{D.} 49$\sqrt{5}$\par\vspace{2pt}
\textbf{\color{primary} Câu 7.} Xác định giá trị hoặc khẳng định đúng cho bài toán sau: Điều kiện xác định của căn thức $\sqrt{x - 9}$ là:\par
\hspace*{1.5em}\textbf{A.} x $\ne$ 9\quad\quad \textbf{B.} x $\le$ 9\quad\quad \textbf{C.} x $\ge$ 9\quad\quad \textbf{D.} x > 9\par\vspace{2pt}
\textbf{\color{primary} Câu 8.} Áp dụng trực tiếp định nghĩa và quy tắc trong Chương III: Giá trị của căn bậc ba $\sqrt[3]{-729}$ bằng:\par
\hspace*{1.5em}\textbf{A.} -81\quad\quad \textbf{B.} 9\quad\quad \textbf{C.} -9\quad\quad \textbf{D.} Không tồn tại\par\vspace{2pt}
\textbf{\color{primary} Câu 9.} Kiểm tra nhanh kiến thức nền tảng: Đưa thừa số ra ngoài dấu căn của $\sqrt{100  \cdot  5}$ ta được:\par
\hspace*{1.5em}\textbf{A.} 5$\sqrt{10}$\quad\quad \textbf{B.} 10$\sqrt{5}$\quad\quad \textbf{C.} 20$\sqrt{5}$\quad\quad \textbf{D.} 100$\sqrt{5}$\par\vspace{2pt}
\textbf{\color{primary} Câu 10.} Cho dữ kiện nhận biết sau đây: Điều kiện xác định của căn thức $\sqrt{x - 12}$ là. Phương án chính xác là:\par
\hspace*{1.5em}\textbf{A.} x $\le$ 12\quad\quad \textbf{B.} x $\ne$ 12\quad\quad \textbf{C.} x > 12\quad\quad \textbf{D.} x $\ge$ 12\par\vspace{2pt}
\textbf{\color{primary} Câu 11.} Hãy chỉ ra kết quả đúng trong bốn lựa chọn A, B, C, D đối với yêu cầu: Giá trị của căn bậc ba $\sqrt[3]{-1728}$ bằng:\par
\hspace*{1.5em}\textbf{A.} 12\quad\quad \textbf{B.} -144\quad\quad \textbf{C.} -12\quad\quad \textbf{D.} Không tồn tại\par\vspace{2pt}
\textbf{\color{primary} Câu 12.} Câu hỏi ôn tập lý thuyết và kỹ năng cơ bản (Chương III): Đưa thừa số ra ngoài dấu căn của $\sqrt{169  \cdot  5}$ ta được:\par
\hspace*{1.5em}\textbf{A.} 5$\sqrt{13}$\quad\quad \textbf{B.} 13$\sqrt{5}$\quad\quad \textbf{C.} 26$\sqrt{5}$\quad\quad \textbf{D.} 169$\sqrt{5}$\par\vspace{2pt}
\textbf{\color{primary} Câu 13.} Theo nội dung sách giáo khoa Kết nối tri thức với cuộc sống (Chương III), điều kiện xác định của căn thức $\sqrt{x - 15}$ là:\par
\hspace*{1.5em}\textbf{A.} x $\le$ 15\quad\quad \textbf{B.} x > 15\quad\quad \textbf{C.} x $\ne$ 15\quad\quad \textbf{D.} x $\ge$ 15\par\vspace{2pt}
\textbf{\color{primary} Câu 14.} Tìm đáp số đúng của bài toán cơ bản sau: Giá trị của căn bậc ba $\sqrt[3]{-3375}$ bằng:\par
\hspace*{1.5em}\textbf{A.} -15\quad\quad \textbf{B.} 15\quad\quad \textbf{C.} Không tồn tại\quad\quad \textbf{D.} -225\par\vspace{2pt}
\textbf{\color{primary} Câu 15.} Xét đối tượng toán học trong Chương III: Đưa thừa số ra ngoài dấu căn của $\sqrt{256  \cdot  5}$ ta được. Giá trị hoặc tính chất tương ứng là:\par
\hspace*{1.5em}\textbf{A.} 5$\sqrt{16}$\quad\quad \textbf{B.} 256$\sqrt{5}$\quad\quad \textbf{C.} 32$\sqrt{5}$\quad\quad \textbf{D.} 16$\sqrt{5}$\par\vspace{2pt}
\textbf{\color{primary} Câu 16.} Hoàn thành phép xác định cơ bản sau đây: Điều kiện xác định của căn thức $\sqrt{x - 18}$ là:\par
\hspace*{1.5em}\textbf{A.} x > 18\quad\quad \textbf{B.} x $\le$ 18\quad\quad \textbf{C.} x $\ne$ 18\quad\quad \textbf{D.} x $\ge$ 18\par\vspace{2pt}
\textbf{\color{primary} Câu 17.} Trong các khẳng định hoặc giá trị dưới đây, đâu là đáp án đúng cho câu hỏi: Giá trị của căn bậc ba $\sqrt[3]{-5832}$ bằng?\par
\hspace*{1.5em}\textbf{A.} -324\quad\quad \textbf{B.} Không tồn tại\quad\quad \textbf{C.} 18\quad\quad \textbf{D.} -18\par\vspace{2pt}
\textbf{\color{primary} Câu 18.} Bài tập khởi động Chương III: Đưa thừa số ra ngoài dấu căn của $\sqrt{361  \cdot  5}$ ta được:\par
\hspace*{1.5em}\textbf{A.} 38$\sqrt{5}$\quad\quad \textbf{B.} 361$\sqrt{5}$\quad\quad \textbf{C.} 5$\sqrt{19}$\quad\quad \textbf{D.} 19$\sqrt{5}$\par\vspace{2pt}
\textbf{\color{primary} Câu 19.} Quan sát dữ kiện đề bài: Điều kiện xác định của căn thức $\sqrt{x - 21}$ là. Hãy chọn phương án đúng:\par
\hspace*{1.5em}\textbf{A.} x > 21\quad\quad \textbf{B.} x $\ne$ 21\quad\quad \textbf{C.} x $\le$ 21\quad\quad \textbf{D.} x $\ge$ 21\par\vspace{2pt}
\textbf{\color{primary} Câu 20.} Sử dụng công thức hoặc tính chất đầu bài của Chương III để giải quyết yêu cầu: Giá trị của căn bậc ba $\sqrt[3]{-9261}$ bằng:\par
\hspace*{1.5em}\textbf{A.} Không tồn tại\quad\quad \textbf{B.} -21\quad\quad \textbf{C.} -441\quad\quad \textbf{D.} 21\par\vspace{2pt}
\subsection*{\color{accent} Phần II. Mức độ Trung bình (Thông hiểu -- Câu 21 đến Câu 40)}
\textbf{\color{primary} Câu 21.} Rút gọn biểu thức A = $\sqrt{(\sqrt{5} - 3)^2}$ (biết $\sqrt{5}$ < 3) ta được:\par
\hspace*{1.5em}\textbf{A.} 3 - $\sqrt{5}$\quad\quad \textbf{B.} $\sqrt{5}$ - 3\quad\quad \textbf{C.} 3 + $\sqrt{5}$\quad\quad \textbf{D.} 4\par\vspace{2pt}
\textbf{\color{primary} Câu 22.} Vận dụng quy tắc trong Chương III (Căn bậc hai và căn bậc ba) để giải bài toán: Thu gọn biểu thức B = 3$\sqrt{3}$ + 4$\sqrt{12}$ - $\sqrt{27}$ ta được:\par
\hspace*{1.5em}\textbf{A.} 6$\sqrt{3}$\quad\quad \textbf{B.} 7$\sqrt{3}$\quad\quad \textbf{C.} 8$\sqrt{3}$\quad\quad \textbf{D.} 10$\sqrt{3}$\par\vspace{2pt}
\textbf{\color{primary} Câu 23.} Biến đổi và tính toán theo kiến thức Chương III: Giải phương trình $\sqrt{2x - 1}$ = 5 (với x $\ge$ 1/2). Kết quả thu được là:\par
\hspace*{1.5em}\textbf{A.} x = 13\quad\quad \textbf{B.} x = 4\quad\quad \textbf{C.} x = 25\quad\quad \textbf{D.} x = 12.0\par\vspace{2pt}
\textbf{\color{primary} Câu 24.} Xét bài toán mức độ thông hiểu sau: Rút gọn biểu thức A = $\sqrt{(\sqrt{8} - 6)^2}$ (biết $\sqrt{8}$ < 6) ta được. Đáp án đúng là:\par
\hspace*{1.5em}\textbf{A.} 6 + $\sqrt{8}$\quad\quad \textbf{B.} 6 - $\sqrt{8}$\quad\quad \textbf{C.} 28\quad\quad \textbf{D.} $\sqrt{8}$ - 6\par\vspace{2pt}
\textbf{\color{primary} Câu 25.} Để hoàn thành phiếu học tập Chương III, hãy xác định: Thu gọn biểu thức B = 3$\sqrt{2}$ + 4$\sqrt{8}$ - $\sqrt{18}$ ta được:\par
\hspace*{1.5em}\textbf{A.} 7$\sqrt{2}$\quad\quad \textbf{B.} 8$\sqrt{2}$\quad\quad \textbf{C.} 10$\sqrt{2}$\quad\quad \textbf{D.} 6$\sqrt{2}$\par\vspace{2pt}
\textbf{\color{primary} Câu 26.} Phân tích dữ kiện bài toán: Giải phương trình $\sqrt{2x - 1}$ = 8 (với x $\ge$ 1/2). Giá trị cần tìm bằng:\par
\hspace*{1.5em}\textbf{A.} x = 7\quad\quad \textbf{B.} x = 64\quad\quad \textbf{C.} x = 32.5\quad\quad \textbf{D.} x = 31.5\par\vspace{2pt}
\textbf{\color{primary} Câu 27.} Áp dụng tính chất đã học ở Chương III vào bài tập sau: Rút gọn biểu thức A = $\sqrt{(\sqrt{11} - 9)^2}$ (biết $\sqrt{11}$ < 9) ta được:\par
\hspace*{1.5em}\textbf{A.} 9 + $\sqrt{11}$\quad\quad \textbf{B.} $\sqrt{11}$ - 9\quad\quad \textbf{C.} 9 - $\sqrt{11}$\quad\quad \textbf{D.} 70\par\vspace{2pt}
\textbf{\color{primary} Câu 28.} Thực hiện các bước suy luận hợp lý để giải quyết yêu cầu: Thu gọn biểu thức B = 3$\sqrt{3}$ + 4$\sqrt{12}$ - $\sqrt{27}$ ta được:\par
\hspace*{1.5em}\textbf{A.} 6$\sqrt{3}$\quad\quad \textbf{B.} 7$\sqrt{3}$\quad\quad \textbf{C.} 8$\sqrt{3}$\quad\quad \textbf{D.} 10$\sqrt{3}$\par\vspace{2pt}
\textbf{\color{primary} Câu 29.} Cho bài toán tự luyện (Căn bậc hai và căn bậc ba): Giải phương trình $\sqrt{2x - 1}$ = 11 (với x $\ge$ 1/2). Phương án nào sau đây chính xác?\par
\hspace*{1.5em}\textbf{A.} x = 60.0\quad\quad \textbf{B.} x = 61\quad\quad \textbf{C.} x = 10\quad\quad \textbf{D.} x = 121\par\vspace{2pt}
\textbf{\color{primary} Câu 30.} Kết hợp các phép biến đổi trong Chương III, hãy tìm kết quả của bài toán: Rút gọn biểu thức A = $\sqrt{(\sqrt{14} - 12)^2}$ (biết $\sqrt{14}$ < 12) ta được:\par
\hspace*{1.5em}\textbf{A.} 12 - $\sqrt{14}$\quad\quad \textbf{B.} 130\quad\quad \textbf{C.} $\sqrt{14}$ - 12\quad\quad \textbf{D.} 12 + $\sqrt{14}$\par\vspace{2pt}
\textbf{\color{primary} Câu 31.} Một bài kiểm tra định kỳ nội dung Chương III nêu câu hỏi: Thu gọn biểu thức B = 3$\sqrt{2}$ + 4$\sqrt{8}$ - $\sqrt{18}$ ta được. Hãy chọn đáp số đúng:\par
\hspace*{1.5em}\textbf{A.} 7$\sqrt{2}$\quad\quad \textbf{B.} 10$\sqrt{2}$\quad\quad \textbf{C.} 6$\sqrt{2}$\quad\quad \textbf{D.} 8$\sqrt{2}$\par\vspace{2pt}
\textbf{\color{primary} Câu 32.} Tìm lời giải đúng cho bài toán thông hiểu: Giải phương trình $\sqrt{2x - 1}$ = 14 (với x $\ge$ 1/2):\par
\hspace*{1.5em}\textbf{A.} x = 98.5\quad\quad \textbf{B.} x = 13\quad\quad \textbf{C.} x = 97.5\quad\quad \textbf{D.} x = 196\par\vspace{2pt}
\textbf{\color{primary} Câu 33.} Sử dụng mối liên hệ giữa các đại lượng trong Chương III: Rút gọn biểu thức A = $\sqrt{(\sqrt{17} - 15)^2}$ (biết $\sqrt{17}$ < 15) ta được:\par
\hspace*{1.5em}\textbf{A.} 15 + $\sqrt{17}$\quad\quad \textbf{B.} 15 - $\sqrt{17}$\quad\quad \textbf{C.} 208\quad\quad \textbf{D.} $\sqrt{17}$ - 15\par\vspace{2pt}
\textbf{\color{primary} Câu 34.} Hãy tính toán và chọn phương án đúng đối với bài tập: Thu gọn biểu thức B = 3$\sqrt{3}$ + 4$\sqrt{12}$ - $\sqrt{27}$ ta được:\par
\hspace*{1.5em}\textbf{A.} 7$\sqrt{3}$\quad\quad \textbf{B.} 6$\sqrt{3}$\quad\quad \textbf{C.} 10$\sqrt{3}$\quad\quad \textbf{D.} 8$\sqrt{3}$\par\vspace{2pt}
\textbf{\color{primary} Câu 35.} Dựa trên các tính chất trọng tâm của Chương III, xét bài toán: Giải phương trình $\sqrt{2x - 1}$ = 17 (với x $\ge$ 1/2):\par
\hspace*{1.5em}\textbf{A.} x = 16\quad\quad \textbf{B.} x = 145\quad\quad \textbf{C.} x = 289\quad\quad \textbf{D.} x = 144.0\par\vspace{2pt}
\textbf{\color{primary} Câu 36.} Thực hành giải toán (Căn bậc hai và căn bậc ba): Rút gọn biểu thức A = $\sqrt{(\sqrt{20} - 18)^2}$ (biết $\sqrt{20}$ < 18) ta được. Đáp án nào dưới đây thỏa mãn?\par
\hspace*{1.5em}\textbf{A.} $\sqrt{20}$ - 18\quad\quad \textbf{B.} 18 - $\sqrt{20}$\quad\quad \textbf{C.} 304\quad\quad \textbf{D.} 18 + $\sqrt{20}$\par\vspace{2pt}
\textbf{\color{primary} Câu 37.} Kiểm tra kỹ năng biến đổi và suy luận qua bài toán: Thu gọn biểu thức B = 3$\sqrt{2}$ + 4$\sqrt{8}$ - $\sqrt{18}$ ta được:\par
\hspace*{1.5em}\textbf{A.} 8$\sqrt{2}$\quad\quad \textbf{B.} 6$\sqrt{2}$\quad\quad \textbf{C.} 7$\sqrt{2}$\quad\quad \textbf{D.} 10$\sqrt{2}$\par\vspace{2pt}
\textbf{\color{primary} Câu 38.} Cho giả thiết thuộc phạm vi Chương III: Giải phương trình $\sqrt{2x - 1}$ = 20 (với x $\ge$ 1/2). Kết luận đúng là:\par
\hspace*{1.5em}\textbf{A.} x = 400\quad\quad \textbf{B.} x = 199.5\quad\quad \textbf{C.} x = 200.5\quad\quad \textbf{D.} x = 19\par\vspace{2pt}
\textbf{\color{primary} Câu 39.} Rút gọn hoặc tính giá trị theo yêu cầu sau của Chương III: Rút gọn biểu thức A = $\sqrt{(\sqrt{23} - 21)^2}$ (biết $\sqrt{23}$ < 21) ta được:\par
\hspace*{1.5em}\textbf{A.} 21 + $\sqrt{23}$\quad\quad \textbf{B.} 21 - $\sqrt{23}$\quad\quad \textbf{C.} 418\quad\quad \textbf{D.} $\sqrt{23}$ - 21\par\vspace{2pt}
\textbf{\color{primary} Câu 40.} Xác định phương án đúng sau khi thực hiện phép tính / suy luận: Thu gọn biểu thức B = 3$\sqrt{3}$ + 4$\sqrt{12}$ - $\sqrt{27}$ ta được:\par
\hspace*{1.5em}\textbf{A.} 7$\sqrt{3}$\quad\quad \textbf{B.} 6$\sqrt{3}$\quad\quad \textbf{C.} 10$\sqrt{3}$\quad\quad \textbf{D.} 8$\sqrt{3}$\par\vspace{2pt}
\subsection*{\color{accent} Phần III. Mức độ Khó (Vận dụng \& Vận dụng cao -- Câu 41 đến Câu 50)}
\textbf{\color{primary} Câu 41.} Với x $\ge$ 0 và x $\ne$ 4, rút gọn biểu thức P = (x - 4) / ($\sqrt{x}$ - 2) ta được:\par
\hspace*{1.5em}\textbf{A.} x + 2\quad\quad \textbf{B.} 1\quad\quad \textbf{C.} $\sqrt{x}$ + 2\quad\quad \textbf{D.} $\sqrt{x}$ - 2\par\vspace{2pt}
\textbf{\color{primary} Câu 42.} Bài toán vận dụng (Chương III – Căn bậc hai và căn bậc ba): Tìm giá trị lớn nhất của biểu thức Q = 13 / ($\sqrt{x}$ + 4) với x $\ge$ 0. Kết quả đúng là:\par
\hspace*{1.5em}\textbf{A.} 13\quad\quad \textbf{B.} 1/4\quad\quad \textbf{C.} 0\quad\quad \textbf{D.} 13/4\par\vspace{2pt}
\textbf{\color{primary} Câu 43.} Vận dụng kiến thức nâng cao trong phạm vi Chương III để giải bài toán: Với x $\ge$ 0 và x $\ne$ 16, rút gọn biểu thức P = (x - 16) / ($\sqrt{x}$ - 4) ta được:\par
\hspace*{1.5em}\textbf{A.} 1\quad\quad \textbf{B.} $\sqrt{x}$ - 4\quad\quad \textbf{C.} x + 4\quad\quad \textbf{D.} $\sqrt{x}$ + 4\par\vspace{2pt}
\textbf{\color{primary} Câu 44.} Phân tích và giải bài toán tổng hợp sau của Chương III: Tìm giá trị lớn nhất của biểu thức Q = 15 / ($\sqrt{x}$ + 6) với x $\ge$ 0:\par
\hspace*{1.5em}\textbf{A.} 15\quad\quad \textbf{B.} 1/6\quad\quad \textbf{C.} 15/6\quad\quad \textbf{D.} 0\par\vspace{2pt}
\textbf{\color{primary} Câu 45.} Thử thách tư duy trong phạm vi kiến thức Chương III: Với x $\ge$ 0 và x $\ne$ 36, rút gọn biểu thức P = (x - 36) / ($\sqrt{x}$ - 6) ta được. Đáp số chính xác là:\par
\hspace*{1.5em}\textbf{A.} 1\quad\quad \textbf{B.} $\sqrt{x}$ + 6\quad\quad \textbf{C.} x + 6\quad\quad \textbf{D.} $\sqrt{x}$ - 6\par\vspace{2pt}
\textbf{\color{primary} Câu 46.} Sử dụng phương pháp giải toán của Chương III để tìm đáp án cho bài toán: Tìm giá trị lớn nhất của biểu thức Q = 17 / ($\sqrt{x}$ + 8) với x $\ge$ 0:\par
\hspace*{1.5em}\textbf{A.} 17/8\quad\quad \textbf{B.} 17\quad\quad \textbf{C.} 0\quad\quad \textbf{D.} 1/8\par\vspace{2pt}
\textbf{\color{primary} Câu 47.} Xét bài toán vận dụng thực tiễn / nâng cao thuộc Chương III: Với x $\ge$ 0 và x $\ne$ 64, rút gọn biểu thức P = (x - 64) / ($\sqrt{x}$ - 8) ta được:\par
\hspace*{1.5em}\textbf{A.} x + 8\quad\quad \textbf{B.} $\sqrt{x}$ + 8\quad\quad \textbf{C.} 1\quad\quad \textbf{D.} $\sqrt{x}$ - 8\par\vspace{2pt}
\textbf{\color{primary} Câu 48.} Tìm kết quả cuối cùng của bài toán phân loại học sinh khá giỏi (Chương III): Tìm giá trị lớn nhất của biểu thức Q = 19 / ($\sqrt{x}$ + 10) với x $\ge$ 0:\par
\hspace*{1.5em}\textbf{A.} 19\quad\quad \textbf{B.} 1/10\quad\quad \textbf{C.} 0\quad\quad \textbf{D.} 19/10\par\vspace{2pt}
\textbf{\color{primary} Câu 49.} Kết hợp các tính chất đã học trong Chương III, hãy giải quyết bài toán: Với x $\ge$ 0 và x $\ne$ 100, rút gọn biểu thức P = (x - 100) / ($\sqrt{x}$ - 10) ta được:\par
\hspace*{1.5em}\textbf{A.} x + 10\quad\quad \textbf{B.} $\sqrt{x}$ + 10\quad\quad \textbf{C.} 1\quad\quad \textbf{D.} $\sqrt{x}$ - 10\par\vspace{2pt}
\textbf{\color{primary} Câu 50.} Câu hỏi vận dụng cao cuối Chương III: Tìm giá trị lớn nhất của biểu thức Q = 21 / ($\sqrt{x}$ + 12) với x $\ge$ 0. Phương án đúng là:\par
\hspace*{1.5em}\textbf{A.} 0\quad\quad \textbf{B.} 21/12\quad\quad \textbf{C.} 1/12\quad\quad \textbf{D.} 21\par\vspace{2pt}
\subsection*{\color{primary} Bảng đáp án Chương III: Căn bậc hai và căn bậc ba}
\begin{center}
\renewcommand{\arraystretch}{1.25}
\begin{tabularx}{\textwidth}{*{10}{>{\centering\arraybackslash}X}}
\toprule
\textbf{\color{primary} Cột 1} & \textbf{\color{primary} Cột 2} & \textbf{\color{primary} Cột 3} & \textbf{\color{primary} Cột 4} & \textbf{\color{primary} Cột 5} & \textbf{\color{primary} Cột 6} & \textbf{\color{primary} Cột 7} & \textbf{\color{primary} Cột 8} & \textbf{\color{primary} Cột 9} & \textbf{\color{primary} Cột 10} \\
\midrule
\textbf{1.B} & \textbf{2.C} & \textbf{3.B} & \textbf{4.B} & \textbf{5.A} & \textbf{6.C} & \textbf{7.C} & \textbf{8.C} & \textbf{9.B} & \textbf{10.D} \\
\rowcolor{rowtint} \textbf{11.C} & \textbf{12.B} & \textbf{13.D} & \textbf{14.A} & \textbf{15.D} & \textbf{16.D} & \textbf{17.D} & \textbf{18.D} & \textbf{19.D} & \textbf{20.B} \\
\textbf{21.A} & \textbf{22.C} & \textbf{23.A} & \textbf{24.B} & \textbf{25.B} & \textbf{26.C} & \textbf{27.C} & \textbf{28.C} & \textbf{29.B} & \textbf{30.A} \\
\rowcolor{rowtint} \textbf{31.D} & \textbf{32.A} & \textbf{33.B} & \textbf{34.D} & \textbf{35.B} & \textbf{36.B} & \textbf{37.A} & \textbf{38.C} & \textbf{39.B} & \textbf{40.D} \\
\textbf{41.C} & \textbf{42.D} & \textbf{43.D} & \textbf{44.C} & \textbf{45.B} & \textbf{46.A} & \textbf{47.B} & \textbf{48.D} & \textbf{49.B} & \textbf{50.B} \\
\bottomrule
\end{tabularx}
\end{center}\vspace{10pt}
\section*{\color{primary} CHƯƠNG IV: HỆ THỨC LƯỢNG TRONG TAM GIÁC VUÔNG}
\textit{Trọng tâm kiến thức SGK: Tỉ số lượng giác của góc nhọn (sin, cos, tan, cot), tỉ số lượng giác của hai góc phụ nhau, một số hệ thức giữa cạnh và góc trong tam giác vuông và ứng dụng thực tế.}\par\vspace{4pt}
\subsection*{\color{accent} Phần I. Mức độ Dễ (Nhận biết -- Câu 1 đến Câu 20)}
\textbf{\color{primary} Câu 1.} Cho hai góc nhọn $\alpha$ và $\beta$ phụ nhau ($\alpha$ + $\beta$ = 90°). Hệ thức nào sau đây là ĐÚNG?\par
\hspace*{1.5em}\textbf{A.} sin $\alpha$ = cos $\beta$\quad\quad \textbf{B.} tan $\alpha$ = tan $\beta$\quad\quad \textbf{C.} sin $\alpha$ = sin $\beta$\quad\quad \textbf{D.} sin $\alpha$ = cot $\beta$\par\vspace{2pt}
\textbf{\color{primary} Câu 2.} Thực hiện yêu cầu nhận biết sau: Cho tam giác ABC vuông tại A có BC = 10 cm, AC = 6 cm. Giá trị của sin B bằng.\par
\hspace*{1.5em}\textbf{A.} 4/5\quad\quad \textbf{B.} 3/4\quad\quad \textbf{C.} 3/5\quad\quad \textbf{D.} 5/3\par\vspace{2pt}
\textbf{\color{primary} Câu 3.} Trong bài học thuộc Chương IV, xét nội dung: Giá trị của biểu thức lượng giác sin$^2$ 35° + cos$^2$ 35° bằng. Đáp án đúng là:\par
\hspace*{1.5em}\textbf{A.} 1/2\quad\quad \textbf{B.} 1\quad\quad \textbf{C.} 2\quad\quad \textbf{D.} 0\par\vspace{2pt}
\textbf{\color{primary} Câu 4.} Dựa vào kiến thức chuẩn SGK của Chương IV (Hệ thức lượng trong tam giác vuông), cho góc nhọn $\alpha$ = 21°. Khi đó tỉ số lượng giác sin 21° bằng tỉ số lượng giác nào dưới đây:\par
\hspace*{1.5em}\textbf{A.} cos 21°\quad\quad \textbf{B.} cos 69°\quad\quad \textbf{C.} sin 69°\quad\quad \textbf{D.} tan 69°\par\vspace{2pt}
\textbf{\color{primary} Câu 5.} Chọn kết quả đúng nhất khi xét: Trong tam giác ABC vuông tại A có cạnh huyền BC = 30 cm và cạnh AB = 24 cm. Giá trị của cos B = AB/BC bằng:\par
\hspace*{1.5em}\textbf{A.} 3/5\quad\quad \textbf{B.} 3/4\quad\quad \textbf{C.} 4/5\quad\quad \textbf{D.} 5/4\par\vspace{2pt}
\textbf{\color{primary} Câu 6.} Khi làm bài tập cơ bản về hệ thức lượng trong tam giác vuông, một học sinh gặp câu hỏi: Cho hai góc nhọn $\alpha$ và $\beta$ phụ nhau ($\alpha$ + $\beta$ = 90°). Hệ thức nào sau đây là ĐÚNG. Kết quả nào dưới đây đúng?\par
\hspace*{1.5em}\textbf{A.} sin $\alpha$ = cos $\beta$\quad\quad \textbf{B.} tan $\alpha$ = tan $\beta$\quad\quad \textbf{C.} sin $\alpha$ = sin $\beta$\quad\quad \textbf{D.} sin $\alpha$ = cot $\beta$\par\vspace{2pt}
\textbf{\color{primary} Câu 7.} Xác định giá trị hoặc khẳng định đúng cho bài toán sau: Cho tam giác ABC vuông tại A có BC = 5 cm, AC = 3 cm. Giá trị của sin B bằng:\par
\hspace*{1.5em}\textbf{A.} 5/3\quad\quad \textbf{B.} 4/5\quad\quad \textbf{C.} 3/5\quad\quad \textbf{D.} 3/4\par\vspace{2pt}
\textbf{\color{primary} Câu 8.} Áp dụng trực tiếp định nghĩa và quy tắc trong Chương IV: Giá trị của biểu thức lượng giác sin$^2$ 35° + cos$^2$ 35° bằng:\par
\hspace*{1.5em}\textbf{A.} 0\quad\quad \textbf{B.} 1\quad\quad \textbf{C.} 1/2\quad\quad \textbf{D.} 2\par\vspace{2pt}
\textbf{\color{primary} Câu 9.} Kiểm tra nhanh kiến thức nền tảng: Cho góc nhọn $\alpha$ = 31°. Khi đó tỉ số lượng giác sin 31° bằng tỉ số lượng giác nào dưới đây?\par
\hspace*{1.5em}\textbf{A.} cos 31°\quad\quad \textbf{B.} sin 59°\quad\quad \textbf{C.} tan 59°\quad\quad \textbf{D.} cos 59°\par\vspace{2pt}
\textbf{\color{primary} Câu 10.} Cho dữ kiện nhận biết sau đây: Trong tam giác ABC vuông tại A có cạnh huyền BC = 55 cm và cạnh AB = 44 cm. Giá trị của cos B = AB/BC bằng. Phương án chính xác là:\par
\hspace*{1.5em}\textbf{A.} 5/4\quad\quad \textbf{B.} 3/4\quad\quad \textbf{C.} 4/5\quad\quad \textbf{D.} 3/5\par\vspace{2pt}
\textbf{\color{primary} Câu 11.} Hãy chỉ ra kết quả đúng trong bốn lựa chọn A, B, C, D đối với yêu cầu: Cho hai góc nhọn $\alpha$ và $\beta$ phụ nhau ($\alpha$ + $\beta$ = 90°). Hệ thức nào sau đây là ĐÚNG:\par
\hspace*{1.5em}\textbf{A.} sin $\alpha$ = cos $\beta$\quad\quad \textbf{B.} sin $\alpha$ = sin $\beta$\quad\quad \textbf{C.} sin $\alpha$ = cot $\beta$\quad\quad \textbf{D.} tan $\alpha$ = tan $\beta$\par\vspace{2pt}
\textbf{\color{primary} Câu 12.} Câu hỏi ôn tập lý thuyết và kỹ năng cơ bản (Chương IV): Cho tam giác ABC vuông tại A có BC = 15 cm, AC = 9 cm. Giá trị của sin B bằng:\par
\hspace*{1.5em}\textbf{A.} 3/5\quad\quad \textbf{B.} 4/5\quad\quad \textbf{C.} 3/4\quad\quad \textbf{D.} 5/3\par\vspace{2pt}
\textbf{\color{primary} Câu 13.} Theo nội dung sách giáo khoa Kết nối tri thức với cuộc sống (Chương IV), giá trị của biểu thức lượng giác sin$^2$ 35° + cos$^2$ 35° bằng:\par
\hspace*{1.5em}\textbf{A.} 0\quad\quad \textbf{B.} 1/2\quad\quad \textbf{C.} 2\quad\quad \textbf{D.} 1\par\vspace{2pt}
\textbf{\color{primary} Câu 14.} Tìm đáp số đúng của bài toán cơ bản sau: Cho góc nhọn $\alpha$ = 41°. Khi đó tỉ số lượng giác sin 41° bằng tỉ số lượng giác nào dưới đây:\par
\hspace*{1.5em}\textbf{A.} cos 49°\quad\quad \textbf{B.} tan 49°\quad\quad \textbf{C.} cos 41°\quad\quad \textbf{D.} sin 49°\par\vspace{2pt}
\textbf{\color{primary} Câu 15.} Xét đối tượng toán học trong Chương IV: Trong tam giác ABC vuông tại A có cạnh huyền BC = 80 cm và cạnh AB = 64 cm. Giá trị của cos B = AB/BC bằng. Giá trị hoặc tính chất tương ứng là:\par
\hspace*{1.5em}\textbf{A.} 5/4\quad\quad \textbf{B.} 3/5\quad\quad \textbf{C.} 4/5\quad\quad \textbf{D.} 3/4\par\vspace{2pt}
\textbf{\color{primary} Câu 16.} Hoàn thành phép xác định cơ bản sau đây: Cho hai góc nhọn $\alpha$ và $\beta$ phụ nhau ($\alpha$ + $\beta$ = 90°). Hệ thức nào sau đây là ĐÚNG:\par
\hspace*{1.5em}\textbf{A.} sin $\alpha$ = sin $\beta$\quad\quad \textbf{B.} sin $\alpha$ = cos $\beta$\quad\quad \textbf{C.} sin $\alpha$ = cot $\beta$\quad\quad \textbf{D.} tan $\alpha$ = tan $\beta$\par\vspace{2pt}
\textbf{\color{primary} Câu 17.} Trong các khẳng định hoặc giá trị dưới đây, đâu là đáp án đúng cho câu hỏi: Cho tam giác ABC vuông tại A có BC = 10 cm, AC = 6 cm. Giá trị của sin B bằng?\par
\hspace*{1.5em}\textbf{A.} 4/5\quad\quad \textbf{B.} 3/4\quad\quad \textbf{C.} 3/5\quad\quad \textbf{D.} 5/3\par\vspace{2pt}
\textbf{\color{primary} Câu 18.} Bài tập khởi động Chương IV: Giá trị của biểu thức lượng giác sin$^2$ 35° + cos$^2$ 35° bằng:\par
\hspace*{1.5em}\textbf{A.} 1/2\quad\quad \textbf{B.} 1\quad\quad \textbf{C.} 2\quad\quad \textbf{D.} 0\par\vspace{2pt}
\textbf{\color{primary} Câu 19.} Quan sát dữ kiện đề bài: Cho góc nhọn $\alpha$ = 51°. Khi đó tỉ số lượng giác sin 51° bằng tỉ số lượng giác nào dưới đây. Hãy chọn phương án đúng:\par
\hspace*{1.5em}\textbf{A.} cos 51°\quad\quad \textbf{B.} tan 39°\quad\quad \textbf{C.} cos 39°\quad\quad \textbf{D.} sin 39°\par\vspace{2pt}
\textbf{\color{primary} Câu 20.} Sử dụng công thức hoặc tính chất đầu bài của Chương IV để giải quyết yêu cầu: Trong tam giác ABC vuông tại A có cạnh huyền BC = 105 cm và cạnh AB = 84 cm. Giá trị của cos B = AB/BC bằng:\par
\hspace*{1.5em}\textbf{A.} 5/4\quad\quad \textbf{B.} 3/5\quad\quad \textbf{C.} 3/4\quad\quad \textbf{D.} 4/5\par\vspace{2pt}
\subsection*{\color{accent} Phần II. Mức độ Trung bình (Thông hiểu -- Câu 21 đến Câu 40)}
\textbf{\color{primary} Câu 21.} Cho tam giác ABC vuông tại A có cạnh huyền BC = 8 cm và góc B = 60°. Độ dài cạnh AB bằng:\par
\hspace*{1.5em}\textbf{A.} 4$\sqrt{2}$ cm\quad\quad \textbf{B.} 4$\sqrt{3}$ cm\quad\quad \textbf{C.} 4 cm\quad\quad \textbf{D.} 8 cm\par\vspace{2pt}
\textbf{\color{primary} Câu 22.} Vận dụng quy tắc trong Chương IV (Hệ thức lượng trong tam giác vuông) để giải bài toán: Cho tam giác ABC vuông tại A có góc C = 30° và cạnh đối AB = 6 cm. Độ dài cạnh kề AC bằng:\par
\hspace*{1.5em}\textbf{A.} 12 cm\quad\quad \textbf{B.} 6$\sqrt{3}$ cm\quad\quad \textbf{C.} 6/$\sqrt{3}$ cm\quad\quad \textbf{D.} 6 cm\par\vspace{2pt}
\textbf{\color{primary} Câu 23.} Biến đổi và tính toán theo kiến thức Chương IV: Rút gọn biểu thức A = sin 22° - cos 68° + tan 32° $\cdot$ cot 32° ta được. Kết quả thu được là:\par
\hspace*{1.5em}\textbf{A.} 2\quad\quad \textbf{B.} 0\quad\quad \textbf{C.} -1\quad\quad \textbf{D.} 1\par\vspace{2pt}
\textbf{\color{primary} Câu 24.} Xét bài toán mức độ thông hiểu sau: Cho tam giác ABC vuông tại A có cạnh huyền BC = 14 cm và góc B = 60°. Độ dài cạnh AB bằng. Đáp án đúng là:\par
\hspace*{1.5em}\textbf{A.} 14 cm\quad\quad \textbf{B.} 7$\sqrt{3}$ cm\quad\quad \textbf{C.} 7$\sqrt{2}$ cm\quad\quad \textbf{D.} 7 cm\par\vspace{2pt}
\textbf{\color{primary} Câu 25.} Để hoàn thành phiếu học tập Chương IV, hãy xác định: Cho tam giác ABC vuông tại A có góc C = 30° và cạnh đối AB = 9 cm. Độ dài cạnh kề AC bằng:\par
\hspace*{1.5em}\textbf{A.} 18 cm\quad\quad \textbf{B.} 9 cm\quad\quad \textbf{C.} 9$\sqrt{3}$ cm\quad\quad \textbf{D.} 9/$\sqrt{3}$ cm\par\vspace{2pt}
\textbf{\color{primary} Câu 26.} Phân tích dữ kiện bài toán: Rút gọn biểu thức A = sin 25° - cos 65° + tan 35° $\cdot$ cot 35° ta được. Giá trị cần tìm bằng:\par
\hspace*{1.5em}\textbf{A.} 2\quad\quad \textbf{B.} 1\quad\quad \textbf{C.} -1\quad\quad \textbf{D.} 0\par\vspace{2pt}
\textbf{\color{primary} Câu 27.} Áp dụng tính chất đã học ở Chương IV vào bài tập sau: Cho tam giác ABC vuông tại A có cạnh huyền BC = 20 cm và góc B = 60°. Độ dài cạnh AB bằng:\par
\hspace*{1.5em}\textbf{A.} 20 cm\quad\quad \textbf{B.} 10 cm\quad\quad \textbf{C.} 10$\sqrt{2}$ cm\quad\quad \textbf{D.} 10$\sqrt{3}$ cm\par\vspace{2pt}
\textbf{\color{primary} Câu 28.} Thực hiện các bước suy luận hợp lý để giải quyết yêu cầu: Cho tam giác ABC vuông tại A có góc C = 30° và cạnh đối AB = 12 cm. Độ dài cạnh kề AC bằng:\par
\hspace*{1.5em}\textbf{A.} 12 cm\quad\quad \textbf{B.} 24 cm\quad\quad \textbf{C.} 12$\sqrt{3}$ cm\quad\quad \textbf{D.} 12/$\sqrt{3}$ cm\par\vspace{2pt}
\textbf{\color{primary} Câu 29.} Cho bài toán tự luyện (Hệ thức lượng trong tam giác vuông): Rút gọn biểu thức A = sin 28° - cos 62° + tan 38° $\cdot$ cot 38° ta được. Phương án nào sau đây chính xác?\par
\hspace*{1.5em}\textbf{A.} 2\quad\quad \textbf{B.} -1\quad\quad \textbf{C.} 0\quad\quad \textbf{D.} 1\par\vspace{2pt}
\textbf{\color{primary} Câu 30.} Kết hợp các phép biến đổi trong Chương IV, hãy tìm kết quả của bài toán: Cho tam giác ABC vuông tại A có cạnh huyền BC = 26 cm và góc B = 60°. Độ dài cạnh AB bằng:\par
\hspace*{1.5em}\textbf{A.} 13 cm\quad\quad \textbf{B.} 13$\sqrt{2}$ cm\quad\quad \textbf{C.} 26 cm\quad\quad \textbf{D.} 13$\sqrt{3}$ cm\par\vspace{2pt}
\textbf{\color{primary} Câu 31.} Một bài kiểm tra định kỳ nội dung Chương IV nêu câu hỏi: Cho tam giác ABC vuông tại A có góc C = 30° và cạnh đối AB = 15 cm. Độ dài cạnh kề AC bằng. Hãy chọn đáp số đúng:\par
\hspace*{1.5em}\textbf{A.} 30 cm\quad\quad \textbf{B.} 15$\sqrt{3}$ cm\quad\quad \textbf{C.} 15/$\sqrt{3}$ cm\quad\quad \textbf{D.} 15 cm\par\vspace{2pt}
\textbf{\color{primary} Câu 32.} Tìm lời giải đúng cho bài toán thông hiểu: Rút gọn biểu thức A = sin 31° - cos 59° + tan 41° $\cdot$ cot 41° ta được:\par
\hspace*{1.5em}\textbf{A.} 1\quad\quad \textbf{B.} -1\quad\quad \textbf{C.} 2\quad\quad \textbf{D.} 0\par\vspace{2pt}
\textbf{\color{primary} Câu 33.} Sử dụng mối liên hệ giữa các đại lượng trong Chương IV: Cho tam giác ABC vuông tại A có cạnh huyền BC = 32 cm và góc B = 60°. Độ dài cạnh AB bằng:\par
\hspace*{1.5em}\textbf{A.} 16$\sqrt{2}$ cm\quad\quad \textbf{B.} 32 cm\quad\quad \textbf{C.} 16$\sqrt{3}$ cm\quad\quad \textbf{D.} 16 cm\par\vspace{2pt}
\textbf{\color{primary} Câu 34.} Hãy tính toán và chọn phương án đúng đối với bài tập: Cho tam giác ABC vuông tại A có góc C = 30° và cạnh đối AB = 18 cm. Độ dài cạnh kề AC bằng:\par
\hspace*{1.5em}\textbf{A.} 18 cm\quad\quad \textbf{B.} 18/$\sqrt{3}$ cm\quad\quad \textbf{C.} 36 cm\quad\quad \textbf{D.} 18$\sqrt{3}$ cm\par\vspace{2pt}
\textbf{\color{primary} Câu 35.} Dựa trên các tính chất trọng tâm của Chương IV, xét bài toán: Rút gọn biểu thức A = sin 34° - cos 56° + tan 44° $\cdot$ cot 44° ta được:\par
\hspace*{1.5em}\textbf{A.} 0\quad\quad \textbf{B.} -1\quad\quad \textbf{C.} 1\quad\quad \textbf{D.} 2\par\vspace{2pt}
\textbf{\color{primary} Câu 36.} Thực hành giải toán (Hệ thức lượng trong tam giác vuông): Cho tam giác ABC vuông tại A có cạnh huyền BC = 38 cm và góc B = 60°. Độ dài cạnh AB bằng. Đáp án nào dưới đây thỏa mãn?\par
\hspace*{1.5em}\textbf{A.} 19 cm\quad\quad \textbf{B.} 38 cm\quad\quad \textbf{C.} 19$\sqrt{2}$ cm\quad\quad \textbf{D.} 19$\sqrt{3}$ cm\par\vspace{2pt}
\textbf{\color{primary} Câu 37.} Kiểm tra kỹ năng biến đổi và suy luận qua bài toán: Cho tam giác ABC vuông tại A có góc C = 30° và cạnh đối AB = 21 cm. Độ dài cạnh kề AC bằng:\par
\hspace*{1.5em}\textbf{A.} 42 cm\quad\quad \textbf{B.} 21$\sqrt{3}$ cm\quad\quad \textbf{C.} 21 cm\quad\quad \textbf{D.} 21/$\sqrt{3}$ cm\par\vspace{2pt}
\textbf{\color{primary} Câu 38.} Cho giả thiết thuộc phạm vi Chương IV: Rút gọn biểu thức A = sin 37° - cos 53° + tan 47° $\cdot$ cot 47° ta được. Kết luận đúng là:\par
\hspace*{1.5em}\textbf{A.} 1\quad\quad \textbf{B.} 0\quad\quad \textbf{C.} -1\quad\quad \textbf{D.} 2\par\vspace{2pt}
\textbf{\color{primary} Câu 39.} Rút gọn hoặc tính giá trị theo yêu cầu sau của Chương IV: Cho tam giác ABC vuông tại A có cạnh huyền BC = 44 cm và góc B = 60°. Độ dài cạnh AB bằng:\par
\hspace*{1.5em}\textbf{A.} 44 cm\quad\quad \textbf{B.} 22$\sqrt{3}$ cm\quad\quad \textbf{C.} 22 cm\quad\quad \textbf{D.} 22$\sqrt{2}$ cm\par\vspace{2pt}
\textbf{\color{primary} Câu 40.} Xác định phương án đúng sau khi thực hiện phép tính / suy luận: Cho tam giác ABC vuông tại A có góc C = 30° và cạnh đối AB = 24 cm. Độ dài cạnh kề AC bằng:\par
\hspace*{1.5em}\textbf{A.} 24$\sqrt{3}$ cm\quad\quad \textbf{B.} 48 cm\quad\quad \textbf{C.} 24 cm\quad\quad \textbf{D.} 24/$\sqrt{3}$ cm\par\vspace{2pt}
\subsection*{\color{accent} Phần III. Mức độ Khó (Vận dụng \& Vận dụng cao -- Câu 41 đến Câu 50)}
\textbf{\color{primary} Câu 41.} Một chiếc thang dài 6 m đặt tựa vào bức tường thẳng đứng tạo với mặt đất nằm ngang một góc 60°. Khoảng cách từ chân thang đến chân tường bằng:\par
\hspace*{1.5em}\textbf{A.} 3$\sqrt{3}$ m\quad\quad \textbf{B.} 6 m\quad\quad \textbf{C.} 3$\sqrt{2}$ m\quad\quad \textbf{D.} 3 m\par\vspace{2pt}
\textbf{\color{primary} Câu 42.} Bài toán vận dụng (Chương IV – Hệ thức lượng trong tam giác vuông): Cho góc nhọn $\alpha$ thỏa mãn sin $\alpha$ = 3/5. Giá trị của biểu thức P = 5$\cdot$cos $\alpha$ + 4$\cdot$tan $\alpha$ bằng. Kết quả đúng là:\par
\hspace*{1.5em}\textbf{A.} 6\quad\quad \textbf{B.} 8\quad\quad \textbf{C.} 7\quad\quad \textbf{D.} 5\par\vspace{2pt}
\textbf{\color{primary} Câu 43.} Vận dụng kiến thức nâng cao trong phạm vi Chương IV để giải bài toán: Một chiếc thang dài 10 m đặt tựa vào bức tường thẳng đứng tạo với mặt đất nằm ngang một góc 60°. Khoảng cách từ chân thang đến chân tường bằng:\par
\hspace*{1.5em}\textbf{A.} 10 m\quad\quad \textbf{B.} 5 m\quad\quad \textbf{C.} 5$\sqrt{2}$ m\quad\quad \textbf{D.} 5$\sqrt{3}$ m\par\vspace{2pt}
\textbf{\color{primary} Câu 44.} Phân tích và giải bài toán tổng hợp sau của Chương IV: Cho góc nhọn $\alpha$ thỏa mãn sin $\alpha$ = 3/5. Giá trị của biểu thức P = 5$\cdot$cos $\alpha$ + 4$\cdot$tan $\alpha$ bằng:\par
\hspace*{1.5em}\textbf{A.} 5\quad\quad \textbf{B.} 8\quad\quad \textbf{C.} 6\quad\quad \textbf{D.} 7\par\vspace{2pt}
\textbf{\color{primary} Câu 45.} Thử thách tư duy trong phạm vi kiến thức Chương IV: Một chiếc thang dài 14 m đặt tựa vào bức tường thẳng đứng tạo với mặt đất nằm ngang một góc 60°. Khoảng cách từ chân thang đến chân tường bằng. Đáp số chính xác là:\par
\hspace*{1.5em}\textbf{A.} 14 m\quad\quad \textbf{B.} 7 m\quad\quad \textbf{C.} 7$\sqrt{2}$ m\quad\quad \textbf{D.} 7$\sqrt{3}$ m\par\vspace{2pt}
\textbf{\color{primary} Câu 46.} Sử dụng phương pháp giải toán của Chương IV để tìm đáp án cho bài toán: Cho góc nhọn $\alpha$ thỏa mãn sin $\alpha$ = 3/5. Giá trị của biểu thức P = 5$\cdot$cos $\alpha$ + 4$\cdot$tan $\alpha$ bằng:\par
\hspace*{1.5em}\textbf{A.} 5\quad\quad \textbf{B.} 7\quad\quad \textbf{C.} 6\quad\quad \textbf{D.} 8\par\vspace{2pt}
\textbf{\color{primary} Câu 47.} Xét bài toán vận dụng thực tiễn / nâng cao thuộc Chương IV: Một chiếc thang dài 18 m đặt tựa vào bức tường thẳng đứng tạo với mặt đất nằm ngang một góc 60°. Khoảng cách từ chân thang đến chân tường bằng:\par
\hspace*{1.5em}\textbf{A.} 9$\sqrt{3}$ m\quad\quad \textbf{B.} 9$\sqrt{2}$ m\quad\quad \textbf{C.} 18 m\quad\quad \textbf{D.} 9 m\par\vspace{2pt}
\textbf{\color{primary} Câu 48.} Tìm kết quả cuối cùng của bài toán phân loại học sinh khá giỏi (Chương IV): Cho góc nhọn $\alpha$ thỏa mãn sin $\alpha$ = 3/5. Giá trị của biểu thức P = 5$\cdot$cos $\alpha$ + 4$\cdot$tan $\alpha$ bằng:\par
\hspace*{1.5em}\textbf{A.} 7\quad\quad \textbf{B.} 8\quad\quad \textbf{C.} 5\quad\quad \textbf{D.} 6\par\vspace{2pt}
\textbf{\color{primary} Câu 49.} Kết hợp các tính chất đã học trong Chương IV, hãy giải quyết bài toán: Một chiếc thang dài 22 m đặt tựa vào bức tường thẳng đứng tạo với mặt đất nằm ngang một góc 60°. Khoảng cách từ chân thang đến chân tường bằng:\par
\hspace*{1.5em}\textbf{A.} 11$\sqrt{3}$ m\quad\quad \textbf{B.} 11$\sqrt{2}$ m\quad\quad \textbf{C.} 11 m\quad\quad \textbf{D.} 22 m\par\vspace{2pt}
\textbf{\color{primary} Câu 50.} Câu hỏi vận dụng cao cuối Chương IV: Cho góc nhọn $\alpha$ thỏa mãn sin $\alpha$ = 3/5. Giá trị của biểu thức P = 5$\cdot$cos $\alpha$ + 4$\cdot$tan $\alpha$ bằng. Phương án đúng là:\par
\hspace*{1.5em}\textbf{A.} 7\quad\quad \textbf{B.} 6\quad\quad \textbf{C.} 5\quad\quad \textbf{D.} 8\par\vspace{2pt}
\subsection*{\color{primary} Bảng đáp án Chương IV: Hệ thức lượng trong tam giác vuông}
\begin{center}
\renewcommand{\arraystretch}{1.25}
\begin{tabularx}{\textwidth}{*{10}{>{\centering\arraybackslash}X}}
\toprule
\textbf{\color{primary} Cột 1} & \textbf{\color{primary} Cột 2} & \textbf{\color{primary} Cột 3} & \textbf{\color{primary} Cột 4} & \textbf{\color{primary} Cột 5} & \textbf{\color{primary} Cột 6} & \textbf{\color{primary} Cột 7} & \textbf{\color{primary} Cột 8} & \textbf{\color{primary} Cột 9} & \textbf{\color{primary} Cột 10} \\
\midrule
\textbf{1.A} & \textbf{2.C} & \textbf{3.B} & \textbf{4.B} & \textbf{5.C} & \textbf{6.A} & \textbf{7.C} & \textbf{8.B} & \textbf{9.D} & \textbf{10.C} \\
\rowcolor{rowtint} \textbf{11.A} & \textbf{12.A} & \textbf{13.D} & \textbf{14.A} & \textbf{15.C} & \textbf{16.B} & \textbf{17.C} & \textbf{18.B} & \textbf{19.C} & \textbf{20.D} \\
\textbf{21.C} & \textbf{22.B} & \textbf{23.D} & \textbf{24.D} & \textbf{25.C} & \textbf{26.B} & \textbf{27.B} & \textbf{28.C} & \textbf{29.D} & \textbf{30.A} \\
\rowcolor{rowtint} \textbf{31.B} & \textbf{32.A} & \textbf{33.D} & \textbf{34.D} & \textbf{35.C} & \textbf{36.A} & \textbf{37.B} & \textbf{38.A} & \textbf{39.C} & \textbf{40.A} \\
\textbf{41.D} & \textbf{42.C} & \textbf{43.B} & \textbf{44.D} & \textbf{45.B} & \textbf{46.B} & \textbf{47.D} & \textbf{48.A} & \textbf{49.C} & \textbf{50.A} \\
\bottomrule
\end{tabularx}
\end{center}\vspace{10pt}
\section*{\color{primary} CHƯƠNG V: ĐƯỜNG TRÒN}
\textit{Trọng tâm kiến thức SGK: Mở đầu về đường tròn, dây và đường kính của đường tròn, góc ở tâm, góc nội tiếp, độ dài của cung tròn, diện tích hình quạt tròn và hình vành khuyên, vị trí tương đối của đường thẳng và hai đường tròn.}\par\vspace{4pt}
\subsection*{\color{accent} Phần I. Mức độ Dễ (Nhận biết -- Câu 1 đến Câu 20)}
\textbf{\color{primary} Câu 1.} Trong đường tròn (O; 5 cm), độ dài của dây cung lớn nhất (đường kính) bằng:\par
\hspace*{1.5em}\textbf{A.} 25$\pi$ cm\quad\quad \textbf{B.} 10 cm\quad\quad \textbf{C.} 5 cm\quad\quad \textbf{D.} 15 cm\par\vspace{2pt}
\textbf{\color{primary} Câu 2.} Thực hiện yêu cầu nhận biết sau: Trong một đường tròn, góc ở tâm chắn cung AB có số đo bằng 56°. Góc nội tiếp cùng chắn cung AB đó có số đo bằng.\par
\hspace*{1.5em}\textbf{A.} 90°\quad\quad \textbf{B.} 112°\quad\quad \textbf{C.} 28°\quad\quad \textbf{D.} 56°\par\vspace{2pt}
\textbf{\color{primary} Câu 3.} Trong bài học thuộc Chương V, xét nội dung: Cho đường tròn (O; 6 cm) và đường thẳng d có khoảng cách từ tâm O đến d là OH = 6 cm. Vị trí tương đối của d và (O) là. Đáp án đúng là:\par
\begin{itemize}\setlength{\itemsep}{1pt}\setlength{\parskip}{0pt}
  \item[] \textbf{A.} Đường thẳng d đi qua tâm O
  \item[] \textbf{B.} Đường thẳng d tiếp xúc với đường tròn (O)
  \item[] \textbf{C.} Đường thẳng d và (O) không giao nhau
  \item[] \textbf{D.} Đường thẳng d cắt đường tròn (O) tại hai điểm
\end{itemize}\vspace{2pt}
\textbf{\color{primary} Câu 4.} Dựa vào kiến thức chuẩn SGK của Chương V (Đường tròn), trong đường tròn (O; 8 cm), độ dài của dây cung lớn nhất (đường kính) bằng:\par
\hspace*{1.5em}\textbf{A.} 8 cm\quad\quad \textbf{B.} 24 cm\quad\quad \textbf{C.} 16 cm\quad\quad \textbf{D.} 64$\pi$ cm\par\vspace{2pt}
\textbf{\color{primary} Câu 5.} Chọn kết quả đúng nhất khi xét: Trong một đường tròn, góc ở tâm chắn cung AB có số đo bằng 74°. Góc nội tiếp cùng chắn cung AB đó có số đo bằng:\par
\hspace*{1.5em}\textbf{A.} 90°\quad\quad \textbf{B.} 37°\quad\quad \textbf{C.} 148°\quad\quad \textbf{D.} 74°\par\vspace{2pt}
\textbf{\color{primary} Câu 6.} Khi làm bài tập cơ bản về đường tròn, một học sinh gặp câu hỏi: Cho đường tròn (O; 9 cm) và đường thẳng d có khoảng cách từ tâm O đến d là OH = 9 cm. Vị trí tương đối của d và (O) là. Kết quả nào dưới đây đúng?\par
\begin{itemize}\setlength{\itemsep}{1pt}\setlength{\parskip}{0pt}
  \item[] \textbf{A.} Đường thẳng d tiếp xúc với đường tròn (O)
  \item[] \textbf{B.} Đường thẳng d cắt đường tròn (O) tại hai điểm
  \item[] \textbf{C.} Đường thẳng d và (O) không giao nhau
  \item[] \textbf{D.} Đường thẳng d đi qua tâm O
\end{itemize}\vspace{2pt}
\textbf{\color{primary} Câu 7.} Xác định giá trị hoặc khẳng định đúng cho bài toán sau: Trong đường tròn (O; 11 cm), độ dài của dây cung lớn nhất (đường kính) bằng:\par
\hspace*{1.5em}\textbf{A.} 121$\pi$ cm\quad\quad \textbf{B.} 33 cm\quad\quad \textbf{C.} 11 cm\quad\quad \textbf{D.} 22 cm\par\vspace{2pt}
\textbf{\color{primary} Câu 8.} Áp dụng trực tiếp định nghĩa và quy tắc trong Chương V: Trong một đường tròn, góc ở tâm chắn cung AB có số đo bằng 92°. Góc nội tiếp cùng chắn cung AB đó có số đo bằng:\par
\hspace*{1.5em}\textbf{A.} 92°\quad\quad \textbf{B.} 46°\quad\quad \textbf{C.} 90°\quad\quad \textbf{D.} 184°\par\vspace{2pt}
\textbf{\color{primary} Câu 9.} Kiểm tra nhanh kiến thức nền tảng: Cho đường tròn (O; 12 cm) và đường thẳng d có khoảng cách từ tâm O đến d là OH = 12 cm. Vị trí tương đối của d và (O) là:\par
\begin{itemize}\setlength{\itemsep}{1pt}\setlength{\parskip}{0pt}
  \item[] \textbf{A.} Đường thẳng d tiếp xúc với đường tròn (O)
  \item[] \textbf{B.} Đường thẳng d và (O) không giao nhau
  \item[] \textbf{C.} Đường thẳng d đi qua tâm O
  \item[] \textbf{D.} Đường thẳng d cắt đường tròn (O) tại hai điểm
\end{itemize}\vspace{2pt}
\textbf{\color{primary} Câu 10.} Cho dữ kiện nhận biết sau đây: Trong đường tròn (O; 14 cm), độ dài của dây cung lớn nhất (đường kính) bằng. Phương án chính xác là:\par
\hspace*{1.5em}\textbf{A.} 14 cm\quad\quad \textbf{B.} 196$\pi$ cm\quad\quad \textbf{C.} 42 cm\quad\quad \textbf{D.} 28 cm\par\vspace{2pt}
\textbf{\color{primary} Câu 11.} Hãy chỉ ra kết quả đúng trong bốn lựa chọn A, B, C, D đối với yêu cầu: Trong một đường tròn, góc ở tâm chắn cung AB có số đo bằng 110°. Góc nội tiếp cùng chắn cung AB đó có số đo bằng:\par
\hspace*{1.5em}\textbf{A.} 55°\quad\quad \textbf{B.} 110°\quad\quad \textbf{C.} 220°\quad\quad \textbf{D.} 90°\par\vspace{2pt}
\textbf{\color{primary} Câu 12.} Câu hỏi ôn tập lý thuyết và kỹ năng cơ bản (Chương V): Cho đường tròn (O; 15 cm) và đường thẳng d có khoảng cách từ tâm O đến d là OH = 15 cm. Vị trí tương đối của d và (O) là:\par
\begin{itemize}\setlength{\itemsep}{1pt}\setlength{\parskip}{0pt}
  \item[] \textbf{A.} Đường thẳng d đi qua tâm O
  \item[] \textbf{B.} Đường thẳng d và (O) không giao nhau
  \item[] \textbf{C.} Đường thẳng d tiếp xúc với đường tròn (O)
  \item[] \textbf{D.} Đường thẳng d cắt đường tròn (O) tại hai điểm
\end{itemize}\vspace{2pt}
\textbf{\color{primary} Câu 13.} Theo nội dung sách giáo khoa Kết nối tri thức với cuộc sống (Chương V), trong đường tròn (O; 17 cm), độ dài của dây cung lớn nhất (đường kính) bằng:\par
\hspace*{1.5em}\textbf{A.} 51 cm\quad\quad \textbf{B.} 17 cm\quad\quad \textbf{C.} 34 cm\quad\quad \textbf{D.} 289$\pi$ cm\par\vspace{2pt}
\textbf{\color{primary} Câu 14.} Tìm đáp số đúng của bài toán cơ bản sau: Trong một đường tròn, góc ở tâm chắn cung AB có số đo bằng 128°. Góc nội tiếp cùng chắn cung AB đó có số đo bằng:\par
\hspace*{1.5em}\textbf{A.} 64°\quad\quad \textbf{B.} 256°\quad\quad \textbf{C.} 90°\quad\quad \textbf{D.} 128°\par\vspace{2pt}
\textbf{\color{primary} Câu 15.} Xét đối tượng toán học trong Chương V: Cho đường tròn (O; 18 cm) và đường thẳng d có khoảng cách từ tâm O đến d là OH = 18 cm. Vị trí tương đối của d và (O) là. Giá trị hoặc tính chất tương ứng là:\par
\begin{itemize}\setlength{\itemsep}{1pt}\setlength{\parskip}{0pt}
  \item[] \textbf{A.} Đường thẳng d cắt đường tròn (O) tại hai điểm
  \item[] \textbf{B.} Đường thẳng d và (O) không giao nhau
  \item[] \textbf{C.} Đường thẳng d đi qua tâm O
  \item[] \textbf{D.} Đường thẳng d tiếp xúc với đường tròn (O)
\end{itemize}\vspace{2pt}
\textbf{\color{primary} Câu 16.} Hoàn thành phép xác định cơ bản sau đây: Trong đường tròn (O; 20 cm), độ dài của dây cung lớn nhất (đường kính) bằng:\par
\hspace*{1.5em}\textbf{A.} 60 cm\quad\quad \textbf{B.} 40 cm\quad\quad \textbf{C.} 20 cm\quad\quad \textbf{D.} 400$\pi$ cm\par\vspace{2pt}
\textbf{\color{primary} Câu 17.} Trong các khẳng định hoặc giá trị dưới đây, đâu là đáp án đúng cho câu hỏi: Trong một đường tròn, góc ở tâm chắn cung AB có số đo bằng 146°. Góc nội tiếp cùng chắn cung AB đó có số đo bằng?\par
\hspace*{1.5em}\textbf{A.} 292°\quad\quad \textbf{B.} 90°\quad\quad \textbf{C.} 146°\quad\quad \textbf{D.} 73°\par\vspace{2pt}
\textbf{\color{primary} Câu 18.} Bài tập khởi động Chương V: Cho đường tròn (O; 21 cm) và đường thẳng d có khoảng cách từ tâm O đến d là OH = 21 cm. Vị trí tương đối của d và (O) là:\par
\begin{itemize}\setlength{\itemsep}{1pt}\setlength{\parskip}{0pt}
  \item[] \textbf{A.} Đường thẳng d tiếp xúc với đường tròn (O)
  \item[] \textbf{B.} Đường thẳng d cắt đường tròn (O) tại hai điểm
  \item[] \textbf{C.} Đường thẳng d và (O) không giao nhau
  \item[] \textbf{D.} Đường thẳng d đi qua tâm O
\end{itemize}\vspace{2pt}
\textbf{\color{primary} Câu 19.} Quan sát dữ kiện đề bài: Trong đường tròn (O; 23 cm), độ dài của dây cung lớn nhất (đường kính) bằng. Hãy chọn phương án đúng:\par
\hspace*{1.5em}\textbf{A.} 46 cm\quad\quad \textbf{B.} 69 cm\quad\quad \textbf{C.} 23 cm\quad\quad \textbf{D.} 529$\pi$ cm\par\vspace{2pt}
\textbf{\color{primary} Câu 20.} Sử dụng công thức hoặc tính chất đầu bài của Chương V để giải quyết yêu cầu: Trong một đường tròn, góc ở tâm chắn cung AB có số đo bằng 164°. Góc nội tiếp cùng chắn cung AB đó có số đo bằng:\par
\hspace*{1.5em}\textbf{A.} 82°\quad\quad \textbf{B.} 164°\quad\quad \textbf{C.} 328°\quad\quad \textbf{D.} 90°\par\vspace{2pt}
\subsection*{\color{accent} Phần II. Mức độ Trung bình (Thông hiểu -- Câu 21 đến Câu 40)}
\textbf{\color{primary} Câu 21.} Trong đường tròn (O; R = 5 cm), cho dây cung AB = 8 cm. Khoảng cách từ tâm O đến dây AB bằng:\par
\hspace*{1.5em}\textbf{A.} 6 cm\quad\quad \textbf{B.} 2 cm\quad\quad \textbf{C.} 4 cm\quad\quad \textbf{D.} 3 cm\par\vspace{2pt}
\textbf{\color{primary} Câu 22.} Vận dụng quy tắc trong Chương V (Đường tròn) để giải bài toán: Độ dài cung 60° của một đường tròn có bán kính R = 9 cm bằng:\par
\hspace*{1.5em}\textbf{A.} 3$\pi$ cm\quad\quad \textbf{B.} 18$\pi$ cm\quad\quad \textbf{C.} 9$\pi$ cm\quad\quad \textbf{D.} 6$\pi$ cm\par\vspace{2pt}
\textbf{\color{primary} Câu 23.} Biến đổi và tính toán theo kiến thức Chương V: Diện tích hình quạt tròn bán kính R = 6 cm ứng với cung 60° bằng. Kết quả thu được là:\par
\hspace*{1.5em}\textbf{A.} 3$\pi$ cm$^2$\quad\quad \textbf{B.} 36$\pi$ cm$^2$\quad\quad \textbf{C.} 6$\pi$ cm$^2$\quad\quad \textbf{D.} 12$\pi$ cm$^2$\par\vspace{2pt}
\textbf{\color{primary} Câu 24.} Xét bài toán mức độ thông hiểu sau: Trong đường tròn (O; R = 5 cm), cho dây cung AB = 8 cm. Khoảng cách từ tâm O đến dây AB bằng. Đáp án đúng là:\par
\hspace*{1.5em}\textbf{A.} 3 cm\quad\quad \textbf{B.} 6 cm\quad\quad \textbf{C.} 2 cm\quad\quad \textbf{D.} 4 cm\par\vspace{2pt}
\textbf{\color{primary} Câu 25.} Để hoàn thành phiếu học tập Chương V, hãy xác định: Độ dài cung 60° của một đường tròn có bán kính R = 18 cm bằng:\par
\hspace*{1.5em}\textbf{A.} 12$\pi$ cm\quad\quad \textbf{B.} 6$\pi$ cm\quad\quad \textbf{C.} 18$\pi$ cm\quad\quad \textbf{D.} 36$\pi$ cm\par\vspace{2pt}
\textbf{\color{primary} Câu 26.} Phân tích dữ kiện bài toán: Diện tích hình quạt tròn bán kính R = 6 cm ứng với cung 60° bằng. Giá trị cần tìm bằng:\par
\hspace*{1.5em}\textbf{A.} 12$\pi$ cm$^2$\quad\quad \textbf{B.} 6$\pi$ cm$^2$\quad\quad \textbf{C.} 3$\pi$ cm$^2$\quad\quad \textbf{D.} 36$\pi$ cm$^2$\par\vspace{2pt}
\textbf{\color{primary} Câu 27.} Áp dụng tính chất đã học ở Chương V vào bài tập sau: Trong đường tròn (O; R = 5 cm), cho dây cung AB = 8 cm. Khoảng cách từ tâm O đến dây AB bằng:\par
\hspace*{1.5em}\textbf{A.} 6 cm\quad\quad \textbf{B.} 3 cm\quad\quad \textbf{C.} 4 cm\quad\quad \textbf{D.} 2 cm\par\vspace{2pt}
\textbf{\color{primary} Câu 28.} Thực hiện các bước suy luận hợp lý để giải quyết yêu cầu: Độ dài cung 60° của một đường tròn có bán kính R = 27 cm bằng:\par
\hspace*{1.5em}\textbf{A.} 18$\pi$ cm\quad\quad \textbf{B.} 27$\pi$ cm\quad\quad \textbf{C.} 9$\pi$ cm\quad\quad \textbf{D.} 54$\pi$ cm\par\vspace{2pt}
\textbf{\color{primary} Câu 29.} Cho bài toán tự luyện (Đường tròn): Diện tích hình quạt tròn bán kính R = 6 cm ứng với cung 60° bằng. Phương án nào sau đây chính xác?\par
\hspace*{1.5em}\textbf{A.} 36$\pi$ cm$^2$\quad\quad \textbf{B.} 3$\pi$ cm$^2$\quad\quad \textbf{C.} 12$\pi$ cm$^2$\quad\quad \textbf{D.} 6$\pi$ cm$^2$\par\vspace{2pt}
\textbf{\color{primary} Câu 30.} Kết hợp các phép biến đổi trong Chương V, hãy tìm kết quả của bài toán: Trong đường tròn (O; R = 5 cm), cho dây cung AB = 8 cm. Khoảng cách từ tâm O đến dây AB bằng:\par
\hspace*{1.5em}\textbf{A.} 6 cm\quad\quad \textbf{B.} 3 cm\quad\quad \textbf{C.} 4 cm\quad\quad \textbf{D.} 2 cm\par\vspace{2pt}
\textbf{\color{primary} Câu 31.} Một bài kiểm tra định kỳ nội dung Chương V nêu câu hỏi: Độ dài cung 60° của một đường tròn có bán kính R = 36 cm bằng. Hãy chọn đáp số đúng:\par
\hspace*{1.5em}\textbf{A.} 12$\pi$ cm\quad\quad \textbf{B.} 24$\pi$ cm\quad\quad \textbf{C.} 36$\pi$ cm\quad\quad \textbf{D.} 72$\pi$ cm\par\vspace{2pt}
\textbf{\color{primary} Câu 32.} Tìm lời giải đúng cho bài toán thông hiểu: Diện tích hình quạt tròn bán kính R = 6 cm ứng với cung 60° bằng:\par
\hspace*{1.5em}\textbf{A.} 36$\pi$ cm$^2$\quad\quad \textbf{B.} 12$\pi$ cm$^2$\quad\quad \textbf{C.} 3$\pi$ cm$^2$\quad\quad \textbf{D.} 6$\pi$ cm$^2$\par\vspace{2pt}
\textbf{\color{primary} Câu 33.} Sử dụng mối liên hệ giữa các đại lượng trong Chương V: Trong đường tròn (O; R = 5 cm), cho dây cung AB = 8 cm. Khoảng cách từ tâm O đến dây AB bằng:\par
\hspace*{1.5em}\textbf{A.} 2 cm\quad\quad \textbf{B.} 6 cm\quad\quad \textbf{C.} 4 cm\quad\quad \textbf{D.} 3 cm\par\vspace{2pt}
\textbf{\color{primary} Câu 34.} Hãy tính toán và chọn phương án đúng đối với bài tập: Độ dài cung 60° của một đường tròn có bán kính R = 45 cm bằng:\par
\hspace*{1.5em}\textbf{A.} 90$\pi$ cm\quad\quad \textbf{B.} 30$\pi$ cm\quad\quad \textbf{C.} 45$\pi$ cm\quad\quad \textbf{D.} 15$\pi$ cm\par\vspace{2pt}
\textbf{\color{primary} Câu 35.} Dựa trên các tính chất trọng tâm của Chương V, xét bài toán: Diện tích hình quạt tròn bán kính R = 6 cm ứng với cung 60° bằng:\par
\hspace*{1.5em}\textbf{A.} 6$\pi$ cm$^2$\quad\quad \textbf{B.} 36$\pi$ cm$^2$\quad\quad \textbf{C.} 12$\pi$ cm$^2$\quad\quad \textbf{D.} 3$\pi$ cm$^2$\par\vspace{2pt}
\textbf{\color{primary} Câu 36.} Thực hành giải toán (Đường tròn): Trong đường tròn (O; R = 5 cm), cho dây cung AB = 8 cm. Khoảng cách từ tâm O đến dây AB bằng. Đáp án nào dưới đây thỏa mãn?\par
\hspace*{1.5em}\textbf{A.} 3 cm\quad\quad \textbf{B.} 4 cm\quad\quad \textbf{C.} 6 cm\quad\quad \textbf{D.} 2 cm\par\vspace{2pt}
\textbf{\color{primary} Câu 37.} Kiểm tra kỹ năng biến đổi và suy luận qua bài toán: Độ dài cung 60° của một đường tròn có bán kính R = 54 cm bằng:\par
\hspace*{1.5em}\textbf{A.} 108$\pi$ cm\quad\quad \textbf{B.} 54$\pi$ cm\quad\quad \textbf{C.} 18$\pi$ cm\quad\quad \textbf{D.} 36$\pi$ cm\par\vspace{2pt}
\textbf{\color{primary} Câu 38.} Cho giả thiết thuộc phạm vi Chương V: Diện tích hình quạt tròn bán kính R = 6 cm ứng với cung 60° bằng. Kết luận đúng là:\par
\hspace*{1.5em}\textbf{A.} 12$\pi$ cm$^2$\quad\quad \textbf{B.} 6$\pi$ cm$^2$\quad\quad \textbf{C.} 36$\pi$ cm$^2$\quad\quad \textbf{D.} 3$\pi$ cm$^2$\par\vspace{2pt}
\textbf{\color{primary} Câu 39.} Rút gọn hoặc tính giá trị theo yêu cầu sau của Chương V: Trong đường tròn (O; R = 5 cm), cho dây cung AB = 8 cm. Khoảng cách từ tâm O đến dây AB bằng:\par
\hspace*{1.5em}\textbf{A.} 2 cm\quad\quad \textbf{B.} 6 cm\quad\quad \textbf{C.} 3 cm\quad\quad \textbf{D.} 4 cm\par\vspace{2pt}
\textbf{\color{primary} Câu 40.} Xác định phương án đúng sau khi thực hiện phép tính / suy luận: Độ dài cung 60° của một đường tròn có bán kính R = 63 cm bằng:\par
\hspace*{1.5em}\textbf{A.} 63$\pi$ cm\quad\quad \textbf{B.} 21$\pi$ cm\quad\quad \textbf{C.} 126$\pi$ cm\quad\quad \textbf{D.} 42$\pi$ cm\par\vspace{2pt}
\subsection*{\color{accent} Phần III. Mức độ Khó (Vận dụng \& Vận dụng cao -- Câu 41 đến Câu 50)}
\textbf{\color{primary} Câu 41.} Từ điểm M nằm ngoài đường tròn (O; R = 3 cm) cách tâm O một khoảng OM = 5 cm, kẻ hai tiếp tuyến MA, MB tới đường tròn (A, B là hai tiếp điểm). Độ dài đoạn tiếp tuyến MA bằng:\par
\hspace*{1.5em}\textbf{A.} 6 cm\quad\quad \textbf{B.} 2 cm\quad\quad \textbf{C.} 8 cm\quad\quad \textbf{D.} 4 cm\par\vspace{2pt}
\textbf{\color{primary} Câu 42.} Bài toán vận dụng (Chương V – Đường tròn): Diện tích hình vành khuyên giới hạn bởi hai đường tròn đồng tâm có bán kính R$_1$ = 7 cm và R$_2$ = 5 cm bằng. Kết quả đúng là:\par
\hspace*{1.5em}\textbf{A.} 4$\pi$ cm$^2$\quad\quad \textbf{B.} 49$\pi$ cm$^2$\quad\quad \textbf{C.} 12$\pi$ cm$^2$\quad\quad \textbf{D.} 24$\pi$ cm$^2$\par\vspace{2pt}
\textbf{\color{primary} Câu 43.} Vận dụng kiến thức nâng cao trong phạm vi Chương V để giải bài toán: Từ điểm M nằm ngoài đường tròn (O; R = 3 cm) cách tâm O một khoảng OM = 5 cm, kẻ hai tiếp tuyến MA, MB tới đường tròn (A, B là hai tiếp điểm). Độ dài đoạn tiếp tuyến MA bằng:\par
\hspace*{1.5em}\textbf{A.} 4 cm\quad\quad \textbf{B.} 8 cm\quad\quad \textbf{C.} 2 cm\quad\quad \textbf{D.} 6 cm\par\vspace{2pt}
\textbf{\color{primary} Câu 44.} Phân tích và giải bài toán tổng hợp sau của Chương V: Diện tích hình vành khuyên giới hạn bởi hai đường tròn đồng tâm có bán kính R$_1$ = 9 cm và R$_2$ = 7 cm bằng:\par
\hspace*{1.5em}\textbf{A.} 4$\pi$ cm$^2$\quad\quad \textbf{B.} 32$\pi$ cm$^2$\quad\quad \textbf{C.} 81$\pi$ cm$^2$\quad\quad \textbf{D.} 16$\pi$ cm$^2$\par\vspace{2pt}
\textbf{\color{primary} Câu 45.} Thử thách tư duy trong phạm vi kiến thức Chương V: Từ điểm M nằm ngoài đường tròn (O; R = 3 cm) cách tâm O một khoảng OM = 5 cm, kẻ hai tiếp tuyến MA, MB tới đường tròn (A, B là hai tiếp điểm). Độ dài đoạn tiếp tuyến MA bằng. Đáp số chính xác là:\par
\hspace*{1.5em}\textbf{A.} 6 cm\quad\quad \textbf{B.} 8 cm\quad\quad \textbf{C.} 2 cm\quad\quad \textbf{D.} 4 cm\par\vspace{2pt}
\textbf{\color{primary} Câu 46.} Sử dụng phương pháp giải toán của Chương V để tìm đáp án cho bài toán: Diện tích hình vành khuyên giới hạn bởi hai đường tròn đồng tâm có bán kính R$_1$ = 11 cm và R$_2$ = 9 cm bằng:\par
\hspace*{1.5em}\textbf{A.} 4$\pi$ cm$^2$\quad\quad \textbf{B.} 121$\pi$ cm$^2$\quad\quad \textbf{C.} 40$\pi$ cm$^2$\quad\quad \textbf{D.} 20$\pi$ cm$^2$\par\vspace{2pt}
\textbf{\color{primary} Câu 47.} Xét bài toán vận dụng thực tiễn / nâng cao thuộc Chương V: Từ điểm M nằm ngoài đường tròn (O; R = 3 cm) cách tâm O một khoảng OM = 5 cm, kẻ hai tiếp tuyến MA, MB tới đường tròn (A, B là hai tiếp điểm). Độ dài đoạn tiếp tuyến MA bằng:\par
\hspace*{1.5em}\textbf{A.} 4 cm\quad\quad \textbf{B.} 6 cm\quad\quad \textbf{C.} 2 cm\quad\quad \textbf{D.} 8 cm\par\vspace{2pt}
\textbf{\color{primary} Câu 48.} Tìm kết quả cuối cùng của bài toán phân loại học sinh khá giỏi (Chương V): Diện tích hình vành khuyên giới hạn bởi hai đường tròn đồng tâm có bán kính R$_1$ = 13 cm và R$_2$ = 11 cm bằng:\par
\hspace*{1.5em}\textbf{A.} 24$\pi$ cm$^2$\quad\quad \textbf{B.} 169$\pi$ cm$^2$\quad\quad \textbf{C.} 4$\pi$ cm$^2$\quad\quad \textbf{D.} 48$\pi$ cm$^2$\par\vspace{2pt}
\textbf{\color{primary} Câu 49.} Kết hợp các tính chất đã học trong Chương V, hãy giải quyết bài toán: Từ điểm M nằm ngoài đường tròn (O; R = 3 cm) cách tâm O một khoảng OM = 5 cm, kẻ hai tiếp tuyến MA, MB tới đường tròn (A, B là hai tiếp điểm). Độ dài đoạn tiếp tuyến MA bằng:\par
\hspace*{1.5em}\textbf{A.} 6 cm\quad\quad \textbf{B.} 8 cm\quad\quad \textbf{C.} 2 cm\quad\quad \textbf{D.} 4 cm\par\vspace{2pt}
\textbf{\color{primary} Câu 50.} Câu hỏi vận dụng cao cuối Chương V: Diện tích hình vành khuyên giới hạn bởi hai đường tròn đồng tâm có bán kính R$_1$ = 15 cm và R$_2$ = 13 cm bằng. Phương án đúng là:\par
\hspace*{1.5em}\textbf{A.} 28$\pi$ cm$^2$\quad\quad \textbf{B.} 4$\pi$ cm$^2$\quad\quad \textbf{C.} 225$\pi$ cm$^2$\quad\quad \textbf{D.} 56$\pi$ cm$^2$\par\vspace{2pt}
\subsection*{\color{primary} Bảng đáp án Chương V: Đường tròn}
\begin{center}
\renewcommand{\arraystretch}{1.25}
\begin{tabularx}{\textwidth}{*{10}{>{\centering\arraybackslash}X}}
\toprule
\textbf{\color{primary} Cột 1} & \textbf{\color{primary} Cột 2} & \textbf{\color{primary} Cột 3} & \textbf{\color{primary} Cột 4} & \textbf{\color{primary} Cột 5} & \textbf{\color{primary} Cột 6} & \textbf{\color{primary} Cột 7} & \textbf{\color{primary} Cột 8} & \textbf{\color{primary} Cột 9} & \textbf{\color{primary} Cột 10} \\
\midrule
\textbf{1.B} & \textbf{2.C} & \textbf{3.B} & \textbf{4.C} & \textbf{5.B} & \textbf{6.A} & \textbf{7.D} & \textbf{8.B} & \textbf{9.A} & \textbf{10.D} \\
\rowcolor{rowtint} \textbf{11.A} & \textbf{12.C} & \textbf{13.C} & \textbf{14.A} & \textbf{15.D} & \textbf{16.B} & \textbf{17.D} & \textbf{18.A} & \textbf{19.A} & \textbf{20.A} \\
\textbf{21.D} & \textbf{22.A} & \textbf{23.C} & \textbf{24.A} & \textbf{25.B} & \textbf{26.B} & \textbf{27.B} & \textbf{28.C} & \textbf{29.D} & \textbf{30.B} \\
\rowcolor{rowtint} \textbf{31.A} & \textbf{32.D} & \textbf{33.D} & \textbf{34.D} & \textbf{35.A} & \textbf{36.A} & \textbf{37.C} & \textbf{38.B} & \textbf{39.C} & \textbf{40.B} \\
\textbf{41.D} & \textbf{42.D} & \textbf{43.A} & \textbf{44.B} & \textbf{45.D} & \textbf{46.C} & \textbf{47.A} & \textbf{48.D} & \textbf{49.D} & \textbf{50.D} \\
\bottomrule
\end{tabularx}
\end{center}\vspace{10pt}
\section*{\color{primary} CHƯƠNG VI: HÀM SỐ Y = AX$^2$ (A $\ne$ 0). PHƯƠNG TRÌNH BẬC HAI MỘT ẨN}
\textit{Trọng tâm kiến thức SGK: Hàm số y = ax$^2$ (a $\ne$ 0), phương trình bậc hai một ẩn, công thức nghiệm, định lí Viète và ứng dụng, giải bài toán bằng cách lập phương trình bậc hai.}\par\vspace{4pt}
\subsection*{\color{accent} Phần I. Mức độ Dễ (Nhận biết -- Câu 1 đến Câu 20)}
\textbf{\color{primary} Câu 1.} Cho hàm số y = 2x$^2$. Giá trị của hàm số tại x = -2 bằng:\par
\hspace*{1.5em}\textbf{A.} 8\quad\quad \textbf{B.} -4\quad\quad \textbf{C.} 4\quad\quad \textbf{D.} -8\par\vspace{2pt}
\textbf{\color{primary} Câu 2.} Thực hiện yêu cầu nhận biết sau: Biệt thức $\Delta$ của phương trình bậc hai x$^2$ - 6x + 3 = 0 bằng.\par
\hspace*{1.5em}\textbf{A.} 24\quad\quad \textbf{B.} 48\quad\quad \textbf{C.} -6\quad\quad \textbf{D.} 6\par\vspace{2pt}
\textbf{\color{primary} Câu 3.} Trong bài học thuộc Chương VI, xét nội dung: Gọi x$_1$, x$_2$ là hai nghiệm của phương trình x$^2$ - 8x + 6 = 0. Theo định lí Viète, tổng x$_1$ + x$_2$ và tích x$_1$$\cdot$x$_2$ lần lượt là. Đáp án đúng là:\par
\hspace*{1.5em}\textbf{A.} x$_1$ + x$_2$ = 8; x$_1$$\cdot$x$_2$ = 6\quad\quad \textbf{B.} x$_1$ + x$_2$ = 8; x$_1$$\cdot$x$_2$ = -6\quad\quad \textbf{C.} x$_1$ + x$_2$ = 6; x$_1$$\cdot$x$_2$ = 8\quad\quad \textbf{D.} x$_1$ + x$_2$ = -8; x$_1$$\cdot$x$_2$ = 6\par\vspace{2pt}
\textbf{\color{primary} Câu 4.} Dựa vào kiến thức chuẩn SGK của Chương VI (Hàm số y = ax$^2$ (a $\ne$ 0). Phương trình bậc hai một ẩn), cho hàm số y = 5x$^2$. Giá trị của hàm số tại x = -2 bằng:\par
\hspace*{1.5em}\textbf{A.} 20\quad\quad \textbf{B.} -20\quad\quad \textbf{C.} 10\quad\quad \textbf{D.} -10\par\vspace{2pt}
\textbf{\color{primary} Câu 5.} Chọn kết quả đúng nhất khi xét: Biệt thức $\Delta$ của phương trình bậc hai x$^2$ - 12x + 3 = 0 bằng:\par
\hspace*{1.5em}\textbf{A.} 33\quad\quad \textbf{B.} 132\quad\quad \textbf{C.} 156\quad\quad \textbf{D.} 0\par\vspace{2pt}
\textbf{\color{primary} Câu 6.} Khi làm bài tập cơ bản về hàm số y = ax$^2$ (a $\ne$ 0). phương trình bậc hai một ẩn, một học sinh gặp câu hỏi: Gọi x$_1$, x$_2$ là hai nghiệm của phương trình x$^2$ - 11x + 9 = 0. Theo định lí Viète, tổng x$_1$ + x$_2$ và tích x$_1$$\cdot$x$_2$ lần lượt là. Kết quả nào dưới đây đúng?\par
\hspace*{1.5em}\textbf{A.} x$_1$ + x$_2$ = -11; x$_1$$\cdot$x$_2$ = 9\quad\quad \textbf{B.} x$_1$ + x$_2$ = 9; x$_1$$\cdot$x$_2$ = 11\quad\quad \textbf{C.} x$_1$ + x$_2$ = 11; x$_1$$\cdot$x$_2$ = 9\quad\quad \textbf{D.} x$_1$ + x$_2$ = 11; x$_1$$\cdot$x$_2$ = -9\par\vspace{2pt}
\textbf{\color{primary} Câu 7.} Xác định giá trị hoặc khẳng định đúng cho bài toán sau: Cho hàm số y = 8x$^2$. Giá trị của hàm số tại x = -2 bằng:\par
\hspace*{1.5em}\textbf{A.} 32\quad\quad \textbf{B.} 16\quad\quad \textbf{C.} -32\quad\quad \textbf{D.} -16\par\vspace{2pt}
\textbf{\color{primary} Câu 8.} Áp dụng trực tiếp định nghĩa và quy tắc trong Chương VI: Biệt thức $\Delta$ của phương trình bậc hai x$^2$ - 18x + 3 = 0 bằng:\par
\hspace*{1.5em}\textbf{A.} 78\quad\quad \textbf{B.} 6\quad\quad \textbf{C.} 336\quad\quad \textbf{D.} 312\par\vspace{2pt}
\textbf{\color{primary} Câu 9.} Kiểm tra nhanh kiến thức nền tảng: Gọi x$_1$, x$_2$ là hai nghiệm của phương trình x$^2$ - 14x + 12 = 0. Theo định lí Viète, tổng x$_1$ + x$_2$ và tích x$_1$$\cdot$x$_2$ lần lượt là:\par
\begin{itemize}\setlength{\itemsep}{1pt}\setlength{\parskip}{0pt}
  \item[] \textbf{A.} x$_1$ + x$_2$ = 14; x$_1$$\cdot$x$_2$ = 12
  \item[] \textbf{B.} x$_1$ + x$_2$ = -14; x$_1$$\cdot$x$_2$ = 12
  \item[] \textbf{C.} x$_1$ + x$_2$ = 12; x$_1$$\cdot$x$_2$ = 14
  \item[] \textbf{D.} x$_1$ + x$_2$ = 14; x$_1$$\cdot$x$_2$ = -12
\end{itemize}\vspace{2pt}
\textbf{\color{primary} Câu 10.} Cho dữ kiện nhận biết sau đây: Cho hàm số y = 11x$^2$. Giá trị của hàm số tại x = -2 bằng. Phương án chính xác là:\par
\hspace*{1.5em}\textbf{A.} 22\quad\quad \textbf{B.} -44\quad\quad \textbf{C.} -22\quad\quad \textbf{D.} 44\par\vspace{2pt}
\textbf{\color{primary} Câu 11.} Hãy chỉ ra kết quả đúng trong bốn lựa chọn A, B, C, D đối với yêu cầu: Biệt thức $\Delta$ của phương trình bậc hai x$^2$ - 24x + 3 = 0 bằng:\par
\hspace*{1.5em}\textbf{A.} 588\quad\quad \textbf{B.} 12\quad\quad \textbf{C.} 564\quad\quad \textbf{D.} 141\par\vspace{2pt}
\textbf{\color{primary} Câu 12.} Câu hỏi ôn tập lý thuyết và kỹ năng cơ bản (Chương VI): Gọi x$_1$, x$_2$ là hai nghiệm của phương trình x$^2$ - 17x + 15 = 0. Theo định lí Viète, tổng x$_1$ + x$_2$ và tích x$_1$$\cdot$x$_2$ lần lượt là:\par
\begin{itemize}\setlength{\itemsep}{1pt}\setlength{\parskip}{0pt}
  \item[] \textbf{A.} x$_1$ + x$_2$ = -17; x$_1$$\cdot$x$_2$ = 15
  \item[] \textbf{B.} x$_1$ + x$_2$ = 17; x$_1$$\cdot$x$_2$ = 15
  \item[] \textbf{C.} x$_1$ + x$_2$ = 17; x$_1$$\cdot$x$_2$ = -15
  \item[] \textbf{D.} x$_1$ + x$_2$ = 15; x$_1$$\cdot$x$_2$ = 17
\end{itemize}\vspace{2pt}
\textbf{\color{primary} Câu 13.} Theo nội dung sách giáo khoa Kết nối tri thức với cuộc sống (Chương VI), cho hàm số y = 14x$^2$. Giá trị của hàm số tại x = -2 bằng:\par
\hspace*{1.5em}\textbf{A.} -56\quad\quad \textbf{B.} 28\quad\quad \textbf{C.} -28\quad\quad \textbf{D.} 56\par\vspace{2pt}
\textbf{\color{primary} Câu 14.} Tìm đáp số đúng của bài toán cơ bản sau: Biệt thức $\Delta$ của phương trình bậc hai x$^2$ - 30x + 3 = 0 bằng:\par
\hspace*{1.5em}\textbf{A.} 18\quad\quad \textbf{B.} 912\quad\quad \textbf{C.} 888\quad\quad \textbf{D.} 222\par\vspace{2pt}
\textbf{\color{primary} Câu 15.} Xét đối tượng toán học trong Chương VI: Gọi x$_1$, x$_2$ là hai nghiệm của phương trình x$^2$ - 20x + 18 = 0. Theo định lí Viète, tổng x$_1$ + x$_2$ và tích x$_1$$\cdot$x$_2$ lần lượt là. Giá trị hoặc tính chất tương ứng là:\par
\begin{itemize}\setlength{\itemsep}{1pt}\setlength{\parskip}{0pt}
  \item[] \textbf{A.} x$_1$ + x$_2$ = 20; x$_1$$\cdot$x$_2$ = -18
  \item[] \textbf{B.} x$_1$ + x$_2$ = 18; x$_1$$\cdot$x$_2$ = 20
  \item[] \textbf{C.} x$_1$ + x$_2$ = 20; x$_1$$\cdot$x$_2$ = 18
  \item[] \textbf{D.} x$_1$ + x$_2$ = -20; x$_1$$\cdot$x$_2$ = 18
\end{itemize}\vspace{2pt}
\textbf{\color{primary} Câu 16.} Hoàn thành phép xác định cơ bản sau đây: Cho hàm số y = 17x$^2$. Giá trị của hàm số tại x = -2 bằng:\par
\hspace*{1.5em}\textbf{A.} -34\quad\quad \textbf{B.} 68\quad\quad \textbf{C.} -68\quad\quad \textbf{D.} 34\par\vspace{2pt}
\textbf{\color{primary} Câu 17.} Trong các khẳng định hoặc giá trị dưới đây, đâu là đáp án đúng cho câu hỏi: Biệt thức $\Delta$ của phương trình bậc hai x$^2$ - 36x + 3 = 0 bằng?\par
\hspace*{1.5em}\textbf{A.} 1308\quad\quad \textbf{B.} 1284\quad\quad \textbf{C.} 321\quad\quad \textbf{D.} 24\par\vspace{2pt}
\textbf{\color{primary} Câu 18.} Bài tập khởi động Chương VI: Gọi x$_1$, x$_2$ là hai nghiệm của phương trình x$^2$ - 23x + 21 = 0. Theo định lí Viète, tổng x$_1$ + x$_2$ và tích x$_1$$\cdot$x$_2$ lần lượt là:\par
\begin{itemize}\setlength{\itemsep}{1pt}\setlength{\parskip}{0pt}
  \item[] \textbf{A.} x$_1$ + x$_2$ = 23; x$_1$$\cdot$x$_2$ = 21
  \item[] \textbf{B.} x$_1$ + x$_2$ = 23; x$_1$$\cdot$x$_2$ = -21
  \item[] \textbf{C.} x$_1$ + x$_2$ = -23; x$_1$$\cdot$x$_2$ = 21
  \item[] \textbf{D.} x$_1$ + x$_2$ = 21; x$_1$$\cdot$x$_2$ = 23
\end{itemize}\vspace{2pt}
\textbf{\color{primary} Câu 19.} Quan sát dữ kiện đề bài: Cho hàm số y = 20x$^2$. Giá trị của hàm số tại x = -2 bằng. Hãy chọn phương án đúng:\par
\hspace*{1.5em}\textbf{A.} -80\quad\quad \textbf{B.} 80\quad\quad \textbf{C.} 40\quad\quad \textbf{D.} -40\par\vspace{2pt}
\textbf{\color{primary} Câu 20.} Sử dụng công thức hoặc tính chất đầu bài của Chương VI để giải quyết yêu cầu: Biệt thức $\Delta$ của phương trình bậc hai x$^2$ - 42x + 3 = 0 bằng:\par
\hspace*{1.5em}\textbf{A.} 1776\quad\quad \textbf{B.} 1752\quad\quad \textbf{C.} 30\quad\quad \textbf{D.} 438\par\vspace{2pt}
\subsection*{\color{accent} Phần II. Mức độ Trung bình (Thông hiểu -- Câu 21 đến Câu 40)}
\textbf{\color{primary} Câu 21.} Đồ thị hàm số y = ax$^2$ (a $\ne$ 0) đi qua điểm M(2; 8) khi hệ số a bằng:\par
\hspace*{1.5em}\textbf{A.} a = 4\quad\quad \textbf{B.} a = 2\quad\quad \textbf{C.} a = 8\quad\quad \textbf{D.} a = -2\par\vspace{2pt}
\textbf{\color{primary} Câu 22.} Vận dụng quy tắc trong Chương VI (Hàm số y = ax$^2$ (a $\ne$ 0). Phương trình bậc hai một ẩn) để giải bài toán: Phương trình bậc hai x$^2$ - 9x + 18 = 0 có hai nghiệm phân biệt là:\par
\hspace*{1.5em}\textbf{A.} x$_1$ = 3; x$_2$ = -6\quad\quad \textbf{B.} x$_1$ = 3; x$_2$ = 6\quad\quad \textbf{C.} x$_1$ = 1; x$_2$ = 18\quad\quad \textbf{D.} x$_1$ = -3; x$_2$ = -6\par\vspace{2pt}
\textbf{\color{primary} Câu 23.} Biến đổi và tính toán theo kiến thức Chương VI: Tìm giá trị của tham số m để phương trình x$^2$ - 10x + m = 0 có nghiệm kép. Kết quả thu được là:\par
\hspace*{1.5em}\textbf{A.} m = -25\quad\quad \textbf{B.} m = 25\quad\quad \textbf{C.} m = 10\quad\quad \textbf{D.} m = 100\par\vspace{2pt}
\textbf{\color{primary} Câu 24.} Xét bài toán mức độ thông hiểu sau: Đồ thị hàm số y = ax$^2$ (a $\ne$ 0) đi qua điểm M(2; 20) khi hệ số a bằng. Đáp án đúng là:\par
\hspace*{1.5em}\textbf{A.} a = -5\quad\quad \textbf{B.} a = 5\quad\quad \textbf{C.} a = 20\quad\quad \textbf{D.} a = 10\par\vspace{2pt}
\textbf{\color{primary} Câu 25.} Để hoàn thành phiếu học tập Chương VI, hãy xác định: Phương trình bậc hai x$^2$ - 15x + 54 = 0 có hai nghiệm phân biệt là:\par
\hspace*{1.5em}\textbf{A.} x$_1$ = 6; x$_2$ = 9\quad\quad \textbf{B.} x$_1$ = 1; x$_2$ = 54\quad\quad \textbf{C.} x$_1$ = 6; x$_2$ = -9\quad\quad \textbf{D.} x$_1$ = -6; x$_2$ = -9\par\vspace{2pt}
\textbf{\color{primary} Câu 26.} Phân tích dữ kiện bài toán: Tìm giá trị của tham số m để phương trình x$^2$ - 16x + m = 0 có nghiệm kép. Giá trị cần tìm bằng:\par
\hspace*{1.5em}\textbf{A.} m = 64\quad\quad \textbf{B.} m = -64\quad\quad \textbf{C.} m = 16\quad\quad \textbf{D.} m = 256\par\vspace{2pt}
\textbf{\color{primary} Câu 27.} Áp dụng tính chất đã học ở Chương VI vào bài tập sau: Đồ thị hàm số y = ax$^2$ (a $\ne$ 0) đi qua điểm M(2; 32) khi hệ số a bằng:\par
\hspace*{1.5em}\textbf{A.} a = -8\quad\quad \textbf{B.} a = 32\quad\quad \textbf{C.} a = 8\quad\quad \textbf{D.} a = 16\par\vspace{2pt}
\textbf{\color{primary} Câu 28.} Thực hiện các bước suy luận hợp lý để giải quyết yêu cầu: Phương trình bậc hai x$^2$ - 21x + 108 = 0 có hai nghiệm phân biệt là:\par
\hspace*{1.5em}\textbf{A.} x$_1$ = 9; x$_2$ = -12\quad\quad \textbf{B.} x$_1$ = -9; x$_2$ = -12\quad\quad \textbf{C.} x$_1$ = 1; x$_2$ = 108\quad\quad \textbf{D.} x$_1$ = 9; x$_2$ = 12\par\vspace{2pt}
\textbf{\color{primary} Câu 29.} Cho bài toán tự luyện (Hàm số y = ax$^2$ (a $\ne$ 0). Phương trình bậc hai một ẩn): Tìm giá trị của tham số m để phương trình x$^2$ - 22x + m = 0 có nghiệm kép. Phương án nào sau đây chính xác?\par
\hspace*{1.5em}\textbf{A.} m = -121\quad\quad \textbf{B.} m = 121\quad\quad \textbf{C.} m = 484\quad\quad \textbf{D.} m = 22\par\vspace{2pt}
\textbf{\color{primary} Câu 30.} Kết hợp các phép biến đổi trong Chương VI, hãy tìm kết quả của bài toán: Đồ thị hàm số y = ax$^2$ (a $\ne$ 0) đi qua điểm M(2; 44) khi hệ số a bằng:\par
\hspace*{1.5em}\textbf{A.} a = 11\quad\quad \textbf{B.} a = 22\quad\quad \textbf{C.} a = 44\quad\quad \textbf{D.} a = -11\par\vspace{2pt}
\textbf{\color{primary} Câu 31.} Một bài kiểm tra định kỳ nội dung Chương VI nêu câu hỏi: Phương trình bậc hai x$^2$ - 27x + 180 = 0 có hai nghiệm phân biệt là. Hãy chọn đáp số đúng:\par
\hspace*{1.5em}\textbf{A.} x$_1$ = -12; x$_2$ = -15\quad\quad \textbf{B.} x$_1$ = 12; x$_2$ = 15\quad\quad \textbf{C.} x$_1$ = 1; x$_2$ = 180\quad\quad \textbf{D.} x$_1$ = 12; x$_2$ = -15\par\vspace{2pt}
\textbf{\color{primary} Câu 32.} Tìm lời giải đúng cho bài toán thông hiểu: Tìm giá trị của tham số m để phương trình x$^2$ - 28x + m = 0 có nghiệm kép:\par
\hspace*{1.5em}\textbf{A.} m = 196\quad\quad \textbf{B.} m = 784\quad\quad \textbf{C.} m = 28\quad\quad \textbf{D.} m = -196\par\vspace{2pt}
\textbf{\color{primary} Câu 33.} Sử dụng mối liên hệ giữa các đại lượng trong Chương VI: Đồ thị hàm số y = ax$^2$ (a $\ne$ 0) đi qua điểm M(2; 56) khi hệ số a bằng:\par
\hspace*{1.5em}\textbf{A.} a = 14\quad\quad \textbf{B.} a = 28\quad\quad \textbf{C.} a = 56\quad\quad \textbf{D.} a = -14\par\vspace{2pt}
\textbf{\color{primary} Câu 34.} Hãy tính toán và chọn phương án đúng đối với bài tập: Phương trình bậc hai x$^2$ - 33x + 270 = 0 có hai nghiệm phân biệt là:\par
\hspace*{1.5em}\textbf{A.} x$_1$ = 15; x$_2$ = 18\quad\quad \textbf{B.} x$_1$ = -15; x$_2$ = -18\quad\quad \textbf{C.} x$_1$ = 1; x$_2$ = 270\quad\quad \textbf{D.} x$_1$ = 15; x$_2$ = -18\par\vspace{2pt}
\textbf{\color{primary} Câu 35.} Dựa trên các tính chất trọng tâm của Chương VI, xét bài toán: Tìm giá trị của tham số m để phương trình x$^2$ - 34x + m = 0 có nghiệm kép:\par
\hspace*{1.5em}\textbf{A.} m = 1156\quad\quad \textbf{B.} m = -289\quad\quad \textbf{C.} m = 289\quad\quad \textbf{D.} m = 34\par\vspace{2pt}
\textbf{\color{primary} Câu 36.} Thực hành giải toán (Hàm số y = ax$^2$ (a $\ne$ 0). Phương trình bậc hai một ẩn): Đồ thị hàm số y = ax$^2$ (a $\ne$ 0) đi qua điểm M(2; 68) khi hệ số a bằng. Đáp án nào dưới đây thỏa mãn?\par
\hspace*{1.5em}\textbf{A.} a = 34\quad\quad \textbf{B.} a = 68\quad\quad \textbf{C.} a = 17\quad\quad \textbf{D.} a = -17\par\vspace{2pt}
\textbf{\color{primary} Câu 37.} Kiểm tra kỹ năng biến đổi và suy luận qua bài toán: Phương trình bậc hai x$^2$ - 39x + 378 = 0 có hai nghiệm phân biệt là:\par
\hspace*{1.5em}\textbf{A.} x$_1$ = -18; x$_2$ = -21\quad\quad \textbf{B.} x$_1$ = 1; x$_2$ = 378\quad\quad \textbf{C.} x$_1$ = 18; x$_2$ = -21\quad\quad \textbf{D.} x$_1$ = 18; x$_2$ = 21\par\vspace{2pt}
\textbf{\color{primary} Câu 38.} Cho giả thiết thuộc phạm vi Chương VI: Tìm giá trị của tham số m để phương trình x$^2$ - 40x + m = 0 có nghiệm kép. Kết luận đúng là:\par
\hspace*{1.5em}\textbf{A.} m = 400\quad\quad \textbf{B.} m = 40\quad\quad \textbf{C.} m = 1600\quad\quad \textbf{D.} m = -400\par\vspace{2pt}
\textbf{\color{primary} Câu 39.} Rút gọn hoặc tính giá trị theo yêu cầu sau của Chương VI: Đồ thị hàm số y = ax$^2$ (a $\ne$ 0) đi qua điểm M(2; 80) khi hệ số a bằng:\par
\hspace*{1.5em}\textbf{A.} a = 80\quad\quad \textbf{B.} a = -20\quad\quad \textbf{C.} a = 20\quad\quad \textbf{D.} a = 40\par\vspace{2pt}
\textbf{\color{primary} Câu 40.} Xác định phương án đúng sau khi thực hiện phép tính / suy luận: Phương trình bậc hai x$^2$ - 45x + 504 = 0 có hai nghiệm phân biệt là:\par
\hspace*{1.5em}\textbf{A.} x$_1$ = 1; x$_2$ = 504\quad\quad \textbf{B.} x$_1$ = 21; x$_2$ = 24\quad\quad \textbf{C.} x$_1$ = 21; x$_2$ = -24\quad\quad \textbf{D.} x$_1$ = -21; x$_2$ = -24\par\vspace{2pt}
\subsection*{\color{accent} Phần III. Mức độ Khó (Vận dụng \& Vận dụng cao -- Câu 41 đến Câu 50)}
\textbf{\color{primary} Câu 41.} Gọi x$_1$, x$_2$ là hai nghiệm của phương trình x$^2$ - 5x + 2 = 0. Giá trị của biểu thức A = x$_1$$^2$ + x$_2$$^2$ bằng:\par
\hspace*{1.5em}\textbf{A.} 29\quad\quad \textbf{B.} 21\quad\quad \textbf{C.} 23\quad\quad \textbf{D.} 25\par\vspace{2pt}
\textbf{\color{primary} Câu 42.} Bài toán vận dụng (Chương VI – Hàm số y = ax$^2$ (a $\ne$ 0). Phương trình bậc hai một ẩn): Tìm m để phương trình x$^2$ - 6x + m = 0 có hai nghiệm phân biệt x$_1$, x$_2$ thỏa mãn x$_1$$^2$ + x$_2$$^2$ = 20. Kết quả đúng là:\par
\hspace*{1.5em}\textbf{A.} m = 10\quad\quad \textbf{B.} m = 6\quad\quad \textbf{C.} m = -8\quad\quad \textbf{D.} m = 8\par\vspace{2pt}
\textbf{\color{primary} Câu 43.} Vận dụng kiến thức nâng cao trong phạm vi Chương VI để giải bài toán: Gọi x$_1$, x$_2$ là hai nghiệm của phương trình x$^2$ - 7x + 2 = 0. Giá trị của biểu thức A = x$_1$$^2$ + x$_2$$^2$ bằng:\par
\hspace*{1.5em}\textbf{A.} 49\quad\quad \textbf{B.} 53\quad\quad \textbf{C.} 47\quad\quad \textbf{D.} 45\par\vspace{2pt}
\textbf{\color{primary} Câu 44.} Phân tích và giải bài toán tổng hợp sau của Chương VI: Tìm m để phương trình x$^2$ - 6x + m = 0 có hai nghiệm phân biệt x$_1$, x$_2$ thỏa mãn x$_1$$^2$ + x$_2$$^2$ = 20:\par
\hspace*{1.5em}\textbf{A.} m = -8\quad\quad \textbf{B.} m = 6\quad\quad \textbf{C.} m = 10\quad\quad \textbf{D.} m = 8\par\vspace{2pt}
\textbf{\color{primary} Câu 45.} Thử thách tư duy trong phạm vi kiến thức Chương VI: Gọi x$_1$, x$_2$ là hai nghiệm của phương trình x$^2$ - 9x + 2 = 0. Giá trị của biểu thức A = x$_1$$^2$ + x$_2$$^2$ bằng. Đáp số chính xác là:\par
\hspace*{1.5em}\textbf{A.} 79\quad\quad \textbf{B.} 85\quad\quad \textbf{C.} 81\quad\quad \textbf{D.} 77\par\vspace{2pt}
\textbf{\color{primary} Câu 46.} Sử dụng phương pháp giải toán của Chương VI để tìm đáp án cho bài toán: Tìm m để phương trình x$^2$ - 6x + m = 0 có hai nghiệm phân biệt x$_1$, x$_2$ thỏa mãn x$_1$$^2$ + x$_2$$^2$ = 20:\par
\hspace*{1.5em}\textbf{A.} m = -8\quad\quad \textbf{B.} m = 6\quad\quad \textbf{C.} m = 8\quad\quad \textbf{D.} m = 10\par\vspace{2pt}
\textbf{\color{primary} Câu 47.} Xét bài toán vận dụng thực tiễn / nâng cao thuộc Chương VI: Gọi x$_1$, x$_2$ là hai nghiệm của phương trình x$^2$ - 11x + 2 = 0. Giá trị của biểu thức A = x$_1$$^2$ + x$_2$$^2$ bằng:\par
\hspace*{1.5em}\textbf{A.} 121\quad\quad \textbf{B.} 117\quad\quad \textbf{C.} 119\quad\quad \textbf{D.} 125\par\vspace{2pt}
\textbf{\color{primary} Câu 48.} Tìm kết quả cuối cùng của bài toán phân loại học sinh khá giỏi (Chương VI): Tìm m để phương trình x$^2$ - 6x + m = 0 có hai nghiệm phân biệt x$_1$, x$_2$ thỏa mãn x$_1$$^2$ + x$_2$$^2$ = 20:\par
\hspace*{1.5em}\textbf{A.} m = -8\quad\quad \textbf{B.} m = 10\quad\quad \textbf{C.} m = 8\quad\quad \textbf{D.} m = 6\par\vspace{2pt}
\textbf{\color{primary} Câu 49.} Kết hợp các tính chất đã học trong Chương VI, hãy giải quyết bài toán: Gọi x$_1$, x$_2$ là hai nghiệm của phương trình x$^2$ - 13x + 2 = 0. Giá trị của biểu thức A = x$_1$$^2$ + x$_2$$^2$ bằng:\par
\hspace*{1.5em}\textbf{A.} 173\quad\quad \textbf{B.} 167\quad\quad \textbf{C.} 169\quad\quad \textbf{D.} 165\par\vspace{2pt}
\textbf{\color{primary} Câu 50.} Câu hỏi vận dụng cao cuối Chương VI: Tìm m để phương trình x$^2$ - 6x + m = 0 có hai nghiệm phân biệt x$_1$, x$_2$ thỏa mãn x$_1$$^2$ + x$_2$$^2$ = 20. Phương án đúng là:\par
\hspace*{1.5em}\textbf{A.} m = -8\quad\quad \textbf{B.} m = 10\quad\quad \textbf{C.} m = 8\quad\quad \textbf{D.} m = 6\par\vspace{2pt}
\subsection*{\color{primary} Bảng đáp án Chương VI: Hàm số y = ax$^2$ (a $\ne$ 0). Phương trình bậc hai một ẩn}
\begin{center}
\renewcommand{\arraystretch}{1.25}
\begin{tabularx}{\textwidth}{*{10}{>{\centering\arraybackslash}X}}
\toprule
\textbf{\color{primary} Cột 1} & \textbf{\color{primary} Cột 2} & \textbf{\color{primary} Cột 3} & \textbf{\color{primary} Cột 4} & \textbf{\color{primary} Cột 5} & \textbf{\color{primary} Cột 6} & \textbf{\color{primary} Cột 7} & \textbf{\color{primary} Cột 8} & \textbf{\color{primary} Cột 9} & \textbf{\color{primary} Cột 10} \\
\midrule
\textbf{1.A} & \textbf{2.A} & \textbf{3.A} & \textbf{4.A} & \textbf{5.B} & \textbf{6.C} & \textbf{7.A} & \textbf{8.D} & \textbf{9.A} & \textbf{10.D} \\
\rowcolor{rowtint} \textbf{11.C} & \textbf{12.B} & \textbf{13.D} & \textbf{14.C} & \textbf{15.C} & \textbf{16.B} & \textbf{17.B} & \textbf{18.A} & \textbf{19.B} & \textbf{20.B} \\
\textbf{21.B} & \textbf{22.B} & \textbf{23.B} & \textbf{24.B} & \textbf{25.A} & \textbf{26.A} & \textbf{27.C} & \textbf{28.D} & \textbf{29.B} & \textbf{30.A} \\
\rowcolor{rowtint} \textbf{31.B} & \textbf{32.A} & \textbf{33.A} & \textbf{34.A} & \textbf{35.C} & \textbf{36.C} & \textbf{37.D} & \textbf{38.A} & \textbf{39.C} & \textbf{40.B} \\
\textbf{41.B} & \textbf{42.D} & \textbf{43.D} & \textbf{44.D} & \textbf{45.D} & \textbf{46.C} & \textbf{47.B} & \textbf{48.C} & \textbf{49.D} & \textbf{50.C} \\
\bottomrule
\end{tabularx}
\end{center}\vspace{10pt}
\section*{\color{primary} CHƯƠNG VII: TẦN SỐ VÀ TẦN SỐ TƯƠNG ĐỐI}
\textit{Trọng tâm kiến thức SGK: Bảng tần số và biểu đồ tần số, bảng tần số tương đối và biểu đồ tần số tương đối, bảng tần số và tần số tương đối ghép nhóm.}\par\vspace{4pt}
\subsection*{\color{accent} Phần I. Mức độ Dễ (Nhận biết -- Câu 1 đến Câu 20)}
\textbf{\color{primary} Câu 1.} Trong mẫu dữ liệu kích thước N = 40, một giá trị xuất hiện với tần số m = 8. Tần số tương đối f của giá trị đó bằng:\par
\hspace*{1.5em}\textbf{A.} 25\%\quad\quad \textbf{B.} 20\%\quad\quad \textbf{C.} 10\%\quad\quad \textbf{D.} 15\%\par\vspace{2pt}
\textbf{\color{primary} Câu 2.} Thực hiện yêu cầu nhận biết sau: Tổng tất cả các tần số của các giá trị khác nhau trong bảng tần số luôn bằng.\par
\begin{itemize}\setlength{\itemsep}{1pt}\setlength{\parskip}{0pt}
  \item[] \textbf{A.} 100\%
  \item[] \textbf{B.} 1
  \item[] \textbf{C.} Số các giá trị khác nhau
  \item[] \textbf{D.} Kích thước mẫu (tổng số dữ liệu thu thập N)
\end{itemize}\vspace{2pt}
\textbf{\color{primary} Câu 3.} Trong bài học thuộc Chương VII, xét nội dung: Đối với nhóm số liệu nửa khoảng [20; 25), độ dài (bề rộng) của nhóm này bằng. Đáp án đúng là:\par
\hspace*{1.5em}\textbf{A.} 10\quad\quad \textbf{B.} 45\quad\quad \textbf{C.} 5\quad\quad \textbf{D.} 22.5\par\vspace{2pt}
\textbf{\color{primary} Câu 4.} Dựa vào kiến thức chuẩn SGK của Chương VII (Tần số và tần số tương đối), trong mẫu dữ liệu kích thước N = 40, một giá trị xuất hiện với tần số m = 8. Tần số tương đối f của giá trị đó bằng:\par
\hspace*{1.5em}\textbf{A.} 20\%\quad\quad \textbf{B.} 10\%\quad\quad \textbf{C.} 15\%\quad\quad \textbf{D.} 25\%\par\vspace{2pt}
\textbf{\color{primary} Câu 5.} Chọn kết quả đúng nhất khi xét: Tổng tất cả các tần số của các giá trị khác nhau trong bảng tần số luôn bằng:\par
\begin{itemize}\setlength{\itemsep}{1pt}\setlength{\parskip}{0pt}
  \item[] \textbf{A.} Số các giá trị khác nhau
  \item[] \textbf{B.} Kích thước mẫu (tổng số dữ liệu thu thập N)
  \item[] \textbf{C.} 1
  \item[] \textbf{D.} 100\%
\end{itemize}\vspace{2pt}
\textbf{\color{primary} Câu 6.} Khi làm bài tập cơ bản về tần số và tần số tương đối, một học sinh gặp câu hỏi: Đối với nhóm số liệu nửa khoảng [35; 40), độ dài (bề rộng) của nhóm này bằng. Kết quả nào dưới đây đúng?\par
\hspace*{1.5em}\textbf{A.} 37.5\quad\quad \textbf{B.} 75\quad\quad \textbf{C.} 10\quad\quad \textbf{D.} 5\par\vspace{2pt}
\textbf{\color{primary} Câu 7.} Xác định giá trị hoặc khẳng định đúng cho bài toán sau: Trong mẫu dữ liệu kích thước N = 40, một giá trị xuất hiện với tần số m = 8. Tần số tương đối f của giá trị đó bằng:\par
\hspace*{1.5em}\textbf{A.} 20\%\quad\quad \textbf{B.} 25\%\quad\quad \textbf{C.} 15\%\quad\quad \textbf{D.} 10\%\par\vspace{2pt}
\textbf{\color{primary} Câu 8.} Áp dụng trực tiếp định nghĩa và quy tắc trong Chương VII: Tổng tất cả các tần số của các giá trị khác nhau trong bảng tần số luôn bằng:\par
\begin{itemize}\setlength{\itemsep}{1pt}\setlength{\parskip}{0pt}
  \item[] \textbf{A.} Số các giá trị khác nhau
  \item[] \textbf{B.} 100\%
  \item[] \textbf{C.} Kích thước mẫu (tổng số dữ liệu thu thập N)
  \item[] \textbf{D.} 1
\end{itemize}\vspace{2pt}
\textbf{\color{primary} Câu 9.} Kiểm tra nhanh kiến thức nền tảng: Đối với nhóm số liệu nửa khoảng [50; 55), độ dài (bề rộng) của nhóm này bằng:\par
\hspace*{1.5em}\textbf{A.} 52.5\quad\quad \textbf{B.} 105\quad\quad \textbf{C.} 10\quad\quad \textbf{D.} 5\par\vspace{2pt}
\textbf{\color{primary} Câu 10.} Cho dữ kiện nhận biết sau đây: Trong mẫu dữ liệu kích thước N = 40, một giá trị xuất hiện với tần số m = 8. Tần số tương đối f của giá trị đó bằng. Phương án chính xác là:\par
\hspace*{1.5em}\textbf{A.} 20\%\quad\quad \textbf{B.} 10\%\quad\quad \textbf{C.} 15\%\quad\quad \textbf{D.} 25\%\par\vspace{2pt}
\textbf{\color{primary} Câu 11.} Hãy chỉ ra kết quả đúng trong bốn lựa chọn A, B, C, D đối với yêu cầu: Tổng tất cả các tần số của các giá trị khác nhau trong bảng tần số luôn bằng:\par
\begin{itemize}\setlength{\itemsep}{1pt}\setlength{\parskip}{0pt}
  \item[] \textbf{A.} Số các giá trị khác nhau
  \item[] \textbf{B.} Kích thước mẫu (tổng số dữ liệu thu thập N)
  \item[] \textbf{C.} 1
  \item[] \textbf{D.} 100\%
\end{itemize}\vspace{2pt}
\textbf{\color{primary} Câu 12.} Câu hỏi ôn tập lý thuyết và kỹ năng cơ bản (Chương VII): Đối với nhóm số liệu nửa khoảng [65; 70), độ dài (bề rộng) của nhóm này bằng:\par
\hspace*{1.5em}\textbf{A.} 135\quad\quad \textbf{B.} 67.5\quad\quad \textbf{C.} 10\quad\quad \textbf{D.} 5\par\vspace{2pt}
\textbf{\color{primary} Câu 13.} Theo nội dung sách giáo khoa Kết nối tri thức với cuộc sống (Chương VII), trong mẫu dữ liệu kích thước N = 40, một giá trị xuất hiện với tần số m = 8. Tần số tương đối f của giá trị đó bằng:\par
\hspace*{1.5em}\textbf{A.} 10\%\quad\quad \textbf{B.} 20\%\quad\quad \textbf{C.} 15\%\quad\quad \textbf{D.} 25\%\par\vspace{2pt}
\textbf{\color{primary} Câu 14.} Tìm đáp số đúng của bài toán cơ bản sau: Tổng tất cả các tần số của các giá trị khác nhau trong bảng tần số luôn bằng:\par
\begin{itemize}\setlength{\itemsep}{1pt}\setlength{\parskip}{0pt}
  \item[] \textbf{A.} 100\%
  \item[] \textbf{B.} Kích thước mẫu (tổng số dữ liệu thu thập N)
  \item[] \textbf{C.} 1
  \item[] \textbf{D.} Số các giá trị khác nhau
\end{itemize}\vspace{2pt}
\textbf{\color{primary} Câu 15.} Xét đối tượng toán học trong Chương VII: Đối với nhóm số liệu nửa khoảng [80; 85), độ dài (bề rộng) của nhóm này bằng. Giá trị hoặc tính chất tương ứng là:\par
\hspace*{1.5em}\textbf{A.} 82.5\quad\quad \textbf{B.} 5\quad\quad \textbf{C.} 10\quad\quad \textbf{D.} 165\par\vspace{2pt}
\textbf{\color{primary} Câu 16.} Hoàn thành phép xác định cơ bản sau đây: Trong mẫu dữ liệu kích thước N = 40, một giá trị xuất hiện với tần số m = 8. Tần số tương đối f của giá trị đó bằng:\par
\hspace*{1.5em}\textbf{A.} 25\%\quad\quad \textbf{B.} 20\%\quad\quad \textbf{C.} 15\%\quad\quad \textbf{D.} 10\%\par\vspace{2pt}
\textbf{\color{primary} Câu 17.} Trong các khẳng định hoặc giá trị dưới đây, đâu là đáp án đúng cho câu hỏi: Tổng tất cả các tần số của các giá trị khác nhau trong bảng tần số luôn bằng?\par
\begin{itemize}\setlength{\itemsep}{1pt}\setlength{\parskip}{0pt}
  \item[] \textbf{A.} Số các giá trị khác nhau
  \item[] \textbf{B.} 1
  \item[] \textbf{C.} Kích thước mẫu (tổng số dữ liệu thu thập N)
  \item[] \textbf{D.} 100\%
\end{itemize}\vspace{2pt}
\textbf{\color{primary} Câu 18.} Bài tập khởi động Chương VII: Đối với nhóm số liệu nửa khoảng [95; 100), độ dài (bề rộng) của nhóm này bằng:\par
\hspace*{1.5em}\textbf{A.} 97.5\quad\quad \textbf{B.} 5\quad\quad \textbf{C.} 195\quad\quad \textbf{D.} 10\par\vspace{2pt}
\textbf{\color{primary} Câu 19.} Quan sát dữ kiện đề bài: Trong mẫu dữ liệu kích thước N = 40, một giá trị xuất hiện với tần số m = 8. Tần số tương đối f của giá trị đó bằng. Hãy chọn phương án đúng:\par
\hspace*{1.5em}\textbf{A.} 20\%\quad\quad \textbf{B.} 25\%\quad\quad \textbf{C.} 10\%\quad\quad \textbf{D.} 15\%\par\vspace{2pt}
\textbf{\color{primary} Câu 20.} Sử dụng công thức hoặc tính chất đầu bài của Chương VII để giải quyết yêu cầu: Tổng tất cả các tần số của các giá trị khác nhau trong bảng tần số luôn bằng:\par
\begin{itemize}\setlength{\itemsep}{1pt}\setlength{\parskip}{0pt}
  \item[] \textbf{A.} 100\%
  \item[] \textbf{B.} 1
  \item[] \textbf{C.} Số các giá trị khác nhau
  \item[] \textbf{D.} Kích thước mẫu (tổng số dữ liệu thu thập N)
\end{itemize}\vspace{2pt}
\subsection*{\color{accent} Phần II. Mức độ Trung bình (Thông hiểu -- Câu 21 đến Câu 40)}
\textbf{\color{primary} Câu 21.} Khảo sát 50 học sinh, giá trị điểm 9 có tần số tương đối là 18\%. Tần số (số học sinh đạt điểm 9) bằng:\par
\hspace*{1.5em}\textbf{A.} 25 học sinh\quad\quad \textbf{B.} 12 học sinh\quad\quad \textbf{C.} 18 học sinh\quad\quad \textbf{D.} 9 học sinh\par\vspace{2pt}
\textbf{\color{primary} Câu 22.} Vận dụng quy tắc trong Chương VII (Tần số và tần số tương đối) để giải bài toán: Một bảng tần số tương đối gồm 4 nhóm A, B, C, D. Biết tần số tương đối của các nhóm A, B, C lần lượt là 21\%, 30\%, 24\%. Tần số tương đối của nhóm D là:\par
\hspace*{1.5em}\textbf{A.} 20\%\quad\quad \textbf{B.} 30\%\quad\quad \textbf{C.} 25\%\quad\quad \textbf{D.} 15\%\par\vspace{2pt}
\textbf{\color{primary} Câu 23.} Biến đổi và tính toán theo kiến thức Chương VII: Để biểu diễn tần số tương đối ghép nhóm của một mẫu số liệu liên tục, loại biểu đồ nào thường được sử dụng trong SGK Toán 9. Kết quả thu được là:\par
\begin{itemize}\setlength{\itemsep}{1pt}\setlength{\parskip}{0pt}
  \item[] \textbf{A.} Sơ đồ hình cây
  \item[] \textbf{B.} Biểu đồ tần số tương đối ghép nhóm dạng cột (có các cột kề sát nhau) hoặc dạng đoạn thẳng
  \item[] \textbf{C.} Biểu đồ tranh
  \item[] \textbf{D.} Bảng kiểm đếm thô
\end{itemize}\vspace{2pt}
\textbf{\color{primary} Câu 24.} Xét bài toán mức độ thông hiểu sau: Khảo sát 100 học sinh, giá trị điểm 9 có tần số tương đối là 18\%. Tần số (số học sinh đạt điểm 9) bằng. Đáp án đúng là:\par
\hspace*{1.5em}\textbf{A.} 21 học sinh\quad\quad \textbf{B.} 18 học sinh (1)\quad\quad \textbf{C.} 50 học sinh\quad\quad \textbf{D.} 18 học sinh\par\vspace{2pt}
\textbf{\color{primary} Câu 25.} Để hoàn thành phiếu học tập Chương VII, hãy xác định: Một bảng tần số tương đối gồm 4 nhóm A, B, C, D. Biết tần số tương đối của các nhóm A, B, C lần lượt là 24\%, 30\%, 21\%. Tần số tương đối của nhóm D là:\par
\hspace*{1.5em}\textbf{A.} 30\%\quad\quad \textbf{B.} 15\%\quad\quad \textbf{C.} 20\%\quad\quad \textbf{D.} 25\%\par\vspace{2pt}
\textbf{\color{primary} Câu 26.} Phân tích dữ kiện bài toán: Để biểu diễn tần số tương đối ghép nhóm của một mẫu số liệu liên tục, loại biểu đồ nào thường được sử dụng trong SGK Toán 9. Giá trị cần tìm bằng:\par
\begin{itemize}\setlength{\itemsep}{1pt}\setlength{\parskip}{0pt}
  \item[] \textbf{A.} Sơ đồ hình cây
  \item[] \textbf{B.} Biểu đồ tranh
  \item[] \textbf{C.} Biểu đồ tần số tương đối ghép nhóm dạng cột (có các cột kề sát nhau) hoặc dạng đoạn thẳng
  \item[] \textbf{D.} Bảng kiểm đếm thô
\end{itemize}\vspace{2pt}
\textbf{\color{primary} Câu 27.} Áp dụng tính chất đã học ở Chương VII vào bài tập sau: Khảo sát 50 học sinh, giá trị điểm 9 có tần số tương đối là 18\%. Tần số (số học sinh đạt điểm 9) bằng:\par
\hspace*{1.5em}\textbf{A.} 12 học sinh\quad\quad \textbf{B.} 18 học sinh\quad\quad \textbf{C.} 9 học sinh\quad\quad \textbf{D.} 25 học sinh\par\vspace{2pt}
\textbf{\color{primary} Câu 28.} Thực hiện các bước suy luận hợp lý để giải quyết yêu cầu: Một bảng tần số tương đối gồm 4 nhóm A, B, C, D. Biết tần số tương đối của các nhóm A, B, C lần lượt là 27\%, 30\%, 18\%. Tần số tương đối của nhóm D là:\par
\hspace*{1.5em}\textbf{A.} 20\%\quad\quad \textbf{B.} 30\%\quad\quad \textbf{C.} 15\%\quad\quad \textbf{D.} 25\%\par\vspace{2pt}
\textbf{\color{primary} Câu 29.} Cho bài toán tự luyện (Tần số và tần số tương đối): Để biểu diễn tần số tương đối ghép nhóm của một mẫu số liệu liên tục, loại biểu đồ nào thường được sử dụng trong SGK Toán 9. Phương án nào sau đây chính xác?\par
\begin{itemize}\setlength{\itemsep}{1pt}\setlength{\parskip}{0pt}
  \item[] \textbf{A.} Biểu đồ tần số tương đối ghép nhóm dạng cột (có các cột kề sát nhau) hoặc dạng đoạn thẳng
  \item[] \textbf{B.} Sơ đồ hình cây
  \item[] \textbf{C.} Biểu đồ tranh
  \item[] \textbf{D.} Bảng kiểm đếm thô
\end{itemize}\vspace{2pt}
\textbf{\color{primary} Câu 30.} Kết hợp các phép biến đổi trong Chương VII, hãy tìm kết quả của bài toán: Khảo sát 100 học sinh, giá trị điểm 9 có tần số tương đối là 18\%. Tần số (số học sinh đạt điểm 9) bằng:\par
\hspace*{1.5em}\textbf{A.} 50 học sinh\quad\quad \textbf{B.} 21 học sinh\quad\quad \textbf{C.} 18 học sinh (1)\quad\quad \textbf{D.} 18 học sinh\par\vspace{2pt}
\textbf{\color{primary} Câu 31.} Một bài kiểm tra định kỳ nội dung Chương VII nêu câu hỏi: Một bảng tần số tương đối gồm 4 nhóm A, B, C, D. Biết tần số tương đối của các nhóm A, B, C lần lượt là 30\%, 30\%, 15\%. Tần số tương đối của nhóm D là. Hãy chọn đáp số đúng:\par
\hspace*{1.5em}\textbf{A.} 20\%\quad\quad \textbf{B.} 25\%\quad\quad \textbf{C.} 30\%\quad\quad \textbf{D.} 15\%\par\vspace{2pt}
\textbf{\color{primary} Câu 32.} Tìm lời giải đúng cho bài toán thông hiểu: Để biểu diễn tần số tương đối ghép nhóm của một mẫu số liệu liên tục, loại biểu đồ nào thường được sử dụng trong SGK Toán 9:\par
\begin{itemize}\setlength{\itemsep}{1pt}\setlength{\parskip}{0pt}
  \item[] \textbf{A.} Biểu đồ tranh
  \item[] \textbf{B.} Biểu đồ tần số tương đối ghép nhóm dạng cột (có các cột kề sát nhau) hoặc dạng đoạn thẳng
  \item[] \textbf{C.} Sơ đồ hình cây
  \item[] \textbf{D.} Bảng kiểm đếm thô
\end{itemize}\vspace{2pt}
\textbf{\color{primary} Câu 33.} Sử dụng mối liên hệ giữa các đại lượng trong Chương VII: Khảo sát 50 học sinh, giá trị điểm 9 có tần số tương đối là 18\%. Tần số (số học sinh đạt điểm 9) bằng:\par
\hspace*{1.5em}\textbf{A.} 25 học sinh\quad\quad \textbf{B.} 12 học sinh\quad\quad \textbf{C.} 9 học sinh\quad\quad \textbf{D.} 18 học sinh\par\vspace{2pt}
\textbf{\color{primary} Câu 34.} Hãy tính toán và chọn phương án đúng đối với bài tập: Một bảng tần số tương đối gồm 4 nhóm A, B, C, D. Biết tần số tương đối của các nhóm A, B, C lần lượt là 33\%, 30\%, 12\%. Tần số tương đối của nhóm D là:\par
\hspace*{1.5em}\textbf{A.} 30\%\quad\quad \textbf{B.} 25\%\quad\quad \textbf{C.} 20\%\quad\quad \textbf{D.} 15\%\par\vspace{2pt}
\textbf{\color{primary} Câu 35.} Dựa trên các tính chất trọng tâm của Chương VII, xét bài toán: Để biểu diễn tần số tương đối ghép nhóm của một mẫu số liệu liên tục, loại biểu đồ nào thường được sử dụng trong SGK Toán 9:\par
\begin{itemize}\setlength{\itemsep}{1pt}\setlength{\parskip}{0pt}
  \item[] \textbf{A.} Sơ đồ hình cây
  \item[] \textbf{B.} Biểu đồ tranh
  \item[] \textbf{C.} Biểu đồ tần số tương đối ghép nhóm dạng cột (có các cột kề sát nhau) hoặc dạng đoạn thẳng
  \item[] \textbf{D.} Bảng kiểm đếm thô
\end{itemize}\vspace{2pt}
\textbf{\color{primary} Câu 36.} Thực hành giải toán (Tần số và tần số tương đối): Khảo sát 100 học sinh, giá trị điểm 9 có tần số tương đối là 18\%. Tần số (số học sinh đạt điểm 9) bằng. Đáp án nào dưới đây thỏa mãn?\par
\hspace*{1.5em}\textbf{A.} 18 học sinh (1)\quad\quad \textbf{B.} 18 học sinh\quad\quad \textbf{C.} 50 học sinh\quad\quad \textbf{D.} 21 học sinh\par\vspace{2pt}
\textbf{\color{primary} Câu 37.} Kiểm tra kỹ năng biến đổi và suy luận qua bài toán: Một bảng tần số tương đối gồm 4 nhóm A, B, C, D. Biết tần số tương đối của các nhóm A, B, C lần lượt là 36\%, 30\%, 9\%. Tần số tương đối của nhóm D là:\par
\hspace*{1.5em}\textbf{A.} 15\%\quad\quad \textbf{B.} 30\%\quad\quad \textbf{C.} 20\%\quad\quad \textbf{D.} 25\%\par\vspace{2pt}
\textbf{\color{primary} Câu 38.} Cho giả thiết thuộc phạm vi Chương VII: Để biểu diễn tần số tương đối ghép nhóm của một mẫu số liệu liên tục, loại biểu đồ nào thường được sử dụng trong SGK Toán 9. Kết luận đúng là:\par
\begin{itemize}\setlength{\itemsep}{1pt}\setlength{\parskip}{0pt}
  \item[] \textbf{A.} Biểu đồ tần số tương đối ghép nhóm dạng cột (có các cột kề sát nhau) hoặc dạng đoạn thẳng
  \item[] \textbf{B.} Sơ đồ hình cây
  \item[] \textbf{C.} Biểu đồ tranh
  \item[] \textbf{D.} Bảng kiểm đếm thô
\end{itemize}\vspace{2pt}
\textbf{\color{primary} Câu 39.} Rút gọn hoặc tính giá trị theo yêu cầu sau của Chương VII: Khảo sát 50 học sinh, giá trị điểm 9 có tần số tương đối là 18\%. Tần số (số học sinh đạt điểm 9) bằng:\par
\hspace*{1.5em}\textbf{A.} 9 học sinh\quad\quad \textbf{B.} 12 học sinh\quad\quad \textbf{C.} 25 học sinh\quad\quad \textbf{D.} 18 học sinh\par\vspace{2pt}
\textbf{\color{primary} Câu 40.} Xác định phương án đúng sau khi thực hiện phép tính / suy luận: Một bảng tần số tương đối gồm 4 nhóm A, B, C, D. Biết tần số tương đối của các nhóm A, B, C lần lượt là 39\%, 30\%, 6\%. Tần số tương đối của nhóm D là:\par
\hspace*{1.5em}\textbf{A.} 25\%\quad\quad \textbf{B.} 20\%\quad\quad \textbf{C.} 30\%\quad\quad \textbf{D.} 15\%\par\vspace{2pt}
\subsection*{\color{accent} Phần III. Mức độ Khó (Vận dụng \& Vận dụng cao -- Câu 41 đến Câu 50)}
\textbf{\color{primary} Câu 41.} Trong một mẫu số liệu gồm 4 giá trị có tần số lần lượt tỉ lệ với 1 : 2 : 3 : 4. Tần số tương đối của giá trị có tần số lớn nhất bằng:\par
\hspace*{1.5em}\textbf{A.} 25\%\quad\quad \textbf{B.} 30\%\quad\quad \textbf{C.} 40\%\quad\quad \textbf{D.} 50\%\par\vspace{2pt}
\textbf{\color{primary} Câu 42.} Bài toán vận dụng (Chương VII – Tần số và tần số tương đối): Khảo sát thời gian chạy của một nhóm học sinh chia thành 3 nhóm: [12; 14) có 8 bạn, [14; 16) có 18 bạn, [16; 18) có 6 bạn. Tần số tương đối của nhóm [14; 16) bằng. Kết quả đúng là:\par
\hspace*{1.5em}\textbf{A.} 44\%\quad\quad \textbf{B.} 60\%\quad\quad \textbf{C.} 50\%\quad\quad \textbf{D.} 56\%\par\vspace{2pt}
\textbf{\color{primary} Câu 43.} Vận dụng kiến thức nâng cao trong phạm vi Chương VII để giải bài toán: Trong một mẫu số liệu gồm 4 giá trị có tần số lần lượt tỉ lệ với 1 : 2 : 3 : 4. Tần số tương đối của giá trị có tần số lớn nhất bằng:\par
\hspace*{1.5em}\textbf{A.} 50\%\quad\quad \textbf{B.} 30\%\quad\quad \textbf{C.} 40\%\quad\quad \textbf{D.} 25\%\par\vspace{2pt}
\textbf{\color{primary} Câu 44.} Phân tích và giải bài toán tổng hợp sau của Chương VII: Khảo sát thời gian chạy của một nhóm học sinh chia thành 3 nhóm: [12; 14) có 12 bạn, [14; 16) có 26 bạn, [16; 18) có 8 bạn. Tần số tương đối của nhóm [14; 16) bằng:\par
\hspace*{1.5em}\textbf{A.} 50\%\quad\quad \textbf{B.} 60\%\quad\quad \textbf{C.} 44\%\quad\quad \textbf{D.} 56\%\par\vspace{2pt}
\textbf{\color{primary} Câu 45.} Thử thách tư duy trong phạm vi kiến thức Chương VII: Trong một mẫu số liệu gồm 4 giá trị có tần số lần lượt tỉ lệ với 1 : 2 : 3 : 4. Tần số tương đối của giá trị có tần số lớn nhất bằng. Đáp số chính xác là:\par
\hspace*{1.5em}\textbf{A.} 50\%\quad\quad \textbf{B.} 30\%\quad\quad \textbf{C.} 40\%\quad\quad \textbf{D.} 25\%\par\vspace{2pt}
\textbf{\color{primary} Câu 46.} Sử dụng phương pháp giải toán của Chương VII để tìm đáp án cho bài toán: Khảo sát thời gian chạy của một nhóm học sinh chia thành 3 nhóm: [12; 14) có 16 bạn, [14; 16) có 34 bạn, [16; 18) có 10 bạn. Tần số tương đối của nhóm [14; 16) bằng:\par
\hspace*{1.5em}\textbf{A.} 56\%\quad\quad \textbf{B.} 50\%\quad\quad \textbf{C.} 60\%\quad\quad \textbf{D.} 44\%\par\vspace{2pt}
\textbf{\color{primary} Câu 47.} Xét bài toán vận dụng thực tiễn / nâng cao thuộc Chương VII: Trong một mẫu số liệu gồm 4 giá trị có tần số lần lượt tỉ lệ với 1 : 2 : 3 : 4. Tần số tương đối của giá trị có tần số lớn nhất bằng:\par
\hspace*{1.5em}\textbf{A.} 40\%\quad\quad \textbf{B.} 50\%\quad\quad \textbf{C.} 30\%\quad\quad \textbf{D.} 25\%\par\vspace{2pt}
\textbf{\color{primary} Câu 48.} Tìm kết quả cuối cùng của bài toán phân loại học sinh khá giỏi (Chương VII): Khảo sát thời gian chạy của một nhóm học sinh chia thành 3 nhóm: [12; 14) có 20 bạn, [14; 16) có 42 bạn, [16; 18) có 12 bạn. Tần số tương đối của nhóm [14; 16) bằng:\par
\hspace*{1.5em}\textbf{A.} 56\%\quad\quad \textbf{B.} 44\%\quad\quad \textbf{C.} 50\%\quad\quad \textbf{D.} 60\%\par\vspace{2pt}
\textbf{\color{primary} Câu 49.} Kết hợp các tính chất đã học trong Chương VII, hãy giải quyết bài toán: Trong một mẫu số liệu gồm 4 giá trị có tần số lần lượt tỉ lệ với 1 : 2 : 3 : 4. Tần số tương đối của giá trị có tần số lớn nhất bằng:\par
\hspace*{1.5em}\textbf{A.} 40\%\quad\quad \textbf{B.} 25\%\quad\quad \textbf{C.} 30\%\quad\quad \textbf{D.} 50\%\par\vspace{2pt}
\textbf{\color{primary} Câu 50.} Câu hỏi vận dụng cao cuối Chương VII: Khảo sát thời gian chạy của một nhóm học sinh chia thành 3 nhóm: [12; 14) có 24 bạn, [14; 16) có 50 bạn, [16; 18) có 14 bạn. Tần số tương đối của nhóm [14; 16) bằng. Phương án đúng là:\par
\hspace*{1.5em}\textbf{A.} 56\%\quad\quad \textbf{B.} 60\%\quad\quad \textbf{C.} 50\%\quad\quad \textbf{D.} 44\%\par\vspace{2pt}
\subsection*{\color{primary} Bảng đáp án Chương VII: Tần số và tần số tương đối}
\begin{center}
\renewcommand{\arraystretch}{1.25}
\begin{tabularx}{\textwidth}{*{10}{>{\centering\arraybackslash}X}}
\toprule
\textbf{\color{primary} Cột 1} & \textbf{\color{primary} Cột 2} & \textbf{\color{primary} Cột 3} & \textbf{\color{primary} Cột 4} & \textbf{\color{primary} Cột 5} & \textbf{\color{primary} Cột 6} & \textbf{\color{primary} Cột 7} & \textbf{\color{primary} Cột 8} & \textbf{\color{primary} Cột 9} & \textbf{\color{primary} Cột 10} \\
\midrule
\textbf{1.B} & \textbf{2.D} & \textbf{3.C} & \textbf{4.A} & \textbf{5.B} & \textbf{6.D} & \textbf{7.A} & \textbf{8.C} & \textbf{9.D} & \textbf{10.A} \\
\rowcolor{rowtint} \textbf{11.B} & \textbf{12.D} & \textbf{13.B} & \textbf{14.B} & \textbf{15.B} & \textbf{16.B} & \textbf{17.C} & \textbf{18.B} & \textbf{19.A} & \textbf{20.D} \\
\textbf{21.D} & \textbf{22.C} & \textbf{23.B} & \textbf{24.D} & \textbf{25.D} & \textbf{26.C} & \textbf{27.C} & \textbf{28.D} & \textbf{29.A} & \textbf{30.D} \\
\rowcolor{rowtint} \textbf{31.B} & \textbf{32.B} & \textbf{33.C} & \textbf{34.B} & \textbf{35.C} & \textbf{36.B} & \textbf{37.D} & \textbf{38.A} & \textbf{39.A} & \textbf{40.A} \\
\textbf{41.C} & \textbf{42.D} & \textbf{43.C} & \textbf{44.D} & \textbf{45.C} & \textbf{46.A} & \textbf{47.A} & \textbf{48.A} & \textbf{49.A} & \textbf{50.A} \\
\bottomrule
\end{tabularx}
\end{center}\vspace{10pt}
\section*{\color{primary} CHƯƠNG VIII: XÁC SUẤT CỦA BIẾN CỐ TRONG MỘT SỐ MÔ HÌNH XÁC SUẤT ĐƠN GIẢN}
\textit{Trọng tâm kiến thức SGK: Không gian mẫu và biến cố trong các phép thử ngẫu nhiên đơn giản, tính xác suất của biến cố liên quan đến phép thử gồm một hoặc hai bước.}\par\vspace{4pt}
\subsection*{\color{accent} Phần I. Mức độ Dễ (Nhận biết -- Câu 1 đến Câu 20)}
\textbf{\color{primary} Câu 1.} Gieo đồng thời hai con xúc xắc cân đối và đồng chất. Số phần tử của không gian mẫu $\Omega$ là:\par
\hspace*{1.5em}\textbf{A.} 18 phần tử\quad\quad \textbf{B.} 12 phần tử\quad\quad \textbf{C.} 36 phần tử\quad\quad \textbf{D.} 6 phần tử\par\vspace{2pt}
\textbf{\color{primary} Câu 2.} Hãy trả lời câu hỏi cơ bản sau: Tung một đồng xu cân đối hai lần liên tiếp. Không gian mẫu $\Omega$ có bao nhiêu phần tử?\par
\hspace*{1.5em}\textbf{A.} 6 phần tử\quad\quad \textbf{B.} 8 phần tử\quad\quad \textbf{C.} 4 phần tử\quad\quad \textbf{D.} 2 phần tử\par\vspace{2pt}
\textbf{\color{primary} Câu 3.} Trong bài học thuộc Chương VIII, xét nội dung: Một hộp có 7 quả bóng được đánh số từ 1 đến 7. Lấy ngẫu nhiên 1 quả bóng. Không gian mẫu có số phần tử là. Đáp án đúng là:\par
\hspace*{1.5em}\textbf{A.} 14 phần tử\quad\quad \textbf{B.} 7 phần tử\quad\quad \textbf{C.} 49 phần tử\quad\quad \textbf{D.} 1 phần tử\par\vspace{2pt}
\textbf{\color{primary} Câu 4.} Dựa vào kiến thức chuẩn SGK của Chương VIII (Xác suất của biến cố trong một số mô hình xác suất đơn giản), gieo đồng thời hai con xúc xắc cân đối và đồng chất. Số phần tử của không gian mẫu $\Omega$ là:\par
\hspace*{1.5em}\textbf{A.} 12 phần tử\quad\quad \textbf{B.} 18 phần tử\quad\quad \textbf{C.} 6 phần tử\quad\quad \textbf{D.} 36 phần tử\par\vspace{2pt}
\textbf{\color{primary} Câu 5.} Chọn phương án trả lời đúng nhất: Tung một đồng xu cân đối hai lần liên tiếp. Không gian mẫu $\Omega$ có bao nhiêu phần tử?\par
\hspace*{1.5em}\textbf{A.} 6 phần tử\quad\quad \textbf{B.} 2 phần tử\quad\quad \textbf{C.} 4 phần tử\quad\quad \textbf{D.} 8 phần tử\par\vspace{2pt}
\textbf{\color{primary} Câu 6.} Khi làm bài tập cơ bản về xác suất của biến cố trong một số mô hình xác suất đơn giản, một học sinh gặp câu hỏi: Một hộp có 10 quả bóng được đánh số từ 1 đến 10. Lấy ngẫu nhiên 1 quả bóng. Không gian mẫu có số phần tử là. Kết quả nào dưới đây đúng?\par
\hspace*{1.5em}\textbf{A.} 1 phần tử\quad\quad \textbf{B.} 10 phần tử\quad\quad \textbf{C.} 20 phần tử\quad\quad \textbf{D.} 100 phần tử\par\vspace{2pt}
\textbf{\color{primary} Câu 7.} Xác định giá trị hoặc khẳng định đúng cho bài toán sau: Gieo đồng thời hai con xúc xắc cân đối và đồng chất. Số phần tử của không gian mẫu $\Omega$ là:\par
\hspace*{1.5em}\textbf{A.} 36 phần tử\quad\quad \textbf{B.} 18 phần tử\quad\quad \textbf{C.} 12 phần tử\quad\quad \textbf{D.} 6 phần tử\par\vspace{2pt}
\textbf{\color{primary} Câu 8.} Áp dụng trực tiếp định nghĩa và quy tắc trong Chương VIII: Tung một đồng xu cân đối hai lần liên tiếp. Không gian mẫu $\Omega$ có bao nhiêu phần tử:\par
\hspace*{1.5em}\textbf{A.} 6 phần tử\quad\quad \textbf{B.} 4 phần tử\quad\quad \textbf{C.} 8 phần tử\quad\quad \textbf{D.} 2 phần tử\par\vspace{2pt}
\textbf{\color{primary} Câu 9.} Kiểm tra nhanh kiến thức nền tảng: Một hộp có 13 quả bóng được đánh số từ 1 đến 13. Lấy ngẫu nhiên 1 quả bóng. Không gian mẫu có số phần tử là:\par
\hspace*{1.5em}\textbf{A.} 13 phần tử\quad\quad \textbf{B.} 26 phần tử\quad\quad \textbf{C.} 169 phần tử\quad\quad \textbf{D.} 1 phần tử\par\vspace{2pt}
\textbf{\color{primary} Câu 10.} Cho dữ kiện nhận biết sau đây: Gieo đồng thời hai con xúc xắc cân đối và đồng chất. Số phần tử của không gian mẫu $\Omega$ là. Phương án chính xác là:\par
\hspace*{1.5em}\textbf{A.} 36 phần tử\quad\quad \textbf{B.} 6 phần tử\quad\quad \textbf{C.} 18 phần tử\quad\quad \textbf{D.} 12 phần tử\par\vspace{2pt}
\textbf{\color{primary} Câu 11.} Hãy chỉ ra kết quả đúng trong bốn lựa chọn A, B, C, D đối với yêu cầu: Tung một đồng xu cân đối hai lần liên tiếp. Không gian mẫu $\Omega$ có bao nhiêu phần tử:\par
\hspace*{1.5em}\textbf{A.} 2 phần tử\quad\quad \textbf{B.} 6 phần tử\quad\quad \textbf{C.} 8 phần tử\quad\quad \textbf{D.} 4 phần tử\par\vspace{2pt}
\textbf{\color{primary} Câu 12.} Câu hỏi ôn tập lý thuyết và kỹ năng cơ bản (Chương VIII): Một hộp có 16 quả bóng được đánh số từ 1 đến 16. Lấy ngẫu nhiên 1 quả bóng. Không gian mẫu có số phần tử là:\par
\hspace*{1.5em}\textbf{A.} 32 phần tử\quad\quad \textbf{B.} 256 phần tử\quad\quad \textbf{C.} 16 phần tử\quad\quad \textbf{D.} 1 phần tử\par\vspace{2pt}
\textbf{\color{primary} Câu 13.} Theo nội dung sách giáo khoa Kết nối tri thức với cuộc sống (Chương VIII), gieo đồng thời hai con xúc xắc cân đối và đồng chất. Số phần tử của không gian mẫu $\Omega$ là:\par
\hspace*{1.5em}\textbf{A.} 12 phần tử\quad\quad \textbf{B.} 18 phần tử\quad\quad \textbf{C.} 36 phần tử\quad\quad \textbf{D.} 6 phần tử\par\vspace{2pt}
\textbf{\color{primary} Câu 14.} Tìm đáp số đúng của bài toán cơ bản sau: Tung một đồng xu cân đối hai lần liên tiếp. Không gian mẫu $\Omega$ có bao nhiêu phần tử:\par
\hspace*{1.5em}\textbf{A.} 4 phần tử\quad\quad \textbf{B.} 2 phần tử\quad\quad \textbf{C.} 8 phần tử\quad\quad \textbf{D.} 6 phần tử\par\vspace{2pt}
\textbf{\color{primary} Câu 15.} Xét đối tượng toán học trong Chương VIII: Một hộp có 19 quả bóng được đánh số từ 1 đến 19. Lấy ngẫu nhiên 1 quả bóng. Không gian mẫu có số phần tử là. Giá trị hoặc tính chất tương ứng là:\par
\hspace*{1.5em}\textbf{A.} 38 phần tử\quad\quad \textbf{B.} 361 phần tử\quad\quad \textbf{C.} 19 phần tử\quad\quad \textbf{D.} 1 phần tử\par\vspace{2pt}
\textbf{\color{primary} Câu 16.} Hoàn thành phép xác định cơ bản sau đây: Gieo đồng thời hai con xúc xắc cân đối và đồng chất. Số phần tử của không gian mẫu $\Omega$ là:\par
\hspace*{1.5em}\textbf{A.} 18 phần tử\quad\quad \textbf{B.} 12 phần tử\quad\quad \textbf{C.} 6 phần tử\quad\quad \textbf{D.} 36 phần tử\par\vspace{2pt}
\textbf{\color{primary} Câu 17.} Trong các khẳng định hoặc giá trị dưới đây, đâu là đáp án đúng cho câu hỏi: Tung một đồng xu cân đối hai lần liên tiếp. Không gian mẫu $\Omega$ có bao nhiêu phần tử?\par
\hspace*{1.5em}\textbf{A.} 6 phần tử\quad\quad \textbf{B.} 8 phần tử\quad\quad \textbf{C.} 4 phần tử\quad\quad \textbf{D.} 2 phần tử\par\vspace{2pt}
\textbf{\color{primary} Câu 18.} Bài tập khởi động Chương VIII: Một hộp có 22 quả bóng được đánh số từ 1 đến 22. Lấy ngẫu nhiên 1 quả bóng. Không gian mẫu có số phần tử là:\par
\hspace*{1.5em}\textbf{A.} 484 phần tử\quad\quad \textbf{B.} 44 phần tử\quad\quad \textbf{C.} 22 phần tử\quad\quad \textbf{D.} 1 phần tử\par\vspace{2pt}
\textbf{\color{primary} Câu 19.} Quan sát dữ kiện đề bài: Gieo đồng thời hai con xúc xắc cân đối và đồng chất. Số phần tử của không gian mẫu $\Omega$ là. Hãy chọn phương án đúng:\par
\hspace*{1.5em}\textbf{A.} 12 phần tử\quad\quad \textbf{B.} 36 phần tử\quad\quad \textbf{C.} 6 phần tử\quad\quad \textbf{D.} 18 phần tử\par\vspace{2pt}
\textbf{\color{primary} Câu 20.} Sử dụng công thức hoặc tính chất đầu bài của Chương VIII để giải quyết yêu cầu: Tung một đồng xu cân đối hai lần liên tiếp. Không gian mẫu $\Omega$ có bao nhiêu phần tử:\par
\hspace*{1.5em}\textbf{A.} 8 phần tử\quad\quad \textbf{B.} 4 phần tử\quad\quad \textbf{C.} 2 phần tử\quad\quad \textbf{D.} 6 phần tử\par\vspace{2pt}
\subsection*{\color{accent} Phần II. Mức độ Trung bình (Thông hiểu -- Câu 21 đến Câu 40)}
\textbf{\color{primary} Câu 21.} Gieo đồng thời hai con xúc xắc cân đối. Xác suất của biến cố 'Hai con xúc xắc xuất hiện số chấm bằng nhau' bằng:\par
\hspace*{1.5em}\textbf{A.} 1/12\quad\quad \textbf{B.} 1/36\quad\quad \textbf{C.} 1/3\quad\quad \textbf{D.} 1/6\par\vspace{2pt}
\textbf{\color{primary} Câu 22.} Vận dụng quy tắc trong Chương VIII (Xác suất của biến cố trong một số mô hình xác suất đơn giản) để giải bài toán: Gieo đồng thời hai con xúc xắc cân đối. Số kết quả thuận lợi cho biến cố 'Tổng số chấm xuất hiện trên hai con xúc xắc bằng 7' là:\par
\hspace*{1.5em}\textbf{A.} 5 kết quả\quad\quad \textbf{B.} 12 kết quả\quad\quad \textbf{C.} 6 kết quả\quad\quad \textbf{D.} 7 kết quả\par\vspace{2pt}
\textbf{\color{primary} Câu 23.} Biến đổi và tính toán theo kiến thức Chương VIII: Có hai túi: túi I đựng 2 thẻ ghi số 1; 2 và túi II đựng 3 thẻ ghi số 3; 4; 5. Lấy ngẫu nhiên mỗi túi 1 thẻ. Xác suất để tích hai số trên hai thẻ là số chẵn bằng. Kết quả thu được là:\par
\hspace*{1.5em}\textbf{A.} 2/3\quad\quad \textbf{B.} 1/2\quad\quad \textbf{C.} 5/6\quad\quad \textbf{D.} 1/3\par\vspace{2pt}
\textbf{\color{primary} Câu 24.} Xét bài toán mức độ thông hiểu sau: Hộp thứ nhất có 5 quả bóng và hộp thứ hai có 6 quả bóng (các quả bóng đều khác nhau). Lấy ngẫu nhiên từ mỗi hộp 1 quả bóng. Số phần tử của không gian mẫu $\Omega$ là. Đáp án đúng là:\par
\hspace*{1.5em}\textbf{A.} 11 phần tử\quad\quad \textbf{B.} 60 phần tử\quad\quad \textbf{C.} 25 phần tử\quad\quad \textbf{D.} 30 phần tử\par\vspace{2pt}
\textbf{\color{primary} Câu 25.} Để hoàn thành phiếu học tập Chương VIII, hãy xác định: Gieo đồng thời hai con xúc xắc cân đối. Xác suất của biến cố 'Hai con xúc xắc xuất hiện số chấm bằng nhau' bằng:\par
\hspace*{1.5em}\textbf{A.} 1/36\quad\quad \textbf{B.} 1/3\quad\quad \textbf{C.} 1/12\quad\quad \textbf{D.} 1/6\par\vspace{2pt}
\textbf{\color{primary} Câu 26.} Phân tích dữ kiện bài toán: Gieo đồng thời hai con xúc xắc cân đối. Số kết quả thuận lợi cho biến cố 'Tổng số chấm xuất hiện trên hai con xúc xắc bằng 7' là. Giá trị cần tìm bằng:\par
\hspace*{1.5em}\textbf{A.} 12 kết quả\quad\quad \textbf{B.} 6 kết quả\quad\quad \textbf{C.} 5 kết quả\quad\quad \textbf{D.} 7 kết quả\par\vspace{2pt}
\textbf{\color{primary} Câu 27.} Áp dụng tính chất đã học ở Chương VIII vào bài tập sau: Có hai túi: túi I đựng 2 thẻ ghi số 1; 2 và túi II đựng 3 thẻ ghi số 3; 4; 5. Lấy ngẫu nhiên mỗi túi 1 thẻ. Xác suất để tích hai số trên hai thẻ là số chẵn bằng:\par
\hspace*{1.5em}\textbf{A.} 5/6\quad\quad \textbf{B.} 2/3\quad\quad \textbf{C.} 1/2\quad\quad \textbf{D.} 1/3\par\vspace{2pt}
\textbf{\color{primary} Câu 28.} Thực hiện các bước suy luận hợp lý để giải quyết yêu cầu: Hộp thứ nhất có 9 quả bóng và hộp thứ hai có 10 quả bóng (các quả bóng đều khác nhau). Lấy ngẫu nhiên từ mỗi hộp 1 quả bóng. Số phần tử của không gian mẫu $\Omega$ là:\par
\hspace*{1.5em}\textbf{A.} 180 phần tử\quad\quad \textbf{B.} 19 phần tử\quad\quad \textbf{C.} 81 phần tử\quad\quad \textbf{D.} 90 phần tử\par\vspace{2pt}
\textbf{\color{primary} Câu 29.} Cho bài toán tự luyện (Xác suất của biến cố trong một số mô hình xác suất đơn giản): Gieo đồng thời hai con xúc xắc cân đối. Xác suất của biến cố 'Hai con xúc xắc xuất hiện số chấm bằng nhau' bằng. Phương án nào sau đây chính xác?\par
\hspace*{1.5em}\textbf{A.} 1/36\quad\quad \textbf{B.} 1/6\quad\quad \textbf{C.} 1/12\quad\quad \textbf{D.} 1/3\par\vspace{2pt}
\textbf{\color{primary} Câu 30.} Kết hợp các phép biến đổi trong Chương VIII, hãy tìm kết quả của bài toán: Gieo đồng thời hai con xúc xắc cân đối. Số kết quả thuận lợi cho biến cố 'Tổng số chấm xuất hiện trên hai con xúc xắc bằng 7' là:\par
\hspace*{1.5em}\textbf{A.} 7 kết quả\quad\quad \textbf{B.} 12 kết quả\quad\quad \textbf{C.} 6 kết quả\quad\quad \textbf{D.} 5 kết quả\par\vspace{2pt}
\textbf{\color{primary} Câu 31.} Một bài kiểm tra định kỳ nội dung Chương VIII nêu câu hỏi: Có hai túi: túi I đựng 2 thẻ ghi số 1; 2 và túi II đựng 3 thẻ ghi số 3; 4; 5. Lấy ngẫu nhiên mỗi túi 1 thẻ. Xác suất để tích hai số trên hai thẻ là số chẵn bằng. Hãy chọn đáp số đúng:\par
\hspace*{1.5em}\textbf{A.} 5/6\quad\quad \textbf{B.} 1/3\quad\quad \textbf{C.} 2/3\quad\quad \textbf{D.} 1/2\par\vspace{2pt}
\textbf{\color{primary} Câu 32.} Tìm lời giải đúng cho bài toán thông hiểu: Hộp thứ nhất có 13 quả bóng và hộp thứ hai có 14 quả bóng (các quả bóng đều khác nhau). Lấy ngẫu nhiên từ mỗi hộp 1 quả bóng. Số phần tử của không gian mẫu $\Omega$ là:\par
\hspace*{1.5em}\textbf{A.} 169 phần tử\quad\quad \textbf{B.} 27 phần tử\quad\quad \textbf{C.} 364 phần tử\quad\quad \textbf{D.} 182 phần tử\par\vspace{2pt}
\textbf{\color{primary} Câu 33.} Sử dụng mối liên hệ giữa các đại lượng trong Chương VIII: Gieo đồng thời hai con xúc xắc cân đối. Xác suất của biến cố 'Hai con xúc xắc xuất hiện số chấm bằng nhau' bằng:\par
\hspace*{1.5em}\textbf{A.} 1/36\quad\quad \textbf{B.} 1/6\quad\quad \textbf{C.} 1/3\quad\quad \textbf{D.} 1/12\par\vspace{2pt}
\textbf{\color{primary} Câu 34.} Hãy tính toán và chọn phương án đúng đối với bài tập: Gieo đồng thời hai con xúc xắc cân đối. Số kết quả thuận lợi cho biến cố 'Tổng số chấm xuất hiện trên hai con xúc xắc bằng 7' là:\par
\hspace*{1.5em}\textbf{A.} 7 kết quả\quad\quad \textbf{B.} 5 kết quả\quad\quad \textbf{C.} 6 kết quả\quad\quad \textbf{D.} 12 kết quả\par\vspace{2pt}
\textbf{\color{primary} Câu 35.} Dựa trên các tính chất trọng tâm của Chương VIII, xét bài toán: Có hai túi: túi I đựng 2 thẻ ghi số 1; 2 và túi II đựng 3 thẻ ghi số 3; 4; 5. Lấy ngẫu nhiên mỗi túi 1 thẻ. Xác suất để tích hai số trên hai thẻ là số chẵn bằng:\par
\hspace*{1.5em}\textbf{A.} 1/2\quad\quad \textbf{B.} 2/3\quad\quad \textbf{C.} 5/6\quad\quad \textbf{D.} 1/3\par\vspace{2pt}
\textbf{\color{primary} Câu 36.} Thực hành giải toán (Xác suất của biến cố trong một số mô hình xác suất đơn giản): Hộp thứ nhất có 17 quả bóng và hộp thứ hai có 18 quả bóng (các quả bóng đều khác nhau). Lấy ngẫu nhiên từ mỗi hộp 1 quả bóng. Số phần tử của không gian mẫu $\Omega$ là. Đáp án nào dưới đây thỏa mãn?\par
\hspace*{1.5em}\textbf{A.} 289 phần tử\quad\quad \textbf{B.} 306 phần tử\quad\quad \textbf{C.} 612 phần tử\quad\quad \textbf{D.} 35 phần tử\par\vspace{2pt}
\textbf{\color{primary} Câu 37.} Kiểm tra kỹ năng biến đổi và suy luận qua bài toán: Gieo đồng thời hai con xúc xắc cân đối. Xác suất của biến cố 'Hai con xúc xắc xuất hiện số chấm bằng nhau' bằng:\par
\hspace*{1.5em}\textbf{A.} 1/12\quad\quad \textbf{B.} 1/36\quad\quad \textbf{C.} 1/3\quad\quad \textbf{D.} 1/6\par\vspace{2pt}
\textbf{\color{primary} Câu 38.} Cho giả thiết thuộc phạm vi Chương VIII: Gieo đồng thời hai con xúc xắc cân đối. Số kết quả thuận lợi cho biến cố 'Tổng số chấm xuất hiện trên hai con xúc xắc bằng 7' là. Kết luận đúng là:\par
\hspace*{1.5em}\textbf{A.} 6 kết quả\quad\quad \textbf{B.} 12 kết quả\quad\quad \textbf{C.} 7 kết quả\quad\quad \textbf{D.} 5 kết quả\par\vspace{2pt}
\textbf{\color{primary} Câu 39.} Rút gọn hoặc tính giá trị theo yêu cầu sau của Chương VIII: Có hai túi: túi I đựng 2 thẻ ghi số 1; 2 và túi II đựng 3 thẻ ghi số 3; 4; 5. Lấy ngẫu nhiên mỗi túi 1 thẻ. Xác suất để tích hai số trên hai thẻ là số chẵn bằng:\par
\hspace*{1.5em}\textbf{A.} 1/3\quad\quad \textbf{B.} 2/3\quad\quad \textbf{C.} 1/2\quad\quad \textbf{D.} 5/6\par\vspace{2pt}
\textbf{\color{primary} Câu 40.} Xác định phương án đúng sau khi thực hiện phép tính / suy luận: Hộp thứ nhất có 21 quả bóng và hộp thứ hai có 22 quả bóng (các quả bóng đều khác nhau). Lấy ngẫu nhiên từ mỗi hộp 1 quả bóng. Số phần tử của không gian mẫu $\Omega$ là:\par
\hspace*{1.5em}\textbf{A.} 924 phần tử\quad\quad \textbf{B.} 462 phần tử\quad\quad \textbf{C.} 441 phần tử\quad\quad \textbf{D.} 43 phần tử\par\vspace{2pt}
\subsection*{\color{accent} Phần III. Mức độ Khó (Vận dụng \& Vận dụng cao -- Câu 41 đến Câu 50)}
\textbf{\color{primary} Câu 41.} Từ một nhóm gồm 3 học sinh nam và 2 học sinh nữ, chọn ngẫu nhiên đồng thời 2 học sinh. Xác suất để chọn được đúng 1 học sinh nam và 1 học sinh nữ bằng:\par
\hspace*{1.5em}\textbf{A.} 3/10\quad\quad \textbf{B.} 3/5\quad\quad \textbf{C.} 1/2\quad\quad \textbf{D.} 2/5\par\vspace{2pt}
\textbf{\color{primary} Câu 42.} Bài toán vận dụng (Chương VIII – Xác suất của biến cố trong một số mô hình xác suất đơn giản): Gieo đồng thời hai con xúc xắc cân đối và đồng chất. Xác suất của biến cố 'Tổng số chấm xuất hiện trên hai con xúc xắc lớn hơn hoặc bằng 10' bằng. Kết quả đúng là:\par
\hspace*{1.5em}\textbf{A.} 7/36\quad\quad \textbf{B.} 1/6\quad\quad \textbf{C.} 1/4\quad\quad \textbf{D.} 5/36\par\vspace{2pt}
\textbf{\color{primary} Câu 43.} Vận dụng kiến thức nâng cao trong phạm vi Chương VIII để giải bài toán: Từ hai túi: túi I chứa các thẻ ghi số 1; 2; 3 và túi II chứa 5 thẻ ghi các số từ 1 đến 5. Rút ngẫu nhiên mỗi túi 1 thẻ, xác suất để cả hai thẻ rút ra đều ghi số 1 là:\par
\hspace*{1.5em}\textbf{A.} 1/8\quad\quad \textbf{B.} 1/5\quad\quad \textbf{C.} 2/15\quad\quad \textbf{D.} 1/15\par\vspace{2pt}
\textbf{\color{primary} Câu 44.} Phân tích và giải bài toán tổng hợp sau của Chương VIII: Từ một nhóm gồm 3 học sinh nam và 2 học sinh nữ, chọn ngẫu nhiên đồng thời 2 học sinh. Xác suất để chọn được đúng 1 học sinh nam và 1 học sinh nữ bằng:\par
\hspace*{1.5em}\textbf{A.} 2/5\quad\quad \textbf{B.} 1/2\quad\quad \textbf{C.} 3/5\quad\quad \textbf{D.} 3/10\par\vspace{2pt}
\textbf{\color{primary} Câu 45.} Thử thách tư duy trong phạm vi kiến thức Chương VIII: Gieo đồng thời hai con xúc xắc cân đối và đồng chất. Xác suất của biến cố 'Tổng số chấm xuất hiện trên hai con xúc xắc lớn hơn hoặc bằng 10' bằng. Đáp số chính xác là:\par
\hspace*{1.5em}\textbf{A.} 1/4\quad\quad \textbf{B.} 1/6\quad\quad \textbf{C.} 7/36\quad\quad \textbf{D.} 5/36\par\vspace{2pt}
\textbf{\color{primary} Câu 46.} Sử dụng phương pháp giải toán của Chương VIII để tìm đáp án cho bài toán: Từ hai túi: túi I chứa các thẻ ghi số 1; 2; 3 và túi II chứa 8 thẻ ghi các số từ 1 đến 8. Rút ngẫu nhiên mỗi túi 1 thẻ, xác suất để cả hai thẻ rút ra đều ghi số 1 là:\par
\hspace*{1.5em}\textbf{A.} 2/24\quad\quad \textbf{B.} 1/24\quad\quad \textbf{C.} 1/11\quad\quad \textbf{D.} 1/8\par\vspace{2pt}
\textbf{\color{primary} Câu 47.} Xét bài toán vận dụng thực tiễn / nâng cao thuộc Chương VIII: Từ một nhóm gồm 3 học sinh nam và 2 học sinh nữ, chọn ngẫu nhiên đồng thời 2 học sinh. Xác suất để chọn được đúng 1 học sinh nam và 1 học sinh nữ bằng:\par
\hspace*{1.5em}\textbf{A.} 3/5\quad\quad \textbf{B.} 1/2\quad\quad \textbf{C.} 3/10\quad\quad \textbf{D.} 2/5\par\vspace{2pt}
\textbf{\color{primary} Câu 48.} Tìm kết quả cuối cùng của bài toán phân loại học sinh khá giỏi (Chương VIII): Gieo đồng thời hai con xúc xắc cân đối và đồng chất. Xác suất của biến cố 'Tổng số chấm xuất hiện trên hai con xúc xắc lớn hơn hoặc bằng 10' bằng:\par
\hspace*{1.5em}\textbf{A.} 1/4\quad\quad \textbf{B.} 1/6\quad\quad \textbf{C.} 7/36\quad\quad \textbf{D.} 5/36\par\vspace{2pt}
\textbf{\color{primary} Câu 49.} Kết hợp các tính chất đã học trong Chương VIII, hãy giải quyết bài toán: Từ hai túi: túi I chứa các thẻ ghi số 1; 2; 3 và túi II chứa 11 thẻ ghi các số từ 1 đến 11. Rút ngẫu nhiên mỗi túi 1 thẻ, xác suất để cả hai thẻ rút ra đều ghi số 1 là:\par
\hspace*{1.5em}\textbf{A.} 1/14\quad\quad \textbf{B.} 2/33\quad\quad \textbf{C.} 1/33\quad\quad \textbf{D.} 1/11\par\vspace{2pt}
\textbf{\color{primary} Câu 50.} Câu hỏi vận dụng cao cuối Chương VIII: Từ một nhóm gồm 3 học sinh nam và 2 học sinh nữ, chọn ngẫu nhiên đồng thời 2 học sinh. Xác suất để chọn được đúng 1 học sinh nam và 1 học sinh nữ bằng. Phương án đúng là:\par
\hspace*{1.5em}\textbf{A.} 3/5\quad\quad \textbf{B.} 2/5\quad\quad \textbf{C.} 3/10\quad\quad \textbf{D.} 1/2\par\vspace{2pt}
\subsection*{\color{primary} Bảng đáp án Chương VIII: Xác suất của biến cố trong một số mô hình xác suất đơn giản}
\begin{center}
\renewcommand{\arraystretch}{1.25}
\begin{tabularx}{\textwidth}{*{10}{>{\centering\arraybackslash}X}}
\toprule
\textbf{\color{primary} Cột 1} & \textbf{\color{primary} Cột 2} & \textbf{\color{primary} Cột 3} & \textbf{\color{primary} Cột 4} & \textbf{\color{primary} Cột 5} & \textbf{\color{primary} Cột 6} & \textbf{\color{primary} Cột 7} & \textbf{\color{primary} Cột 8} & \textbf{\color{primary} Cột 9} & \textbf{\color{primary} Cột 10} \\
\midrule
\textbf{1.C} & \textbf{2.C} & \textbf{3.B} & \textbf{4.D} & \textbf{5.C} & \textbf{6.B} & \textbf{7.A} & \textbf{8.B} & \textbf{9.A} & \textbf{10.A} \\
\rowcolor{rowtint} \textbf{11.D} & \textbf{12.C} & \textbf{13.C} & \textbf{14.A} & \textbf{15.C} & \textbf{16.D} & \textbf{17.C} & \textbf{18.C} & \textbf{19.B} & \textbf{20.B} \\
\textbf{21.D} & \textbf{22.C} & \textbf{23.A} & \textbf{24.D} & \textbf{25.D} & \textbf{26.B} & \textbf{27.B} & \textbf{28.D} & \textbf{29.B} & \textbf{30.C} \\
\rowcolor{rowtint} \textbf{31.C} & \textbf{32.D} & \textbf{33.B} & \textbf{34.C} & \textbf{35.B} & \textbf{36.B} & \textbf{37.D} & \textbf{38.A} & \textbf{39.B} & \textbf{40.B} \\
\textbf{41.B} & \textbf{42.B} & \textbf{43.D} & \textbf{44.C} & \textbf{45.B} & \textbf{46.B} & \textbf{47.A} & \textbf{48.B} & \textbf{49.C} & \textbf{50.A} \\
\bottomrule
\end{tabularx}
\end{center}\vspace{10pt}
\section*{\color{primary} CHƯƠNG IX: ĐƯỜNG TRÒN NGOẠI TIẾP VÀ ĐƯỜNG TRÒN NỘI TIẾP}
\textit{Trọng tâm kiến thức SGK: Đường tròn ngoại tiếp và đường tròn nội tiếp một tam giác, tứ giác nội tiếp đường tròn, đa giác đều và phép quay.}\par\vspace{4pt}
\subsection*{\color{accent} Phần I. Mức độ Dễ (Nhận biết -- Câu 1 đến Câu 20)}
\textbf{\color{primary} Câu 1.} Cho tam giác ABC vuông tại A có cạnh huyền BC = 8 cm. Bán kính R của đường tròn ngoại tiếp tam giác ABC bằng:\par
\hspace*{1.5em}\textbf{A.} 4$\sqrt{2}$ cm\quad\quad \textbf{B.} 4 cm\quad\quad \textbf{C.} 8 cm\quad\quad \textbf{D.} 4/2 cm\par\vspace{2pt}
\textbf{\color{primary} Câu 2.} Thực hiện yêu cầu nhận biết sau: Cho tứ giác ABCD nội tiếp một đường tròn. Biết góc A = 77°, số đo góc đối diện C bằng.\par
\hspace*{1.5em}\textbf{A.} 103°\quad\quad \textbf{B.} 90°\quad\quad \textbf{C.} 13°\quad\quad \textbf{D.} 77°\par\vspace{2pt}
\textbf{\color{primary} Câu 3.} Trong bài học thuộc Chương IX, xét nội dung: Tâm của đường tròn ngoại tiếp một tam giác là giao điểm của ba đường nào trong tam giác đó. Đáp án đúng là:\par
\hspace*{1.5em}\textbf{A.} Ba đường trung tuyến\quad\quad \textbf{B.} Ba đường trung trực\quad\quad \textbf{C.} Ba đường cao\quad\quad \textbf{D.} Ba đường phân giác trong\par\vspace{2pt}
\textbf{\color{primary} Câu 4.} Dựa vào kiến thức chuẩn SGK của Chương IX (Đường tròn ngoại tiếp và đường tròn nội tiếp), cho tam giác ABC vuông tại A có cạnh huyền BC = 14 cm. Bán kính R của đường tròn ngoại tiếp tam giác ABC bằng:\par
\hspace*{1.5em}\textbf{A.} 7/2 cm\quad\quad \textbf{B.} 7$\sqrt{2}$ cm\quad\quad \textbf{C.} 14 cm\quad\quad \textbf{D.} 7 cm\par\vspace{2pt}
\textbf{\color{primary} Câu 5.} Chọn kết quả đúng nhất khi xét: Cho tứ giác ABCD nội tiếp một đường tròn. Biết góc A = 83°, số đo góc đối diện C bằng:\par
\hspace*{1.5em}\textbf{A.} 83°\quad\quad \textbf{B.} 97°\quad\quad \textbf{C.} 90°\quad\quad \textbf{D.} 7°\par\vspace{2pt}
\textbf{\color{primary} Câu 6.} Khi làm bài tập cơ bản về đường tròn ngoại tiếp và đường tròn nội tiếp, một học sinh gặp câu hỏi: Tâm của đường tròn ngoại tiếp một tam giác là giao điểm của ba đường nào trong tam giác đó. Kết quả nào dưới đây đúng?\par
\hspace*{1.5em}\textbf{A.} Ba đường trung trực\quad\quad \textbf{B.} Ba đường phân giác trong\quad\quad \textbf{C.} Ba đường cao\quad\quad \textbf{D.} Ba đường trung tuyến\par\vspace{2pt}
\textbf{\color{primary} Câu 7.} Xác định giá trị hoặc khẳng định đúng cho bài toán sau: Cho tam giác ABC vuông tại A có cạnh huyền BC = 20 cm. Bán kính R của đường tròn ngoại tiếp tam giác ABC bằng:\par
\hspace*{1.5em}\textbf{A.} 10 cm\quad\quad \textbf{B.} 20 cm\quad\quad \textbf{C.} 10$\sqrt{2}$ cm\quad\quad \textbf{D.} 10/2 cm\par\vspace{2pt}
\textbf{\color{primary} Câu 8.} Áp dụng trực tiếp định nghĩa và quy tắc trong Chương IX: Cho tứ giác ABCD nội tiếp một đường tròn. Biết góc A = 89°, số đo góc đối diện C bằng:\par
\hspace*{1.5em}\textbf{A.} 89°\quad\quad \textbf{B.} 91°\quad\quad \textbf{C.} 1°\quad\quad \textbf{D.} 90°\par\vspace{2pt}
\textbf{\color{primary} Câu 9.} Kiểm tra nhanh kiến thức nền tảng: Tâm của đường tròn ngoại tiếp một tam giác là giao điểm của ba đường nào trong tam giác đó?\par
\hspace*{1.5em}\textbf{A.} Ba đường trung tuyến\quad\quad \textbf{B.} Ba đường cao\quad\quad \textbf{C.} Ba đường trung trực\quad\quad \textbf{D.} Ba đường phân giác trong\par\vspace{2pt}
\textbf{\color{primary} Câu 10.} Cho dữ kiện nhận biết sau đây: Cho tam giác ABC vuông tại A có cạnh huyền BC = 26 cm. Bán kính R của đường tròn ngoại tiếp tam giác ABC bằng. Phương án chính xác là:\par
\hspace*{1.5em}\textbf{A.} 13$\sqrt{2}$ cm\quad\quad \textbf{B.} 26 cm\quad\quad \textbf{C.} 13/2 cm\quad\quad \textbf{D.} 13 cm\par\vspace{2pt}
\textbf{\color{primary} Câu 11.} Hãy chỉ ra kết quả đúng trong bốn lựa chọn A, B, C, D đối với yêu cầu: Cho tứ giác ABCD nội tiếp một đường tròn. Biết góc A = 95°, số đo góc đối diện C bằng:\par
\hspace*{1.5em}\textbf{A.} 85°\quad\quad \textbf{B.} 95°\quad\quad \textbf{C.} -5°\quad\quad \textbf{D.} 90°\par\vspace{2pt}
\textbf{\color{primary} Câu 12.} Câu hỏi ôn tập lý thuyết và kỹ năng cơ bản (Chương IX): Tâm của đường tròn ngoại tiếp một tam giác là giao điểm của ba đường nào trong tam giác đó:\par
\hspace*{1.5em}\textbf{A.} Ba đường cao\quad\quad \textbf{B.} Ba đường trung tuyến\quad\quad \textbf{C.} Ba đường phân giác trong\quad\quad \textbf{D.} Ba đường trung trực\par\vspace{2pt}
\textbf{\color{primary} Câu 13.} Theo nội dung sách giáo khoa Kết nối tri thức với cuộc sống (Chương IX), cho tam giác ABC vuông tại A có cạnh huyền BC = 32 cm. Bán kính R của đường tròn ngoại tiếp tam giác ABC bằng:\par
\hspace*{1.5em}\textbf{A.} 16$\sqrt{2}$ cm\quad\quad \textbf{B.} 32 cm\quad\quad \textbf{C.} 16/2 cm\quad\quad \textbf{D.} 16 cm\par\vspace{2pt}
\textbf{\color{primary} Câu 14.} Tìm đáp số đúng của bài toán cơ bản sau: Cho tứ giác ABCD nội tiếp một đường tròn. Biết góc A = 101°, số đo góc đối diện C bằng:\par
\hspace*{1.5em}\textbf{A.} -11°\quad\quad \textbf{B.} 101°\quad\quad \textbf{C.} 79°\quad\quad \textbf{D.} 90°\par\vspace{2pt}
\textbf{\color{primary} Câu 15.} Xét đối tượng toán học trong Chương IX: Tâm của đường tròn ngoại tiếp một tam giác là giao điểm của ba đường nào trong tam giác đó. Giá trị hoặc tính chất tương ứng là:\par
\hspace*{1.5em}\textbf{A.} Ba đường trung tuyến\quad\quad \textbf{B.} Ba đường phân giác trong\quad\quad \textbf{C.} Ba đường trung trực\quad\quad \textbf{D.} Ba đường cao\par\vspace{2pt}
\textbf{\color{primary} Câu 16.} Hoàn thành phép xác định cơ bản sau đây: Cho tam giác ABC vuông tại A có cạnh huyền BC = 38 cm. Bán kính R của đường tròn ngoại tiếp tam giác ABC bằng:\par
\hspace*{1.5em}\textbf{A.} 19/2 cm\quad\quad \textbf{B.} 19 cm\quad\quad \textbf{C.} 19$\sqrt{2}$ cm\quad\quad \textbf{D.} 38 cm\par\vspace{2pt}
\textbf{\color{primary} Câu 17.} Trong các khẳng định hoặc giá trị dưới đây, đâu là đáp án đúng cho câu hỏi: Cho tứ giác ABCD nội tiếp một đường tròn. Biết góc A = 107°, số đo góc đối diện C bằng?\par
\hspace*{1.5em}\textbf{A.} 107°\quad\quad \textbf{B.} -17°\quad\quad \textbf{C.} 90°\quad\quad \textbf{D.} 73°\par\vspace{2pt}
\textbf{\color{primary} Câu 18.} Bài tập khởi động Chương IX: Tâm của đường tròn ngoại tiếp một tam giác là giao điểm của ba đường nào trong tam giác đó:\par
\hspace*{1.5em}\textbf{A.} Ba đường phân giác trong\quad\quad \textbf{B.} Ba đường trung tuyến\quad\quad \textbf{C.} Ba đường cao\quad\quad \textbf{D.} Ba đường trung trực\par\vspace{2pt}
\textbf{\color{primary} Câu 19.} Quan sát dữ kiện đề bài: Cho tam giác ABC vuông tại A có cạnh huyền BC = 44 cm. Bán kính R của đường tròn ngoại tiếp tam giác ABC bằng. Hãy chọn phương án đúng:\par
\hspace*{1.5em}\textbf{A.} 44 cm\quad\quad \textbf{B.} 22 cm\quad\quad \textbf{C.} 22/2 cm\quad\quad \textbf{D.} 22$\sqrt{2}$ cm\par\vspace{2pt}
\textbf{\color{primary} Câu 20.} Sử dụng công thức hoặc tính chất đầu bài của Chương IX để giải quyết yêu cầu: Cho tứ giác ABCD nội tiếp một đường tròn. Biết góc A = 113°, số đo góc đối diện C bằng:\par
\hspace*{1.5em}\textbf{A.} 90°\quad\quad \textbf{B.} 67°\quad\quad \textbf{C.} -23°\quad\quad \textbf{D.} 113°\par\vspace{2pt}
\subsection*{\color{accent} Phần II. Mức độ Trung bình (Thông hiểu -- Câu 21 đến Câu 40)}
\textbf{\color{primary} Câu 21.} Cho tam giác đều ABC có độ dài cạnh a = 6 cm. Bán kính R của đường tròn ngoại tiếp tam giác đều ABC bằng:\par
\hspace*{1.5em}\textbf{A.} 6$\sqrt{3}$ cm\quad\quad \textbf{B.} 2$\sqrt{3}$ cm\quad\quad \textbf{C.} 2 cm\quad\quad \textbf{D.} 2$\sqrt{3}$/2 cm\par\vspace{2pt}
\textbf{\color{primary} Câu 22.} Vận dụng quy tắc trong Chương IX (Đường tròn ngoại tiếp và đường tròn nội tiếp) để giải bài toán: Một hình vuông ABCD nội tiếp đường tròn có bán kính R = 4$\sqrt{2}$ cm. Độ dài cạnh của hình vuông ABCD bằng:\par
\hspace*{1.5em}\textbf{A.} 8 cm\quad\quad \textbf{B.} 4$\sqrt{2}$ cm\quad\quad \textbf{C.} 4 cm\quad\quad \textbf{D.} 16 cm\par\vspace{2pt}
\textbf{\color{primary} Câu 23.} Biến đổi và tính toán theo kiến thức Chương IX: Trong các tứ giác sau đây, tứ giác nào LUÔN nội tiếp được một đường tròn. Kết quả thu được là:\par
\begin{itemize}\setlength{\itemsep}{1pt}\setlength{\parskip}{0pt}
  \item[] \textbf{A.} Hình bình hành bất kỳ
  \item[] \textbf{B.} Hình thang vuông bất kỳ
  \item[] \textbf{C.} Hình thoi có một góc bằng 60°
  \item[] \textbf{D.} Hình thang cân, hình chữ nhật, hình vuông
\end{itemize}\vspace{2pt}
\textbf{\color{primary} Câu 24.} Xét bài toán mức độ thông hiểu sau: Cho tam giác đều ABC có độ dài cạnh a = 15 cm. Bán kính R của đường tròn ngoại tiếp tam giác đều ABC bằng. Đáp án đúng là:\par
\hspace*{1.5em}\textbf{A.} 5$\sqrt{3}$/2 cm\quad\quad \textbf{B.} 15$\sqrt{3}$ cm\quad\quad \textbf{C.} 5 cm\quad\quad \textbf{D.} 5$\sqrt{3}$ cm\par\vspace{2pt}
\textbf{\color{primary} Câu 25.} Để hoàn thành phiếu học tập Chương IX, hãy xác định: Một hình vuông ABCD nội tiếp đường tròn có bán kính R = 7$\sqrt{2}$ cm. Độ dài cạnh của hình vuông ABCD bằng:\par
\hspace*{1.5em}\textbf{A.} 7 cm\quad\quad \textbf{B.} 7$\sqrt{2}$ cm\quad\quad \textbf{C.} 28 cm\quad\quad \textbf{D.} 14 cm\par\vspace{2pt}
\textbf{\color{primary} Câu 26.} Phân tích dữ kiện bài toán: Trong các tứ giác sau đây, tứ giác nào LUÔN nội tiếp được một đường tròn. Giá trị cần tìm bằng:\par
\begin{itemize}\setlength{\itemsep}{1pt}\setlength{\parskip}{0pt}
  \item[] \textbf{A.} Hình thoi có một góc bằng 60°
  \item[] \textbf{B.} Hình thang cân, hình chữ nhật, hình vuông
  \item[] \textbf{C.} Hình bình hành bất kỳ
  \item[] \textbf{D.} Hình thang vuông bất kỳ
\end{itemize}\vspace{2pt}
\textbf{\color{primary} Câu 27.} Áp dụng tính chất đã học ở Chương IX vào bài tập sau: Cho tam giác đều ABC có độ dài cạnh a = 24 cm. Bán kính R của đường tròn ngoại tiếp tam giác đều ABC bằng:\par
\hspace*{1.5em}\textbf{A.} 8$\sqrt{3}$/2 cm\quad\quad \textbf{B.} 8 cm\quad\quad \textbf{C.} 24$\sqrt{3}$ cm\quad\quad \textbf{D.} 8$\sqrt{3}$ cm\par\vspace{2pt}
\textbf{\color{primary} Câu 28.} Thực hiện các bước suy luận hợp lý để giải quyết yêu cầu: Một hình vuông ABCD nội tiếp đường tròn có bán kính R = 10$\sqrt{2}$ cm. Độ dài cạnh của hình vuông ABCD bằng:\par
\hspace*{1.5em}\textbf{A.} 10$\sqrt{2}$ cm\quad\quad \textbf{B.} 10 cm\quad\quad \textbf{C.} 20 cm\quad\quad \textbf{D.} 40 cm\par\vspace{2pt}
\textbf{\color{primary} Câu 29.} Cho bài toán tự luyện (Đường tròn ngoại tiếp và đường tròn nội tiếp): Trong các tứ giác sau đây, tứ giác nào LUÔN nội tiếp được một đường tròn. Phương án nào sau đây chính xác?\par
\begin{itemize}\setlength{\itemsep}{1pt}\setlength{\parskip}{0pt}
  \item[] \textbf{A.} Hình thang cân, hình chữ nhật, hình vuông
  \item[] \textbf{B.} Hình thang vuông bất kỳ
  \item[] \textbf{C.} Hình bình hành bất kỳ
  \item[] \textbf{D.} Hình thoi có một góc bằng 60°
\end{itemize}\vspace{2pt}
\textbf{\color{primary} Câu 30.} Kết hợp các phép biến đổi trong Chương IX, hãy tìm kết quả của bài toán: Cho tam giác đều ABC có độ dài cạnh a = 33 cm. Bán kính R của đường tròn ngoại tiếp tam giác đều ABC bằng:\par
\hspace*{1.5em}\textbf{A.} 11 cm\quad\quad \textbf{B.} 33$\sqrt{3}$ cm\quad\quad \textbf{C.} 11$\sqrt{3}$/2 cm\quad\quad \textbf{D.} 11$\sqrt{3}$ cm\par\vspace{2pt}
\textbf{\color{primary} Câu 31.} Một bài kiểm tra định kỳ nội dung Chương IX nêu câu hỏi: Một hình vuông ABCD nội tiếp đường tròn có bán kính R = 13$\sqrt{2}$ cm. Độ dài cạnh của hình vuông ABCD bằng. Hãy chọn đáp số đúng:\par
\hspace*{1.5em}\textbf{A.} 52 cm\quad\quad \textbf{B.} 13$\sqrt{2}$ cm\quad\quad \textbf{C.} 26 cm\quad\quad \textbf{D.} 13 cm\par\vspace{2pt}
\textbf{\color{primary} Câu 32.} Tìm lời giải đúng cho bài toán thông hiểu: Trong các tứ giác sau đây, tứ giác nào LUÔN nội tiếp được một đường tròn:\par
\begin{itemize}\setlength{\itemsep}{1pt}\setlength{\parskip}{0pt}
  \item[] \textbf{A.} Hình thang cân, hình chữ nhật, hình vuông
  \item[] \textbf{B.} Hình thang vuông bất kỳ
  \item[] \textbf{C.} Hình bình hành bất kỳ
  \item[] \textbf{D.} Hình thoi có một góc bằng 60°
\end{itemize}\vspace{2pt}
\textbf{\color{primary} Câu 33.} Sử dụng mối liên hệ giữa các đại lượng trong Chương IX: Cho tam giác đều ABC có độ dài cạnh a = 42 cm. Bán kính R của đường tròn ngoại tiếp tam giác đều ABC bằng:\par
\hspace*{1.5em}\textbf{A.} 14$\sqrt{3}$ cm\quad\quad \textbf{B.} 42$\sqrt{3}$ cm\quad\quad \textbf{C.} 14 cm\quad\quad \textbf{D.} 14$\sqrt{3}$/2 cm\par\vspace{2pt}
\textbf{\color{primary} Câu 34.} Hãy tính toán và chọn phương án đúng đối với bài tập: Một hình vuông ABCD nội tiếp đường tròn có bán kính R = 16$\sqrt{2}$ cm. Độ dài cạnh của hình vuông ABCD bằng:\par
\hspace*{1.5em}\textbf{A.} 16$\sqrt{2}$ cm\quad\quad \textbf{B.} 16 cm\quad\quad \textbf{C.} 64 cm\quad\quad \textbf{D.} 32 cm\par\vspace{2pt}
\textbf{\color{primary} Câu 35.} Dựa trên các tính chất trọng tâm của Chương IX, xét bài toán: Trong các tứ giác sau đây, tứ giác nào LUÔN nội tiếp được một đường tròn:\par
\begin{itemize}\setlength{\itemsep}{1pt}\setlength{\parskip}{0pt}
  \item[] \textbf{A.} Hình thang cân, hình chữ nhật, hình vuông
  \item[] \textbf{B.} Hình bình hành bất kỳ
  \item[] \textbf{C.} Hình thoi có một góc bằng 60°
  \item[] \textbf{D.} Hình thang vuông bất kỳ
\end{itemize}\vspace{2pt}
\textbf{\color{primary} Câu 36.} Thực hành giải toán (Đường tròn ngoại tiếp và đường tròn nội tiếp): Cho tam giác đều ABC có độ dài cạnh a = 51 cm. Bán kính R của đường tròn ngoại tiếp tam giác đều ABC bằng. Đáp án nào dưới đây thỏa mãn?\par
\hspace*{1.5em}\textbf{A.} 17$\sqrt{3}$/2 cm\quad\quad \textbf{B.} 51$\sqrt{3}$ cm\quad\quad \textbf{C.} 17$\sqrt{3}$ cm\quad\quad \textbf{D.} 17 cm\par\vspace{2pt}
\textbf{\color{primary} Câu 37.} Kiểm tra kỹ năng biến đổi và suy luận qua bài toán: Một hình vuông ABCD nội tiếp đường tròn có bán kính R = 19$\sqrt{2}$ cm. Độ dài cạnh của hình vuông ABCD bằng:\par
\hspace*{1.5em}\textbf{A.} 76 cm\quad\quad \textbf{B.} 19 cm\quad\quad \textbf{C.} 38 cm\quad\quad \textbf{D.} 19$\sqrt{2}$ cm\par\vspace{2pt}
\textbf{\color{primary} Câu 38.} Cho giả thiết thuộc phạm vi Chương IX: Trong các tứ giác sau đây, tứ giác nào LUÔN nội tiếp được một đường tròn. Kết luận đúng là:\par
\begin{itemize}\setlength{\itemsep}{1pt}\setlength{\parskip}{0pt}
  \item[] \textbf{A.} Hình bình hành bất kỳ
  \item[] \textbf{B.} Hình thang vuông bất kỳ
  \item[] \textbf{C.} Hình thang cân, hình chữ nhật, hình vuông
  \item[] \textbf{D.} Hình thoi có một góc bằng 60°
\end{itemize}\vspace{2pt}
\textbf{\color{primary} Câu 39.} Rút gọn hoặc tính giá trị theo yêu cầu sau của Chương IX: Cho tam giác đều ABC có độ dài cạnh a = 60 cm. Bán kính R của đường tròn ngoại tiếp tam giác đều ABC bằng:\par
\hspace*{1.5em}\textbf{A.} 20$\sqrt{3}$/2 cm\quad\quad \textbf{B.} 20 cm\quad\quad \textbf{C.} 20$\sqrt{3}$ cm\quad\quad \textbf{D.} 60$\sqrt{3}$ cm\par\vspace{2pt}
\textbf{\color{primary} Câu 40.} Xác định phương án đúng sau khi thực hiện phép tính / suy luận: Một hình vuông ABCD nội tiếp đường tròn có bán kính R = 22$\sqrt{2}$ cm. Độ dài cạnh của hình vuông ABCD bằng:\par
\hspace*{1.5em}\textbf{A.} 22$\sqrt{2}$ cm\quad\quad \textbf{B.} 44 cm\quad\quad \textbf{C.} 88 cm\quad\quad \textbf{D.} 22 cm\par\vspace{2pt}
\subsection*{\color{accent} Phần III. Mức độ Khó (Vận dụng \& Vận dụng cao -- Câu 41 đến Câu 50)}
\textbf{\color{primary} Câu 41.} Cho lục giác đều ABCDEF tâm O. Phép quay thuận chiều tâm O góc quay $\alpha$ nào dưới đây biến lục giác đều ABCDEF thành chính nó?\par
\hspace*{1.5em}\textbf{A.} 60°\quad\quad \textbf{B.} 45°\quad\quad \textbf{C.} 150°\quad\quad \textbf{D.} 90°\par\vspace{2pt}
\textbf{\color{primary} Câu 42.} Bài toán vận dụng (Chương IX – Đường tròn ngoại tiếp và đường tròn nội tiếp): Cho tam giác ABC vuông tại A có hai cạnh góc vuông AB = 6 cm, AC = 8 cm. Bán kính r của đường tròn nội tiếp tam giác ABC (r = (AB + AC - BC)/2) bằng. Kết quả đúng là:\par
\hspace*{1.5em}\textbf{A.} 2 cm\quad\quad \textbf{B.} 6 cm\quad\quad \textbf{C.} 5.0 cm\quad\quad \textbf{D.} 4 cm\par\vspace{2pt}
\textbf{\color{primary} Câu 43.} Vận dụng kiến thức nâng cao trong phạm vi Chương IX để giải bài toán: Cho tam giác ABC vuông tại A có hai cạnh góc vuông AB = 12 cm và AC = 16 cm (cạnh huyền BC = 20 cm). Tổng bán kính đường tròn ngoại tiếp R và bán kính đường tròn nội tiếp r của tam giác ABC bằng:\par
\hspace*{1.5em}\textbf{A.} 18.0 cm\quad\quad \textbf{B.} 12.0 cm\quad\quad \textbf{C.} 10.0 cm\quad\quad \textbf{D.} 14.0 cm\par\vspace{2pt}
\textbf{\color{primary} Câu 44.} Phân tích và giải bài toán tổng hợp sau của Chương IX: Cho lục giác đều ABCDEF tâm O. Phép quay thuận chiều tâm O góc quay $\alpha$ nào dưới đây biến lục giác đều ABCDEF thành chính nó:\par
\hspace*{1.5em}\textbf{A.} 240°\quad\quad \textbf{B.} 90°\quad\quad \textbf{C.} 45°\quad\quad \textbf{D.} 150°\par\vspace{2pt}
\textbf{\color{primary} Câu 45.} Thử thách tư duy trong phạm vi kiến thức Chương IX: Cho tam giác ABC vuông tại A có hai cạnh góc vuông AB = 3 cm, AC = 4 cm. Bán kính r của đường tròn nội tiếp tam giác ABC (r = (AB + AC - BC)/2) bằng. Đáp số chính xác là:\par
\hspace*{1.5em}\textbf{A.} 2.5 cm\quad\quad \textbf{B.} 1 cm\quad\quad \textbf{C.} 3 cm\quad\quad \textbf{D.} 2 cm\par\vspace{2pt}
\textbf{\color{primary} Câu 46.} Sử dụng phương pháp giải toán của Chương IX để tìm đáp án cho bài toán: Cho tam giác ABC vuông tại A có hai cạnh góc vuông AB = 21 cm và AC = 28 cm (cạnh huyền BC = 35 cm). Tổng bán kính đường tròn ngoại tiếp R và bán kính đường tròn nội tiếp r của tam giác ABC bằng:\par
\hspace*{1.5em}\textbf{A.} 17.5 cm\quad\quad \textbf{B.} 31.5 cm\quad\quad \textbf{C.} 21.0 cm\quad\quad \textbf{D.} 24.5 cm\par\vspace{2pt}
\textbf{\color{primary} Câu 47.} Xét bài toán vận dụng thực tiễn / nâng cao thuộc Chương IX: Cho lục giác đều ABCDEF tâm O. Phép quay thuận chiều tâm O góc quay $\alpha$ nào dưới đây biến lục giác đều ABCDEF thành chính nó:\par
\hspace*{1.5em}\textbf{A.} 45°\quad\quad \textbf{B.} 150°\quad\quad \textbf{C.} 120°\quad\quad \textbf{D.} 90°\par\vspace{2pt}
\textbf{\color{primary} Câu 48.} Tìm kết quả cuối cùng của bài toán phân loại học sinh khá giỏi (Chương IX): Cho tam giác ABC vuông tại A có hai cạnh góc vuông AB = 6 cm, AC = 8 cm. Bán kính r của đường tròn nội tiếp tam giác ABC (r = (AB + AC - BC)/2) bằng:\par
\hspace*{1.5em}\textbf{A.} 2 cm\quad\quad \textbf{B.} 4 cm\quad\quad \textbf{C.} 6 cm\quad\quad \textbf{D.} 5.0 cm\par\vspace{2pt}
\textbf{\color{primary} Câu 49.} Kết hợp các tính chất đã học trong Chương IX, hãy giải quyết bài toán: Cho tam giác ABC vuông tại A có hai cạnh góc vuông AB = 30 cm và AC = 40 cm (cạnh huyền BC = 50 cm). Tổng bán kính đường tròn ngoại tiếp R và bán kính đường tròn nội tiếp r của tam giác ABC bằng:\par
\hspace*{1.5em}\textbf{A.} 45.0 cm\quad\quad \textbf{B.} 35.0 cm\quad\quad \textbf{C.} 30.0 cm\quad\quad \textbf{D.} 25.0 cm\par\vspace{2pt}
\textbf{\color{primary} Câu 50.} Câu hỏi vận dụng cao cuối Chương IX: Cho lục giác đều ABCDEF tâm O. Phép quay thuận chiều tâm O góc quay $\alpha$ nào dưới đây biến lục giác đều ABCDEF thành chính nó. Phương án đúng là:\par
\hspace*{1.5em}\textbf{A.} 45°\quad\quad \textbf{B.} 150°\quad\quad \textbf{C.} 90°\quad\quad \textbf{D.} 300°\par\vspace{2pt}
\subsection*{\color{primary} Bảng đáp án Chương IX: Đường tròn ngoại tiếp và đường tròn nội tiếp}
\begin{center}
\renewcommand{\arraystretch}{1.25}
\begin{tabularx}{\textwidth}{*{10}{>{\centering\arraybackslash}X}}
\toprule
\textbf{\color{primary} Cột 1} & \textbf{\color{primary} Cột 2} & \textbf{\color{primary} Cột 3} & \textbf{\color{primary} Cột 4} & \textbf{\color{primary} Cột 5} & \textbf{\color{primary} Cột 6} & \textbf{\color{primary} Cột 7} & \textbf{\color{primary} Cột 8} & \textbf{\color{primary} Cột 9} & \textbf{\color{primary} Cột 10} \\
\midrule
\textbf{1.B} & \textbf{2.A} & \textbf{3.B} & \textbf{4.D} & \textbf{5.B} & \textbf{6.A} & \textbf{7.A} & \textbf{8.B} & \textbf{9.C} & \textbf{10.D} \\
\rowcolor{rowtint} \textbf{11.A} & \textbf{12.D} & \textbf{13.D} & \textbf{14.C} & \textbf{15.C} & \textbf{16.B} & \textbf{17.D} & \textbf{18.D} & \textbf{19.B} & \textbf{20.B} \\
\textbf{21.B} & \textbf{22.A} & \textbf{23.D} & \textbf{24.D} & \textbf{25.D} & \textbf{26.B} & \textbf{27.D} & \textbf{28.C} & \textbf{29.A} & \textbf{30.D} \\
\rowcolor{rowtint} \textbf{31.C} & \textbf{32.A} & \textbf{33.A} & \textbf{34.D} & \textbf{35.A} & \textbf{36.C} & \textbf{37.C} & \textbf{38.C} & \textbf{39.C} & \textbf{40.B} \\
\textbf{41.A} & \textbf{42.A} & \textbf{43.D} & \textbf{44.A} & \textbf{45.B} & \textbf{46.D} & \textbf{47.C} & \textbf{48.A} & \textbf{49.B} & \textbf{50.D} \\
\bottomrule
\end{tabularx}
\end{center}\vspace{10pt}
\section*{\color{primary} CHƯƠNG X: MỘT SỐ HÌNH KHỐI TRONG THỰC TIỄN}
\textit{Trọng tâm kiến thức SGK: Hình trụ, hình nón và mặt cầu/hình cầu (công thức tính diện tích xung quanh, diện tích toàn phần và thể tích).}\par\vspace{4pt}
\subsection*{\color{accent} Phần I. Mức độ Dễ (Nhận biết -- Câu 1 đến Câu 20)}
\textbf{\color{primary} Câu 1.} Một hình trụ có bán kính đáy r = 2 cm và chiều cao h = 5 cm. Diện tích xung quanh của hình trụ đó bằng:\par
\hspace*{1.5em}\textbf{A.} 20$\pi$ cm$^2$\quad\quad \textbf{B.} 20$\pi$ cm$^2$ (1)\quad\quad \textbf{C.} 10$\pi$ cm$^2$\quad\quad \textbf{D.} 40$\pi$ cm$^2$\par\vspace{2pt}
\textbf{\color{primary} Câu 2.} Thực hiện yêu cầu nhận biết sau: Một mặt cầu có bán kính R = 3 cm. Diện tích S của mặt cầu đó (S = 4$\pi$R$^2$) bằng.\par
\hspace*{1.5em}\textbf{A.} 18$\pi$ cm$^2$\quad\quad \textbf{B.} 9$\pi$ cm$^2$\quad\quad \textbf{C.} 36.0$\pi$ cm$^2$\quad\quad \textbf{D.} 36$\pi$ cm$^2$\par\vspace{2pt}
\textbf{\color{primary} Câu 3.} Trong bài học thuộc Chương X, xét nội dung: Một hình nón có bán kính đáy r = 3 cm và chiều cao h = 6 cm. Thể tích V của hình nón đó (V = 1/3 $\cdot$ $\pi$r$^2$h) bằng. Đáp án đúng là:\par
\hspace*{1.5em}\textbf{A.} 6$\pi$ cm$^3$\quad\quad \textbf{B.} 18$\pi$ cm$^3$\quad\quad \textbf{C.} 54$\pi$ cm$^3$\quad\quad \textbf{D.} 36$\pi$ cm$^3$\par\vspace{2pt}
\textbf{\color{primary} Câu 4.} Dựa vào kiến thức chuẩn SGK của Chương X (Một số hình khối trong thực tiễn), một hình trụ có bán kính đáy r = 5 cm và chiều cao h = 5 cm. Diện tích xung quanh của hình trụ đó bằng:\par
\hspace*{1.5em}\textbf{A.} 100$\pi$ cm$^2$\quad\quad \textbf{B.} 125$\pi$ cm$^2$\quad\quad \textbf{C.} 50$\pi$ cm$^2$\quad\quad \textbf{D.} 25$\pi$ cm$^2$\par\vspace{2pt}
\textbf{\color{primary} Câu 5.} Chọn kết quả đúng nhất khi xét: Một mặt cầu có bán kính R = 6 cm. Diện tích S của mặt cầu đó (S = 4$\pi$R$^2$) bằng:\par
\hspace*{1.5em}\textbf{A.} 144$\pi$ cm$^2$\quad\quad \textbf{B.} 72$\pi$ cm$^2$\quad\quad \textbf{C.} 36$\pi$ cm$^2$\quad\quad \textbf{D.} 288.0$\pi$ cm$^2$\par\vspace{2pt}
\textbf{\color{primary} Câu 6.} Khi làm bài tập cơ bản về một số hình khối trong thực tiễn, một học sinh gặp câu hỏi: Một hình nón có bán kính đáy r = 3 cm và chiều cao h = 9 cm. Thể tích V của hình nón đó (V = 1/3 $\cdot$ $\pi$r$^2$h) bằng. Kết quả nào dưới đây đúng?\par
\hspace*{1.5em}\textbf{A.} 54$\pi$ cm$^3$\quad\quad \textbf{B.} 81$\pi$ cm$^3$\quad\quad \textbf{C.} 27$\pi$ cm$^3$\quad\quad \textbf{D.} 9$\pi$ cm$^3$\par\vspace{2pt}
\textbf{\color{primary} Câu 7.} Xác định giá trị hoặc khẳng định đúng cho bài toán sau: Một hình trụ có bán kính đáy r = 8 cm và chiều cao h = 5 cm. Diện tích xung quanh của hình trụ đó bằng:\par
\hspace*{1.5em}\textbf{A.} 80$\pi$ cm$^2$\quad\quad \textbf{B.} 320$\pi$ cm$^2$\quad\quad \textbf{C.} 40$\pi$ cm$^2$\quad\quad \textbf{D.} 160$\pi$ cm$^2$\par\vspace{2pt}
\textbf{\color{primary} Câu 8.} Áp dụng trực tiếp định nghĩa và quy tắc trong Chương X: Một mặt cầu có bán kính R = 9 cm. Diện tích S của mặt cầu đó (S = 4$\pi$R$^2$) bằng:\par
\hspace*{1.5em}\textbf{A.} 972.0$\pi$ cm$^2$\quad\quad \textbf{B.} 324$\pi$ cm$^2$\quad\quad \textbf{C.} 162$\pi$ cm$^2$\quad\quad \textbf{D.} 81$\pi$ cm$^2$\par\vspace{2pt}
\textbf{\color{primary} Câu 9.} Kiểm tra nhanh kiến thức nền tảng: Một hình nón có bán kính đáy r = 3 cm và chiều cao h = 12 cm. Thể tích V của hình nón đó (V = 1/3 $\cdot$ $\pi$r$^2$h) bằng:\par
\hspace*{1.5em}\textbf{A.} 72$\pi$ cm$^3$\quad\quad \textbf{B.} 108$\pi$ cm$^3$\quad\quad \textbf{C.} 36$\pi$ cm$^3$\quad\quad \textbf{D.} 12$\pi$ cm$^3$\par\vspace{2pt}
\textbf{\color{primary} Câu 10.} Cho dữ kiện nhận biết sau đây: Một hình trụ có bán kính đáy r = 11 cm và chiều cao h = 5 cm. Diện tích xung quanh của hình trụ đó bằng. Phương án chính xác là:\par
\hspace*{1.5em}\textbf{A.} 110$\pi$ cm$^2$\quad\quad \textbf{B.} 55$\pi$ cm$^2$\quad\quad \textbf{C.} 220$\pi$ cm$^2$\quad\quad \textbf{D.} 605$\pi$ cm$^2$\par\vspace{2pt}
\textbf{\color{primary} Câu 11.} Hãy chỉ ra kết quả đúng trong bốn lựa chọn A, B, C, D đối với yêu cầu: Một mặt cầu có bán kính R = 12 cm. Diện tích S của mặt cầu đó (S = 4$\pi$R$^2$) bằng:\par
\hspace*{1.5em}\textbf{A.} 144$\pi$ cm$^2$\quad\quad \textbf{B.} 576$\pi$ cm$^2$\quad\quad \textbf{C.} 288$\pi$ cm$^2$\quad\quad \textbf{D.} 2304.0$\pi$ cm$^2$\par\vspace{2pt}
\textbf{\color{primary} Câu 12.} Câu hỏi ôn tập lý thuyết và kỹ năng cơ bản (Chương X): Một hình nón có bán kính đáy r = 3 cm và chiều cao h = 15 cm. Thể tích V của hình nón đó (V = 1/3 $\cdot$ $\pi$r$^2$h) bằng:\par
\hspace*{1.5em}\textbf{A.} 15$\pi$ cm$^3$\quad\quad \textbf{B.} 45$\pi$ cm$^3$\quad\quad \textbf{C.} 135$\pi$ cm$^3$\quad\quad \textbf{D.} 90$\pi$ cm$^3$\par\vspace{2pt}
\textbf{\color{primary} Câu 13.} Theo nội dung sách giáo khoa Kết nối tri thức với cuộc sống (Chương X), một hình trụ có bán kính đáy r = 14 cm và chiều cao h = 5 cm. Diện tích xung quanh của hình trụ đó bằng:\par
\hspace*{1.5em}\textbf{A.} 980$\pi$ cm$^2$\quad\quad \textbf{B.} 280$\pi$ cm$^2$\quad\quad \textbf{C.} 140$\pi$ cm$^2$\quad\quad \textbf{D.} 70$\pi$ cm$^2$\par\vspace{2pt}
\textbf{\color{primary} Câu 14.} Tìm đáp số đúng của bài toán cơ bản sau: Một mặt cầu có bán kính R = 15 cm. Diện tích S của mặt cầu đó (S = 4$\pi$R$^2$) bằng:\par
\hspace*{1.5em}\textbf{A.} 450$\pi$ cm$^2$\quad\quad \textbf{B.} 225$\pi$ cm$^2$\quad\quad \textbf{C.} 900$\pi$ cm$^2$\quad\quad \textbf{D.} 4500.0$\pi$ cm$^2$\par\vspace{2pt}
\textbf{\color{primary} Câu 15.} Xét đối tượng toán học trong Chương X: Một hình nón có bán kính đáy r = 3 cm và chiều cao h = 18 cm. Thể tích V của hình nón đó (V = 1/3 $\cdot$ $\pi$r$^2$h) bằng. Giá trị hoặc tính chất tương ứng là:\par
\hspace*{1.5em}\textbf{A.} 162$\pi$ cm$^3$\quad\quad \textbf{B.} 108$\pi$ cm$^3$\quad\quad \textbf{C.} 54$\pi$ cm$^3$\quad\quad \textbf{D.} 18$\pi$ cm$^3$\par\vspace{2pt}
\textbf{\color{primary} Câu 16.} Hoàn thành phép xác định cơ bản sau đây: Một hình trụ có bán kính đáy r = 17 cm và chiều cao h = 5 cm. Diện tích xung quanh của hình trụ đó bằng:\par
\hspace*{1.5em}\textbf{A.} 340$\pi$ cm$^2$\quad\quad \textbf{B.} 85$\pi$ cm$^2$\quad\quad \textbf{C.} 170$\pi$ cm$^2$\quad\quad \textbf{D.} 1445$\pi$ cm$^2$\par\vspace{2pt}
\textbf{\color{primary} Câu 17.} Trong các khẳng định hoặc giá trị dưới đây, đâu là đáp án đúng cho câu hỏi: Một mặt cầu có bán kính R = 18 cm. Diện tích S của mặt cầu đó (S = 4$\pi$R$^2$) bằng?\par
\hspace*{1.5em}\textbf{A.} 7776.0$\pi$ cm$^2$\quad\quad \textbf{B.} 1296$\pi$ cm$^2$\quad\quad \textbf{C.} 648$\pi$ cm$^2$\quad\quad \textbf{D.} 324$\pi$ cm$^2$\par\vspace{2pt}
\textbf{\color{primary} Câu 18.} Bài tập khởi động Chương X: Một hình nón có bán kính đáy r = 3 cm và chiều cao h = 21 cm. Thể tích V của hình nón đó (V = 1/3 $\cdot$ $\pi$r$^2$h) bằng:\par
\hspace*{1.5em}\textbf{A.} 63$\pi$ cm$^3$\quad\quad \textbf{B.} 189$\pi$ cm$^3$\quad\quad \textbf{C.} 126$\pi$ cm$^3$\quad\quad \textbf{D.} 21$\pi$ cm$^3$\par\vspace{2pt}
\textbf{\color{primary} Câu 19.} Quan sát dữ kiện đề bài: Một hình trụ có bán kính đáy r = 20 cm và chiều cao h = 5 cm. Diện tích xung quanh của hình trụ đó bằng. Hãy chọn phương án đúng:\par
\hspace*{1.5em}\textbf{A.} 100$\pi$ cm$^2$\quad\quad \textbf{B.} 400$\pi$ cm$^2$\quad\quad \textbf{C.} 2000$\pi$ cm$^2$\quad\quad \textbf{D.} 200$\pi$ cm$^2$\par\vspace{2pt}
\textbf{\color{primary} Câu 20.} Sử dụng công thức hoặc tính chất đầu bài của Chương X để giải quyết yêu cầu: Một mặt cầu có bán kính R = 21 cm. Diện tích S của mặt cầu đó (S = 4$\pi$R$^2$) bằng:\par
\hspace*{1.5em}\textbf{A.} 12348.0$\pi$ cm$^2$\quad\quad \textbf{B.} 1764$\pi$ cm$^2$\quad\quad \textbf{C.} 441$\pi$ cm$^2$\quad\quad \textbf{D.} 882$\pi$ cm$^2$\par\vspace{2pt}
\subsection*{\color{accent} Phần II. Mức độ Trung bình (Thông hiểu -- Câu 21 đến Câu 40)}
\textbf{\color{primary} Câu 21.} Một hình nón có bán kính đáy r = 3 cm và chiều cao h = 4 cm. Độ dài đường sinh l và diện tích xung quanh $S_{xq}$ = $\pi$rl của hình nón lần lượt là:\par
\hspace*{1.5em}\textbf{A.} l = 5 cm; $S_{xq}$ = 30$\pi$ cm$^2$\quad\quad \textbf{B.} l = 5 cm; $S_{xq}$ = 12$\pi$ cm$^2$\quad\quad \textbf{C.} l = 7 cm; $S_{xq}$ = 21$\pi$ cm$^2$\quad\quad \textbf{D.} l = 5 cm; $S_{xq}$ = 15$\pi$ cm$^2$\par\vspace{2pt}
\textbf{\color{primary} Câu 22.} Vận dụng quy tắc trong Chương X (Một số hình khối trong thực tiễn) để giải bài toán: Một hình cầu có bán kính R = 3 cm. Thể tích V của hình cầu đó (V = 4/3 $\cdot$ $\pi$R$^3$) bằng:\par
\hspace*{1.5em}\textbf{A.} 108$\pi$ cm$^3$\quad\quad \textbf{B.} 36$\pi$ cm$^3$\quad\quad \textbf{C.} 27$\pi$ cm$^3$\quad\quad \textbf{D.} 12$\pi$ cm$^3$\par\vspace{2pt}
\textbf{\color{primary} Câu 23.} Biến đổi và tính toán theo kiến thức Chương X: Một hình trụ có bán kính đáy r = 5 cm và chiều cao h = 7 cm. Thể tích của hình trụ đó bằng. Kết quả thu được là:\par
\hspace*{1.5em}\textbf{A.} 245$\pi$ cm$^3$\quad\quad \textbf{B.} 175$\pi$ cm$^3$\quad\quad \textbf{C.} 70$\pi$ cm$^3$\quad\quad \textbf{D.} 58.3$\pi$ cm$^3$\par\vspace{2pt}
\textbf{\color{primary} Câu 24.} Xét bài toán mức độ thông hiểu sau: Một hình nón có bán kính đáy r = 6 cm và chiều cao h = 8 cm. Độ dài đường sinh l và diện tích xung quanh $S_{xq}$ = $\pi$rl của hình nón lần lượt là. Đáp án đúng là:\par
\begin{itemize}\setlength{\itemsep}{1pt}\setlength{\parskip}{0pt}
  \item[] \textbf{A.} l = 10 cm; $S_{xq}$ = 120$\pi$ cm$^2$
  \item[] \textbf{B.} l = 14 cm; $S_{xq}$ = 84$\pi$ cm$^2$
  \item[] \textbf{C.} l = 10 cm; $S_{xq}$ = 48$\pi$ cm$^2$
  \item[] \textbf{D.} l = 10 cm; $S_{xq}$ = 60$\pi$ cm$^2$
\end{itemize}\vspace{2pt}
\textbf{\color{primary} Câu 25.} Để hoàn thành phiếu học tập Chương X, hãy xác định: Một hình cầu có bán kính R = 3 cm. Thể tích V của hình cầu đó (V = 4/3 $\cdot$ $\pi$R$^3$) bằng:\par
\hspace*{1.5em}\textbf{A.} 36$\pi$ cm$^3$\quad\quad \textbf{B.} 27$\pi$ cm$^3$\quad\quad \textbf{C.} 12$\pi$ cm$^3$\quad\quad \textbf{D.} 108$\pi$ cm$^3$\par\vspace{2pt}
\textbf{\color{primary} Câu 26.} Phân tích dữ kiện bài toán: Một hình trụ có bán kính đáy r = 8 cm và chiều cao h = 10 cm. Thể tích của hình trụ đó bằng. Giá trị cần tìm bằng:\par
\hspace*{1.5em}\textbf{A.} 640$\pi$ cm$^3$\quad\quad \textbf{B.} 800$\pi$ cm$^3$\quad\quad \textbf{C.} 160$\pi$ cm$^3$\quad\quad \textbf{D.} 213.3$\pi$ cm$^3$\par\vspace{2pt}
\textbf{\color{primary} Câu 27.} Áp dụng tính chất đã học ở Chương X vào bài tập sau: Một hình nón có bán kính đáy r = 3 cm và chiều cao h = 4 cm. Độ dài đường sinh l và diện tích xung quanh $S_{xq}$ = $\pi$rl của hình nón lần lượt là:\par
\hspace*{1.5em}\textbf{A.} l = 5 cm; $S_{xq}$ = 12$\pi$ cm$^2$\quad\quad \textbf{B.} l = 5 cm; $S_{xq}$ = 30$\pi$ cm$^2$\quad\quad \textbf{C.} l = 7 cm; $S_{xq}$ = 21$\pi$ cm$^2$\quad\quad \textbf{D.} l = 5 cm; $S_{xq}$ = 15$\pi$ cm$^2$\par\vspace{2pt}
\textbf{\color{primary} Câu 28.} Thực hiện các bước suy luận hợp lý để giải quyết yêu cầu: Một hình cầu có bán kính R = 3 cm. Thể tích V của hình cầu đó (V = 4/3 $\cdot$ $\pi$R$^3$) bằng:\par
\hspace*{1.5em}\textbf{A.} 108$\pi$ cm$^3$\quad\quad \textbf{B.} 36$\pi$ cm$^3$\quad\quad \textbf{C.} 27$\pi$ cm$^3$\quad\quad \textbf{D.} 12$\pi$ cm$^3$\par\vspace{2pt}
\textbf{\color{primary} Câu 29.} Cho bài toán tự luyện (Một số hình khối trong thực tiễn): Một hình trụ có bán kính đáy r = 11 cm và chiều cao h = 13 cm. Thể tích của hình trụ đó bằng. Phương án nào sau đây chính xác?\par
\hspace*{1.5em}\textbf{A.} 524.3$\pi$ cm$^3$\quad\quad \textbf{B.} 1573$\pi$ cm$^3$\quad\quad \textbf{C.} 1859$\pi$ cm$^3$\quad\quad \textbf{D.} 286$\pi$ cm$^3$\par\vspace{2pt}
\textbf{\color{primary} Câu 30.} Kết hợp các phép biến đổi trong Chương X, hãy tìm kết quả của bài toán: Một hình nón có bán kính đáy r = 6 cm và chiều cao h = 8 cm. Độ dài đường sinh l và diện tích xung quanh $S_{xq}$ = $\pi$rl của hình nón lần lượt là:\par
\begin{itemize}\setlength{\itemsep}{1pt}\setlength{\parskip}{0pt}
  \item[] \textbf{A.} l = 10 cm; $S_{xq}$ = 60$\pi$ cm$^2$
  \item[] \textbf{B.} l = 14 cm; $S_{xq}$ = 84$\pi$ cm$^2$
  \item[] \textbf{C.} l = 10 cm; $S_{xq}$ = 120$\pi$ cm$^2$
  \item[] \textbf{D.} l = 10 cm; $S_{xq}$ = 48$\pi$ cm$^2$
\end{itemize}\vspace{2pt}
\textbf{\color{primary} Câu 31.} Một bài kiểm tra định kỳ nội dung Chương X nêu câu hỏi: Một hình cầu có bán kính R = 3 cm. Thể tích V của hình cầu đó (V = 4/3 $\cdot$ $\pi$R$^3$) bằng. Hãy chọn đáp số đúng:\par
\hspace*{1.5em}\textbf{A.} 12$\pi$ cm$^3$\quad\quad \textbf{B.} 36$\pi$ cm$^3$\quad\quad \textbf{C.} 108$\pi$ cm$^3$\quad\quad \textbf{D.} 27$\pi$ cm$^3$\par\vspace{2pt}
\textbf{\color{primary} Câu 32.} Tìm lời giải đúng cho bài toán thông hiểu: Một hình trụ có bán kính đáy r = 14 cm và chiều cao h = 16 cm. Thể tích của hình trụ đó bằng:\par
\hspace*{1.5em}\textbf{A.} 3584$\pi$ cm$^3$\quad\quad \textbf{B.} 448$\pi$ cm$^3$\quad\quad \textbf{C.} 3136$\pi$ cm$^3$\quad\quad \textbf{D.} 1045.3$\pi$ cm$^3$\par\vspace{2pt}
\textbf{\color{primary} Câu 33.} Sử dụng mối liên hệ giữa các đại lượng trong Chương X: Một hình nón có bán kính đáy r = 3 cm và chiều cao h = 4 cm. Độ dài đường sinh l và diện tích xung quanh $S_{xq}$ = $\pi$rl của hình nón lần lượt là:\par
\hspace*{1.5em}\textbf{A.} l = 5 cm; $S_{xq}$ = 12$\pi$ cm$^2$\quad\quad \textbf{B.} l = 5 cm; $S_{xq}$ = 15$\pi$ cm$^2$\quad\quad \textbf{C.} l = 5 cm; $S_{xq}$ = 30$\pi$ cm$^2$\quad\quad \textbf{D.} l = 7 cm; $S_{xq}$ = 21$\pi$ cm$^2$\par\vspace{2pt}
\textbf{\color{primary} Câu 34.} Hãy tính toán và chọn phương án đúng đối với bài tập: Một hình cầu có bán kính R = 3 cm. Thể tích V của hình cầu đó (V = 4/3 $\cdot$ $\pi$R$^3$) bằng:\par
\hspace*{1.5em}\textbf{A.} 108$\pi$ cm$^3$\quad\quad \textbf{B.} 36$\pi$ cm$^3$\quad\quad \textbf{C.} 27$\pi$ cm$^3$\quad\quad \textbf{D.} 12$\pi$ cm$^3$\par\vspace{2pt}
\textbf{\color{primary} Câu 35.} Dựa trên các tính chất trọng tâm của Chương X, xét bài toán: Một hình trụ có bán kính đáy r = 17 cm và chiều cao h = 19 cm. Thể tích của hình trụ đó bằng:\par
\hspace*{1.5em}\textbf{A.} 6137$\pi$ cm$^3$\quad\quad \textbf{B.} 1830.3$\pi$ cm$^3$\quad\quad \textbf{C.} 5491$\pi$ cm$^3$\quad\quad \textbf{D.} 646$\pi$ cm$^3$\par\vspace{2pt}
\textbf{\color{primary} Câu 36.} Thực hành giải toán (Một số hình khối trong thực tiễn): Một hình nón có bán kính đáy r = 6 cm và chiều cao h = 8 cm. Độ dài đường sinh l và diện tích xung quanh $S_{xq}$ = $\pi$rl của hình nón lần lượt là. Đáp án nào dưới đây thỏa mãn?\par
\begin{itemize}\setlength{\itemsep}{1pt}\setlength{\parskip}{0pt}
  \item[] \textbf{A.} l = 10 cm; $S_{xq}$ = 120$\pi$ cm$^2$
  \item[] \textbf{B.} l = 14 cm; $S_{xq}$ = 84$\pi$ cm$^2$
  \item[] \textbf{C.} l = 10 cm; $S_{xq}$ = 60$\pi$ cm$^2$
  \item[] \textbf{D.} l = 10 cm; $S_{xq}$ = 48$\pi$ cm$^2$
\end{itemize}\vspace{2pt}
\textbf{\color{primary} Câu 37.} Kiểm tra kỹ năng biến đổi và suy luận qua bài toán: Một hình cầu có bán kính R = 3 cm. Thể tích V của hình cầu đó (V = 4/3 $\cdot$ $\pi$R$^3$) bằng:\par
\hspace*{1.5em}\textbf{A.} 36$\pi$ cm$^3$\quad\quad \textbf{B.} 27$\pi$ cm$^3$\quad\quad \textbf{C.} 108$\pi$ cm$^3$\quad\quad \textbf{D.} 12$\pi$ cm$^3$\par\vspace{2pt}
\textbf{\color{primary} Câu 38.} Cho giả thiết thuộc phạm vi Chương X: Một hình trụ có bán kính đáy r = 20 cm và chiều cao h = 22 cm. Thể tích của hình trụ đó bằng. Kết luận đúng là:\par
\hspace*{1.5em}\textbf{A.} 2933.3$\pi$ cm$^3$\quad\quad \textbf{B.} 880$\pi$ cm$^3$\quad\quad \textbf{C.} 8800$\pi$ cm$^3$\quad\quad \textbf{D.} 9680$\pi$ cm$^3$\par\vspace{2pt}
\textbf{\color{primary} Câu 39.} Rút gọn hoặc tính giá trị theo yêu cầu sau của Chương X: Một hình nón có bán kính đáy r = 3 cm và chiều cao h = 4 cm. Độ dài đường sinh l và diện tích xung quanh $S_{xq}$ = $\pi$rl của hình nón lần lượt là:\par
\hspace*{1.5em}\textbf{A.} l = 5 cm; $S_{xq}$ = 30$\pi$ cm$^2$\quad\quad \textbf{B.} l = 7 cm; $S_{xq}$ = 21$\pi$ cm$^2$\quad\quad \textbf{C.} l = 5 cm; $S_{xq}$ = 12$\pi$ cm$^2$\quad\quad \textbf{D.} l = 5 cm; $S_{xq}$ = 15$\pi$ cm$^2$\par\vspace{2pt}
\textbf{\color{primary} Câu 40.} Xác định phương án đúng sau khi thực hiện phép tính / suy luận: Một hình cầu có bán kính R = 3 cm. Thể tích V của hình cầu đó (V = 4/3 $\cdot$ $\pi$R$^3$) bằng:\par
\hspace*{1.5em}\textbf{A.} 27$\pi$ cm$^3$\quad\quad \textbf{B.} 108$\pi$ cm$^3$\quad\quad \textbf{C.} 12$\pi$ cm$^3$\quad\quad \textbf{D.} 36$\pi$ cm$^3$\par\vspace{2pt}
\subsection*{\color{accent} Phần III. Mức độ Khó (Vận dụng \& Vận dụng cao -- Câu 41 đến Câu 50)}
\textbf{\color{primary} Câu 41.} Quay một hình chữ nhật ABCD có AB = 3 cm, BC = 4 cm một vòng quanh cạnh cố định AB, ta thu được một hình trụ có thể tích bằng:\par
\hspace*{1.5em}\textbf{A.} 12$\pi$ cm$^3$\quad\quad \textbf{B.} 36$\pi$ cm$^3$\quad\quad \textbf{C.} 48$\pi$ cm$^3$\quad\quad \textbf{D.} 24$\pi$ cm$^3$\par\vspace{2pt}
\textbf{\color{primary} Câu 42.} Bài toán vận dụng (Chương X – Một số hình khối trong thực tiễn): Một hình trụ và một hình nón có cùng bán kính đáy r và cùng chiều cao h. Tỉ số thể tích của hình trụ so với thể tích của hình nón đó luôn bằng. Kết quả đúng là:\par
\hspace*{1.5em}\textbf{A.} 9\quad\quad \textbf{B.} 1/3\quad\quad \textbf{C.} 2\quad\quad \textbf{D.} 3\par\vspace{2pt}
\textbf{\color{primary} Câu 43.} Vận dụng kiến thức nâng cao trong phạm vi Chương X để giải bài toán: Quay một hình chữ nhật ABCD có AB = 5 cm, BC = 6 cm một vòng quanh cạnh cố định AB, ta thu được một hình trụ có thể tích bằng:\par
\hspace*{1.5em}\textbf{A.} 180$\pi$ cm$^3$\quad\quad \textbf{B.} 30$\pi$ cm$^3$\quad\quad \textbf{C.} 60$\pi$ cm$^3$\quad\quad \textbf{D.} 150$\pi$ cm$^3$\par\vspace{2pt}
\textbf{\color{primary} Câu 44.} Phân tích và giải bài toán tổng hợp sau của Chương X: Một hình trụ và một hình nón có cùng bán kính đáy r và cùng chiều cao h. Tỉ số thể tích của hình trụ so với thể tích của hình nón đó luôn bằng:\par
\hspace*{1.5em}\textbf{A.} 2\quad\quad \textbf{B.} 1/3\quad\quad \textbf{C.} 9\quad\quad \textbf{D.} 3\par\vspace{2pt}
\textbf{\color{primary} Câu 45.} Thử thách tư duy trong phạm vi kiến thức Chương X: Quay một hình chữ nhật ABCD có AB = 7 cm, BC = 8 cm một vòng quanh cạnh cố định AB, ta thu được một hình trụ có thể tích bằng. Đáp số chính xác là:\par
\hspace*{1.5em}\textbf{A.} 56$\pi$ cm$^3$\quad\quad \textbf{B.} 392$\pi$ cm$^3$\quad\quad \textbf{C.} 448$\pi$ cm$^3$\quad\quad \textbf{D.} 112$\pi$ cm$^3$\par\vspace{2pt}
\textbf{\color{primary} Câu 46.} Sử dụng phương pháp giải toán của Chương X để tìm đáp án cho bài toán: Một hình trụ và một hình nón có cùng bán kính đáy r và cùng chiều cao h. Tỉ số thể tích của hình trụ so với thể tích của hình nón đó luôn bằng:\par
\hspace*{1.5em}\textbf{A.} 1/3\quad\quad \textbf{B.} 3\quad\quad \textbf{C.} 2\quad\quad \textbf{D.} 9\par\vspace{2pt}
\textbf{\color{primary} Câu 47.} Xét bài toán vận dụng thực tiễn / nâng cao thuộc Chương X: Quay một hình chữ nhật ABCD có AB = 9 cm, BC = 10 cm một vòng quanh cạnh cố định AB, ta thu được một hình trụ có thể tích bằng:\par
\hspace*{1.5em}\textbf{A.} 810$\pi$ cm$^3$\quad\quad \textbf{B.} 180$\pi$ cm$^3$\quad\quad \textbf{C.} 90$\pi$ cm$^3$\quad\quad \textbf{D.} 900$\pi$ cm$^3$\par\vspace{2pt}
\textbf{\color{primary} Câu 48.} Tìm kết quả cuối cùng của bài toán phân loại học sinh khá giỏi (Chương X): Một hình trụ và một hình nón có cùng bán kính đáy r và cùng chiều cao h. Tỉ số thể tích của hình trụ so với thể tích của hình nón đó luôn bằng:\par
\hspace*{1.5em}\textbf{A.} 9\quad\quad \textbf{B.} 2\quad\quad \textbf{C.} 1/3\quad\quad \textbf{D.} 3\par\vspace{2pt}
\textbf{\color{primary} Câu 49.} Kết hợp các tính chất đã học trong Chương X, hãy giải quyết bài toán: Quay một hình chữ nhật ABCD có AB = 11 cm, BC = 12 cm một vòng quanh cạnh cố định AB, ta thu được một hình trụ có thể tích bằng:\par
\hspace*{1.5em}\textbf{A.} 1584$\pi$ cm$^3$\quad\quad \textbf{B.} 264$\pi$ cm$^3$\quad\quad \textbf{C.} 1452$\pi$ cm$^3$\quad\quad \textbf{D.} 132$\pi$ cm$^3$\par\vspace{2pt}
\textbf{\color{primary} Câu 50.} Câu hỏi vận dụng cao cuối Chương X: Một hình trụ và một hình nón có cùng bán kính đáy r và cùng chiều cao h. Tỉ số thể tích của hình trụ so với thể tích của hình nón đó luôn bằng. Phương án đúng là:\par
\hspace*{1.5em}\textbf{A.} 1/3\quad\quad \textbf{B.} 9\quad\quad \textbf{C.} 2\quad\quad \textbf{D.} 3\par\vspace{2pt}
\subsection*{\color{primary} Bảng đáp án Chương X: Một số hình khối trong thực tiễn}
\begin{center}
\renewcommand{\arraystretch}{1.25}
\begin{tabularx}{\textwidth}{*{10}{>{\centering\arraybackslash}X}}
\toprule
\textbf{\color{primary} Cột 1} & \textbf{\color{primary} Cột 2} & \textbf{\color{primary} Cột 3} & \textbf{\color{primary} Cột 4} & \textbf{\color{primary} Cột 5} & \textbf{\color{primary} Cột 6} & \textbf{\color{primary} Cột 7} & \textbf{\color{primary} Cột 8} & \textbf{\color{primary} Cột 9} & \textbf{\color{primary} Cột 10} \\
\midrule
\textbf{1.A} & \textbf{2.D} & \textbf{3.B} & \textbf{4.C} & \textbf{5.A} & \textbf{6.C} & \textbf{7.A} & \textbf{8.B} & \textbf{9.C} & \textbf{10.A} \\
\rowcolor{rowtint} \textbf{11.B} & \textbf{12.B} & \textbf{13.C} & \textbf{14.C} & \textbf{15.C} & \textbf{16.C} & \textbf{17.B} & \textbf{18.A} & \textbf{19.D} & \textbf{20.B} \\
\textbf{21.D} & \textbf{22.B} & \textbf{23.B} & \textbf{24.D} & \textbf{25.A} & \textbf{26.A} & \textbf{27.D} & \textbf{28.B} & \textbf{29.B} & \textbf{30.A} \\
\rowcolor{rowtint} \textbf{31.B} & \textbf{32.C} & \textbf{33.B} & \textbf{34.B} & \textbf{35.C} & \textbf{36.C} & \textbf{37.A} & \textbf{38.C} & \textbf{39.D} & \textbf{40.D} \\
\textbf{41.C} & \textbf{42.D} & \textbf{43.A} & \textbf{44.D} & \textbf{45.C} & \textbf{46.B} & \textbf{47.D} & \textbf{48.D} & \textbf{49.A} & \textbf{50.D} \\
\bottomrule
\end{tabularx}
\end{center}\vspace{10pt}
\end{document}
